\documentclass[11pt,oneside]{article}
\usepackage[T1]{fontenc}
\usepackage[utf8]{inputenc}
\usepackage{lmodern,microtype}
\usepackage[margin=1in,headheight=15pt]{geometry}
\usepackage{amsmath,amssymb,amsthm,mathtools}
\usepackage{booktabs,longtable,array,enumitem,xcolor}
\usepackage{listings,fancyhdr}
\usepackage{xurl}
\usepackage[colorlinks=true,linkcolor=blue!45!black,citecolor=blue!45!black,urlcolor=blue!45!black]{hyperref}
\usepackage{bookmark}
\newcommand{\ipra}{\ensuremath{\mathrm{iPRA}}}
\newcommand{\N}{\mathbb N}
\newcommand{\bits}{\{0,1\}}
\newcommand{\poly}{\operatorname{poly}}
\newcommand{\len}[1]{\lvert #1\rvert}
\newcommand{\code}[1]{\ulcorner #1\urcorner}
\newcommand{\B}{\mathsf B}
\newcommand{\Run}{\operatorname{Run}}
\newcommand{\Any}{\operatorname{Any}}
\newcommand{\All}{\operatorname{All}}
\newcommand{\sg}{\operatorname{sg}}
\newcommand{\size}{\operatorname{size}}
\newcommand{\Space}{\operatorname{Space}}
\newcommand{\Time}{\operatorname{Time}}
\newcommand{\Pclass}{\mathsf P}
\newcommand{\NP}{\mathsf{NP}}
\newcommand{\coNP}{\mathsf{coNP}}
\newcommand{\PSPACE}{\mathsf{PSPACE}}
\newcommand{\EXPTIME}{\mathsf{EXPTIME}}
\newcommand{\EXPSPACE}{\mathsf{EXPSPACE}}
\newcommand{\DTIME}{\mathsf{DTIME}}
\newcommand{\DSPACE}{\mathsf{DSPACE}}
\newcommand{\NSPACE}{\mathsf{NSPACE}}
\newcommand{\NTIME}{\mathsf{NTIME}}
\newcommand{\SAT}{\mathsf{SAT}}
\newcommand{\QBF}{\mathsf{TQBF}}
\newcommand{\Reach}{\operatorname{Reach}}
\newcommand{\true}{\mathsf{true}}
\newcommand{\false}{\mathsf{false}}
\newtheorem{theorem}{Theorem}[section]
\newtheorem{lemma}[theorem]{Lemma}
\newtheorem{corollary}[theorem]{Corollary}
\newtheorem{proposition}[theorem]{Proposition}
\theoremstyle{definition}
\newtheorem{definition}[theorem]{Definition}
\newtheorem{remark}[theorem]{Remark}
\newtheorem{example}[theorem]{Example}
\setlist{itemsep=3pt,topsep=5pt}
\lstset{basicstyle=\ttfamily\footnotesize,breaklines=true,columns=fullflexible,keepspaces=true,frame=single,rulecolor=\color{black!20},xleftmargin=3pt,xrightmargin=3pt}
\setlength{\emergencystretch}{2em}
\pagestyle{fancy}
\fancyhf{}
\fancyhead[L]{\small Constructive complexity in iPRA}
\fancyhead[R]{\small Vladimir Reshetnikov research packet}
\fancyfoot[C]{\thepage}
\hypersetup{pdftitle={A constructive complexity-theory atlas in intuitionistic primitive recursive arithmetic},pdfauthor={Research prepared for Vladimir Reshetnikov},pdfsubject={Finite proofs, computational witnesses, and Lean formalization plan}}
\begin{document}
\begin{titlepage}
\thispagestyle{empty}
\vspace*{1cm}
{\Large\color{blue!45!black} CONSTRUCTIVE MATHEMATICS / COMPLEXITY THEORY\par}
\vspace{1cm}
{\huge\bfseries A complexity-theory atlas in intuitionistic primitive recursive arithmetic\par}
\vspace{0.7cm}
{\Large Explicit algorithms, finite witnesses, resource bounds, and a route to Lean\par}
\vspace{1.2cm}
{\large Prepared for Vladimir Reshetnikov\par}
{\large 22 November 2026\par}
\vfill
\begin{minipage}{0.94\textwidth}
\textbf{Mathematical deliverable.} A broad collection of fully developed finitary proofs, organized around the computational-complexity article and its immediate subject links. Every asserted arithmetic theorem uses numerical codes, explicit primitive-recursive bounds, and quantifier-free induction through Boolean invariants.

\medskip
\textbf{Verification boundary.} The new Lean files are an uncompiled formalization outline. This packet does not claim completed kernel-checked object-language derivations of its main complexity theorems. Executable finite checks and PDF validation are reported separately. Existing ProveIt endpoints are cited at a fixed revision and are not re-certified here.
\end{minipage}
\vspace{0.7cm}
\end{titlepage}
\tableofcontents
\newpage
\section{Results, scope, and how to read the claims}
\label{sec:scope}

This report develops \textbf{86 labelled theorem and lemma groups} around the mathematical content of the requested computational-complexity article and its immediate subject links \cite{WikiRoot}. Several groups contain more than one conventional named result. The count is a navigational inventory, not a claim of 86 independent discoveries. The classical theorems are established results; our work is their explicit finite constructive reconstruction, the separation of their numerical readings from informal class notation, and a detailed route to object-language proofs in intuitionistic primitive recursive arithmetic.

\begin{center}
\begin{tabular}{p{0.12\textwidth}p{0.76\textwidth}}
\toprule
Labels & Principal content \\
\midrule
F01--F06 & The theory, finite witnesses, bounded induction, coded computation, and simulation. \\
C01--C06 & Class presentations and closures; unary/binary time and space completeness. \\
N01--N22 & Cook--Levin; graph, integer, and counting reductions; 2SAT and Horn SAT; search, approximation, and parameterized vertex cover. \\
S01--S13 & Configuration graphs; Savitch; inductive counting and complement closure; QBF and finite games. \\
H01--H05 & Explicit time/space diagonal witnesses and coarse hierarchy separations. \\
A01--A07 & AP = PSPACE; finite PH levels and conditional collapses; uniform NC and AC constructions. \\
R01--R23 & Randomness, PP, advice, second-level PH containment, zero error, counting, circuit bounds, queries, communication, and amortization. \\
Q01--Q04 & Exact integer path sums and classical containments for the specified Hadamard--Toffoli gate model. \\
\bottomrule
\end{tabular}
\end{center}

The most substantial conclusions include
\begin{gather*}
 \mathsf{NPSPACE}=\mathsf{PSPACE},\quad
 \mathsf{NEXPSPACE}=\mathsf{EXPSPACE},\quad
 \mathsf{NL}=\mathsf{coNL},\quad
 \mathsf{AP}=\mathsf{PSPACE},\\
 \mathsf P\subsetneq\mathsf{EXPTIME},\qquad
 \mathsf L\subsetneq\mathsf{PSPACE}\subsetneq\mathsf{EXPSPACE},\\
 \mathsf{BPP}\subseteq\mathsf P/\poly,\qquad
 \mathsf{BPP}\subseteq\Sigma_2^{\mathsf P}\cap\Pi_2^{\mathsf P},\qquad
 \mathsf{NP}\cup\mathsf{coNP}\cup\mathsf{BPP}\subseteq\mathsf{PP}\subseteq\mathsf{PSPACE}.
\end{gather*}
All class notation is shorthand for the numerical presentation and compiler statements specified in the report. For general space functions, a PR budget and the stated space constructor are hypotheses. Nonuniform advice is identified as such. Quantum conclusions retain the explicit gate-model subscript. None of the conclusions resolves P versus NP or the other unproved equalities in the article.

\subsection{Three different verification statements}
The mathematical proofs are human-readable constructive arguments. They explicitly give the algorithms, witnesses, bounds, and induction predicates that place the stated numerical theorems within \ipra. This is the main deliverable.

The Lean deliverable is \emph{an outline with concrete semantic definitions and precise remaining proof obligations}. No Lean executable was available in this workspace, and the new files have not been compiled. They contain no admitted substantive proofs presented as completed theorems. Their eventual target is an actual derivation in ProveIt's intuitionistic arithmetic, not merely a theorem about natural numbers in Lean's ambient logic.

The Python checks are independent finite implementation checks. They test nontrivial identities, reductions, witness multiplicities, table constructions, and resource-related edge cases on finite examples. They do not establish the universal mathematical claims, prove asymptotic resource bounds, or certify the object-language derivations. The archive records these distinctions per claim in \texttt{CLAIMS.json}.

\subsection{Relationship with the existing ProveIt work}
We inspected ProveIt at revision
\begin{center}\small\texttt{3eb195cab4cc3d552a9eca78880c6d428f3e44f9}.\end{center}
The November 20 packet \texttt{ipra\_complexity\_foundations\_2026-11-20} already developed bounded formulas, verifier circuits, Savitch, NL complementation, compact QBF, an exponential diagonal, and sketches of several extensions \cite{ProveItBaseline}. We retain those compatible foundations and give a self-contained exposition, rather than counting them as new mathematical discoveries. The present expansion includes extensive concrete completeness reductions, parsimonious SUBSET SUM, Horn P-completeness, detailed probability/PP/advice arguments, second-level hierarchy containment, explicit space separations, shallow-circuit simulations, query and communication lower bounds, and exact finite-gate quantum simulations.

Actual reusable repository endpoints include \texttt{PR.lean}, the bounded-formula compiler and bounded induction, and arithmetic estimates in \texttt{DiagonalClock.lean}. A similarly named \texttt{ReachableFrontier.lean} is a metatheoretic theorem about proof graphs; it is not a ready-made object-arithmetic configuration-graph theorem. We preserve that distinction in the source audit. Upstream source was inspected, not rebuilt, and its prior build claims are not substituted for verification of this new packet.

\subsection{Why a finite atlas is the appropriate reading of the request}
The article links very broad subjects such as algorithms, graphs, computability, probability, and quantum computation. Their full mathematical contents have no practical finite maximum. We therefore prioritize a broad connected core of important statements, develop the substantial proofs, and give a selected topic-to-result ledger. That ledger is not an exhaustive crawl of every navigation, biography, language, or subject link. Advanced theorems not proved here are listed explicitly at the end, together with the missing mathematical work. A restricted finite counterpart is never labelled as a proof of an unrestricted theorem.

\section{The arithmetic foundation and its proof discipline}
\label{sec:foundations}

\subsection{Exactly which theory is used}
There is one sort, the natural numbers. The language contains $0$, successor $S$, and a function symbol for each fixed primitive-recursive (PR) definition. Zero, projections, composition, and primitive recursion generate these definitions. If $h$ is defined from $f,g$, its axioms are
\[
 h(0,\vec x)=f(\vec x),\qquad
 h(Sn,\vec x)=g(n,h(n,\vec x),\vec x).
\]
Composition and projection symbols have their defining equations. The remaining arithmetic axioms are $Sx\ne0$, $Sx=Sy\to x=y$, and
\[
 A(0,\vec x)\to
 [\forall n\,(A(n,\vec x)\to A(Sn,\vec x))]\to
 \forall n\,A(n,\vec x)
\]
for \emph{quantifier-free} $A$. The logic is intuitionistic first-order logic with equality: the introduction and elimination rules for $\bot,\land,\lor,\to,\forall,\exists$, equality reflexivity and substitution, and the usual eigenvariable restrictions. Negation abbreviates implication to $\bot$. There is no unrestricted excluded middle, choice, Markov principle, or self-reflection principle. This agrees with ProveIt's actual \texttt{PR}, \texttt{PTm}, \texttt{PRAAx}, and \texttt{Prf} interface \cite{ProveItPR}.

Addition, truncated subtraction, multiplication, exponentiation, comparison, quotient and remainder, finite sums, and finite products have their elementary PR definitions. Their usual laws follow by quantifier-free induction. For quotient by $b>0$, one may maintain a remainder $r<b$ and a quotient $q$ as the dividend counter increases: increment $r$ unless it becomes $b$, in which case reset it to zero and increment $q$. The invariant $a=bq+r$ proves division with remainder using only numerical equalities and inequalities. Set quotient and remainder to zero when $b=0$ so every symbol is total. Signed integers are coded by sign and magnitude with a distinguished zero; their arithmetic reduces to these natural-number operations.

It is essential to distinguish two uses of ``primitive recursive.'' A PR symbol is legitimate in the arithmetic language, but it is not a unit-cost operation in the machine model. Conversely, there is no total PR function $U(e,x)$ evaluating every unary PR function at its numerical code: if there were, the PR function $x\mapsto U(x,x)+1$ would contradict its own alleged index. We use a fixed \emph{clocked} machine interpreter instead.

\subsection{Words and finite data}
A word $w=b_0\cdots b_{n-1}$ is represented by
\[
 \code w=2^n+\sum_{i<n}b_i2^i.
\]
Thus the empty word has code $1$, and $0$ is reserved as invalid. The length is the unique $n$ satisfying $2^n\le\code w<2^{n+1}$. Find it by bounded search below the word code. Bit access uses division and remainder. Appending a word, taking a substring, reversing, and reading a self-delimiting header are bounded loops. The numerical code can be exponentially larger than its bit length, and all resource bounds below refer to \emph{bit length}.

For an array of length $r$ with entries below $b\ge2$, use $\sum_{i<r}a_i b^i<b^r$, with length and base in a header. Its $i$th entry is $(A\mathbin{\mathrm{div}}b^i)\bmod b$. All proof-theoretic arrays are of this form or nested such records. Machine inputs and outputs use binary fields with self-delimiting lengths; a pairing with unnecessarily long unary fields is never silently treated as a polynomial-size serialization. Malformed words, graphs, circuits, clauses, and machines have a specified default interpretation, normally rejection. Every theorem about a meaningful instance includes the corresponding validity test, or uses a total constructor that guarantees validity.

\begin{lemma}[F01: finite decisions and least witnesses]\label{thm:F01}
For every \emph{fixed} PR Boolean predicate $b(i,\vec x)$, there are PR functions
\[
 \Any_b(r,\vec x),\quad \All_b(r,\vec x),\quad \mu_b(r,\vec x)\le r
\]
such that \ipra\ proves
\begin{align*}
 \Any_b(r,\vec x)=1&\ \longleftrightarrow\ \exists i<r\ b(i,\vec x)=1,\\
 \All_b(r,\vec x)=1&\ \longleftrightarrow\ \forall i<r\ b(i,\vec x)=1.
\end{align*}
If a successful index exists, $\mu_b$ is its least value; otherwise $\mu_b=r$.
\end{lemma}
\begin{proof}
Define the first two functions by folding OR and AND, starting from $0$ and $1$. For search maintain a found-bit $f$ and an index $j$. At stage $i$, if $f=0$ and $b(i,\vec x)=1$, set $(f,j)=(1,i)$; otherwise leave them unchanged. At the end return $j$ if $f=1$, and $r$ if $f=0$.

The loop invariants are: $f$ is the OR of the processed bits; $f=1$ implies $j<i$, $b(j)=1$, and the product of the failure bits below $j$ is $1$; $f=0$ implies the product of all failure bits below $i$ is $1$. These are quantifier-free equations or Boolean combinations of equations between PR terms. Each of the two updates preserves them by the fold equations. Induction is therefore permitted.

To recover quantified meanings, a given successful index forces the OR to be one: split the fold before and after that index. A one OR gives the saved witness. Dually, a zero AND gives a least failed factor, and a one AND forbids any failed factor. This proves both directions of the equivalences without induction on an unbounded existential formula. The zero-length cases give an empty OR of zero, an empty AND of one, and sentinel zero.
\end{proof}

\begin{theorem}[F02: bounded formulas and bounded induction]\label{thm:F02}
For each fixed formula $A$ whose quantifiers have PR bounds, there is a fixed PR bit $\chi_A$ and an \ipra\ proof $A\leftrightarrow\chi_A=1$. Induction for $A$ is consequently derivable from quantifier-free induction.
\end{theorem}
\begin{proof}
Equality and order have PR comparison bits. Conjunction and disjunction become Boolean product and OR. Since compiled subformulas are decidable, implication becomes $\neg\chi_B\lor\chi_C$; its equivalence with $B\to C$ is verified in the four Boolean cases. Bounded quantifiers become the folds of Lemma~\ref{thm:F01}. This constructs a finite PR definition and a finite proof by recursion on the \emph{external fixed syntax} of $A$. It is not a numerical truth predicate for arbitrary PR syntax.

Translate the base and successor assumptions for $A(n)$ into those for the quantifier-free equation $\chi_A(n)=1$, apply its induction axiom, and translate back. Unbounded formulas may occur in the surrounding context; they are never put inside the induction predicate. This is exactly the mechanism exposed by ProveIt's \texttt{BoundedSyntaxInduction.induction} \cite{ProveItBounded}.
\end{proof}

\begin{lemma}[F03: bounded-loop correctness]\label{thm:F03}
Suppose a finite program has a PR fuel bound $T(\vec x)$, PR state updates, and a PR bound $B(\vec x)$ on all state codes throughout that run. If an invariant has only quantifiers bounded by explicit PR terms, its base case and preservation prove correctness in \ipra. This includes bounded least search, array dynamic programming, and a depth-first traversal of a finitely bounded tree.
\end{lemma}
\begin{proof}
Replace a while-loop by $T$ iterations that leave the state unchanged after its stopping bit becomes one. Simultaneous recursion on the tuple of variables is primitive recursion on a code. For nested loops, use nested recursions, or a flattened counter below the product of their bounds. Apply Theorem~\ref{thm:F02} to the assertion that either the counter is beyond $T$ or the coded state satisfies the invariant. Branching is a finite Boolean case distinction. A depth-first tree with branching at most $b$ and height at most $h$ has at most $1+b+\cdots+b^h$ nodes, so its stack traversal has a PR clock and a PR stack-code bound. This justifies the proof rule, not an assertion that every boundedly computable program has a correctness proof. Its actual invariant must still be established.
\end{proof}

\begin{lemma}[F04: finite counting and pigeonhole]\label{thm:F04}
For any explicitly coded map $f:\{0,\ldots,r-1\}\to\{0,\ldots,s-1\}$, injectivity implies $r\le s$. If $r>s$, bounded search constructs $i<j<r$ with $f(i)=f(j)$. For any Boolean predicates $P_i$ on a finite set $X$, the size of their union is at most the sum of their sizes.
\end{lemma}
\begin{proof}
Build a bitmap of the values seen in the first $k$ entries. Initially its sum is zero. Inserting a previously unseen entry increments the sum by one; inserting a repeated entry leaves it unchanged and produces a pair of equal values. If all $r$ entries are distinct, its sum is $r$ and is at most the number $s$ of positions. If $r>s$, search all $r^2$ pairs; absence of a found pair would give the impossible bitmap sum. The search succeeds constructively because its success bit is decidable. For the union bound, at each $x$ the OR of $P_i(x)$ is at most the sum of their bits; sum over $x$. All invariants are finite sums or bounded bitmap assertions.
\end{proof}

\subsection{The machine model and clocked semantics}
We use deterministic or binary-branching Turing machines with a fixed finite alphabet, a read-only input tape with endmarkers when work space is measured, finitely many work tapes, and optional one-way write-only output. The output-head position is unavailable to the control and the output tape has no readable boundary signal. Time counts transitions. Space counts work-tape cells visited; input and write-only output do not count. Small-space transducers are additionally required to halt; their output length and time follow from the usual finite-configuration argument because their control never reads the output. Fixed machines can be normalized to one work tape with a fixed alphabet with constant-factor space overhead. Ordinary one-tape machines with input on the work tape are used only for time arguments, where the input contributes $n$ cells.

Tables, states, head positions, tapes, and accepting/rejecting flags have finite numerical codes. Accepting and rejecting states are absorbing. An invalid table denotes an immediately rejecting machine. A branch-controlled interpreter uses an extra finite choice word; one bit is consumed at each branching step. Finite random-tape interpreters use exactly the same mechanism but count choice words instead of existentially quantifying over them.

\begin{theorem}[F05: bounded simulation and trace uniqueness]\label{thm:F05}
There are fixed PR functions $\Run(e,x,t)$ and $\B(e,x,t)\in\bits$ for the configuration and acceptance bit after $t$ transitions. A correctly coded trace of length $t+1$ beginning at the initial configuration agrees with $\Run$ at every position. Accepting or rejecting output remains unchanged under additional fuel. All statements have \ipra\ proofs.
\end{theorem}
\begin{proof}
One transition is table lookup followed by replacement of finitely many fields, so it is PR. Define $\Run(e,x,0)$ as initialization and $\Run(e,x,t+1)$ as one step applied to $\Run(e,x,t)$. For a fixed horizon $T$, a one-tape head cannot leave $[-T,n+T]$, and only $n+2T+1$ tape cells matter. If there are $q$ states and alphabet size $g$, the number of configurations has the elementary bound
\[
 C=(q+2)(n+2T+3)g^{n+2T+3}
\]
after harmless enlargements for markers. This displayed formula is for an ordinary input-on-work-tape machine. With an independently moving read-only input head, include the additional factor $n+2$. A trace can be stored below $\max(2,C)^{T+1}$. For multiple tapes take the product of the corresponding bounds.

The invariant $\forall i\le\min(j,T),\ \operatorname{trace}(z,i)=\Run(e,x,i)$ is bounded by those codes. Its base is initialization and its step is the deterministic transition equation. Lemma~\ref{thm:F03} proves it. Absorption follows by induction on extra steps. With a supplied branch word the transition is again deterministic, so the same argument applies. No assertion that an unclocked program halts was needed.
\end{proof}

\begin{theorem}[F06: concrete polynomial simulation]\label{thm:F06}
For a fixed number of source tapes, time $t$ can be simulated on one tape in $O((n+t+1)^2)$ time. There is also a fixed interpreter that, with a finite source table of length $m$ and a unary clock $t$, computes $\B(e,x,t)$ in at most $C(m+n+t+1)^4$ transitions for a fixed constant $C$. Both simulations and their bounds admit bounded-invariant \ipra\ proofs.
\end{theorem}
\begin{proof}
Put source tapes on separate tracks of one physical tape, with a head marker on each. One full sweep collects the scanned symbols into finite control; a return sweep updates symbols and markers. A bounded number of sweeps simulates one move. At stage $i$, occupied length is $O(n+i+1)$, so summing $O(n+i+1)$ over $i<t$ gives the first bound. Since the number of tapes is fixed for that compiler, its finite control and track alphabet are fixed too.

For a variable program table, use canonical descriptions with one work tape over a fixed alphabet (and, if present, a separate read-only input tape). Other fixed-machine presentations are compiled to this form first. Scan its $m$ bits to locate the relevant rule, using an explicit finite word for the source state. Repeated comparisons, done by scans, cost at most $O((m+1)^2)$ virtual steps per source step. Store a unary fuel counter and the simulated tape. They occupy $O(m+n+t+1)$ cells. One-tape sweeps therefore cost at most
\[
 C\,t(m+1)^2(m+n+t+1)\le C(m+n+t+1)^4.
\]
Increase $C$ for parsing, initialization, and invalid descriptions. The preservation invariant records every represented cell, all head markers, the current table rule, the remaining fuel, and the occupied interval. Its scans and cell indices are bounded. Summing the scan counters gives the displayed costs by the same bounded induction. This supplies the coarse simulation used here; the sharper $t\log t$ simulation is not assumed \cite{HennieStearns1966}.
\end{proof}

The theory can therefore justify computations with immense finite bounds while still making meaningful distinctions between polynomial, exponential, and space-bounded machines. Arithmetic definability and computational efficiency remain separate claims throughout the report.

\section{Presentations of complexity classes and basic constructions}
\label{sec:classes}

\subsection{Numerical presentations, not sets as arithmetic objects}
A \emph{language bit} is a specified PR Boolean function $A(x)$ on word codes $x\ge1$. More generally a displayed first-order predicate is identified with such a bit by a proved equivalence. This restricts the representation, not the usual classes discussed below: every deterministic polynomial/exponential-time machine, polynomial-space machine, or polynomial verifier has a PR acceptance bit. It does not represent all computable or all abstract languages by one PR evaluator.

Put $p_{c,k}(n)=c(n+1)^k$ for $c\ge1$. We write
\[
 A\in\mathsf P\quad\text{for}\quad
 \exists e,c\ge1,k\ \forall x\ge1\quad
 A(x)=\B(e,x,p_{c,k}(|x|)).
\]
The notation $e,c\ge1$ just requires a valid word code $e$ and positive $c$. This is an ordinary first-order formula containing the fixed function symbol $A$. It agrees extensionally with polynomial-time decidability: timeout means rejection, and the clocked interpreter is itself polynomial time for each fixed descriptor by Theorem~\ref{thm:F06}. Conversely enlarge the constant and degree of any halting polynomial-time algorithm.

For a verifier descriptor $v=(e,c_p,k_p,c_t,k_t)$ put $p(n)=p_{c_p,k_p}(n)$ and $t(n)=p_{c_t,k_t}(n+p(n)+1)$. Its NP bit is
\[
 N(v,x)=\Any_{y<2^{p(n)}}\,B(e,\operatorname{pairWord}(x,\operatorname{word}_{p(n)}(y)),t(n)).
\]
Here all $2^{p(n)}$ witnesses have exactly $p(n)$ bits; for variable-length witnesses use an explicit length field and padding convention. A correct variable-length encoding has a designated canonical padding when counts matter. $A\in\mathsf{NP}$ means $\exists v\forall x\ge1, A(x)=N(v,x)$. Define $\mathsf{coNP}$ by the bit $1-N(v,x)$. These definitions identify both yes and no cases positively; an arbitrary predicate with only one one-sided implication is insufficient.

For space classes use the read-only-input model of Section~\ref{sec:foundations}. Given a cell budget $s$, reject on attempted escape, and simulate for a configuration-count bound $C(e,n,s)$, with halting states absorbing. Call the resulting bit $B_s(e,x)$. Section~\ref{sec:space} proves that this finite semantics agrees with ordinary halting space-bounded decision. With $\ell(n)=\lceil\log_2(n+2)\rceil$, use the budgets in the following table. Constants and degrees are finite descriptors.

\begin{center}
\begin{tabular}{lll}
\toprule
Class & Mode & Budget \\
\midrule
$\mathsf L$, $\mathsf{NL}$ & deterministic, nondeterministic & $c\ell(n)$ cells \\
$\mathsf{PSPACE}$, $\mathsf{NPSPACE}$ & deterministic, nondeterministic & $c(n+1)^k$ cells \\
$\mathsf{EXPSPACE}$, $\mathsf{NEXPSPACE}$ & deterministic, nondeterministic & $c2^{(n+1)^k}$ cells \\
$\mathsf{EXPTIME}$, $\mathsf{NEXPTIME}$ & deterministic, nondeterministic & $c2^{(n+1)^k}$ steps \\
\bottomrule
\end{tabular}
\end{center}
Nondeterministic space acceptance means there is an accepting path within the configuration bound, not that every path halts. It is a bounded existential bit. The budget constructor and uniformity qualifications for general functions $s(n)$ are stated in the space section.

A class inclusion is used in two precise ways. For any fixed language bit it is a first-order implication between the corresponding presentation formulas. More strongly, the proofs give PR transformations of descriptors and a proof of equality of their acceptance bits. For example the numerical sentence expressing $\mathsf P=\mathsf{NP}$ is
\[
 \forall v\ \exists d\ \forall x\ge1\quad N(v,x)=P(d,x),
\]
where $P(d,x)$ is the clocked deterministic acceptance bit. Nothing in this definition quantifies over arbitrary subsets of $\N$.

\begin{theorem}[C01: basic time and space inclusions]\label{thm:C01}
The following transformations are provable in \ipra:
\[
 \mathsf P\subseteq\mathsf{NP}\subseteq\mathsf{PSPACE}
 \subseteq\mathsf{EXPTIME},\qquad
 \mathsf{coNP}\subseteq\mathsf{PSPACE},
\]
and $\mathsf P\subseteq\mathsf{coNP}$. Together with the space results this gives
$\mathsf L\subseteq\mathsf{NL}\subseteq\mathsf P$.
\end{theorem}
\begin{proof}
A deterministic verifier ignores its witness, giving $\mathsf P\subseteq\mathsf{NP}$. Flipping its absorbing output gives a polynomial decider for the complement, hence $\mathsf P\subseteq\mathsf{coNP}$. To decide an NP bit deterministically, enumerate all $p(n)$-bit witness words with a binary counter and run the polynomial-time verifier for each. Save the original input, the current witness, verifier work tape, and an OR bit. A machine running for $t(n)$ steps visits at most $n+p(n)+t(n)+O(1)$ cells, a polynomial. Reuse the verifier storage between witnesses. The invariant says that the OR bit equals the bounded OR over processed witnesses. The complement version flips the final bit and uses the same space.

For a fixed polynomial-space decider, the number of configurations is $2^{\poly(n)}$; a deterministic repeated configuration would repeat forever. Thus the explicit configuration clock decides the same language in exponential time. The count and loop removal are proved in Theorems~\ref{thm:S01}--\ref{thm:S02}. For nondeterministic logarithmic space the configuration graph has polynomially many vertices and its reachability bit is computable in polynomial time, as proved there. All simulations have named clocks, so even the brute-force witness enumeration is PR.
\end{proof}

\begin{theorem}[C02: closure constructions for P]\label{thm:C02}
$\mathsf P$ is closed under complement, finite union, finite intersection, difference, concatenation, and Kleene star. The constructors and correctness proofs are effective in \ipra.
\end{theorem}
\begin{proof}
For Boolean operations run the finitely many deciders and combine their bits. Their clocks add, with polynomial copying overhead. For concatenation, test all $n+1$ cuts $i$ of a word $x$ and OR the tests $A(x[0:i])B(x[i:n])$. Restriction and copying cost polynomial time, and each input to a decider has length at most $n$.

For star, ignore empty factors except in the representation of the empty word. Set $D(0)=1$ and, for $1\le j\le n$, set
\[
 D(j)=\bigvee_{i<j}\bigl(D(i)\land A(x[i:j])\bigr).
\]
Every nonempty factorization of a prefix has a last positive-length factor. Conversely an accepting earlier prefix followed by one accepted factor gives a factorization. The invariant states that $D(j)$ equals existence of a partition into positive-length accepted factors, coded as at most $n+1$ cut positions below $n+1$. Its bound is $(n+1)^{n+1}$, independent of efficiency. It is a bounded formula, so its proof is allowed. The algorithm performs $O(n^2)$ decider calls and uses a length-$n+1$ bitmap; it is polynomial time. Return $D(n)$. No unbounded search over factorizations occurs.
\end{proof}

\begin{theorem}[C03: positive closures of NP and coNP]\label{thm:C03}
$\mathsf{NP}$ is closed under finite union, finite intersection, concatenation, and Kleene star. $\mathsf{coNP}$ is closed under finite union and finite intersection.
\end{theorem}
\begin{proof}
For union, a selector bit chooses which verifier and witness are supplied. For intersection supply both witnesses. For concatenation supply a cut $i\le n$ and the two certificates. For star supply at most $n$ positive-length pieces, their cut positions, and one certificate for each piece. Every piece has length at most $n$, so increasing the witness polynomial if necessary bounds the total certificate by $n\,p(n)+O(n\log(n+2))$. Verification checks the cuts are increasing, the first is zero and last $n$, and runs the appropriate verifier on every piece. The empty input is handled separately. Soundness concatenates the accepted pieces; completeness constructs the list from a bounded factorization. The finite witness search of Lemma~\ref{thm:F01} supplies any needed canonical certificates in the proof, not in the efficient verifier.

The coNP conclusions follow by applying Boolean De Morgan laws to the NP constructions. Since these are complemented \emph{bits}, this is finite Boolean algebra and does not invoke a nonconstructive De Morgan law for arbitrary unbounded predicates. Complement closure for NP itself is not asserted.
\end{proof}

\begin{theorem}[C04: encoding invariance and the function/decision distinction]\label{thm:C04}
If two encodings have polynomial-time translations that preserve and recover the represented object, polynomial-time solvability and NP verification transfer in both directions. A polynomial-time function with polynomially long output has a polynomial-time graph predicate. A polynomial-time graph predicate alone gives a bounded-search function algorithm only when a length bound is supplied; it need not give a polynomial-time function algorithm.
\end{theorem}
\begin{proof}
Compose the input translation with the old decider or verifier and, for output problems, with the output translation. A polynomial-time transducer writes at most polynomially many output symbols. Composition of polynomial bounds is polynomial; the resource and witness constructions are detailed in the reduction section.

For a function $f$, compute $f(x)$ and compare it with the proposed $y$, giving its graph test. In the reverse direction, if $|f(x)|\le p(|x|)$, enumerate all words up to that length and take the first passing graph test. Functionality guarantees uniqueness and the length hypothesis guarantees a witness. This is a PR algorithm using up to $2^{p(n)+1}$ candidates. Nothing in this argument makes that search polynomial. The stronger SAT search-to-decision theorem works because its queries ask whether a \emph{prefix has an extension}; they are not merely equality tests for a completed answer. Thus the source article's informal equivalence between a function and its graph requires this qualification.
\end{proof}

\subsection{Canonical complete bounded-computation problems}
The following are useful examples where the entire completeness argument is elementary. Variable source programs use canonical fixed-alphabet, one-work-tape descriptions; fixed machines in other standard presentations are compiled to this form before reduction. A many-one reduction maps $x$ to $f(x)$ and proves equality of the two language bits. It must also exhibit a transducer resource bound. Definitions and closure under composition are developed in Section~\ref{sec:np}.

\begin{theorem}[C05: unary and binary bounded-time acceptance]\label{thm:C05}
The problem ``does deterministic $e$ accept $x$ within $T$ steps?'' is P-complete under logarithmic-space many-one reductions when $T$ is written in unary, and EXPTIME-complete under polynomial-time many-one reductions when $T$ is written in binary. With nondeterministic $e$, the two versions are respectively NP-complete and NEXPTIME-complete.
\end{theorem}
\begin{proof}
In the unary version the input length $N$ is at least $T$, $|e|$, and $|x|$. Theorem~\ref{thm:F06} gives a polynomial deterministic algorithm. For a fixed source polynomial-time machine $M$ with clock $c(n+1)^k$, output $(M,x,1^{c(n+1)^k})$. A counter up to this polynomial occupies $O(\log(n+2))$ bits; stream the header, the input, and that many unary marks. This is a logarithmic-space transducer, and the acceptance equivalence is the defining simulation equation.

In the binary version $T<2^{N+1}$. A unary countdown can be produced and the simulation run in $2^{O(N)}$ time with a generous fixed polynomial simulation overhead. If $M$ runs in time $c2^{(n+1)^k}$, the binary clock has length $(n+1)^k+O(1)$ and is printed in polynomial time by shifting the fixed binary representation of $c$. This is an EXPTIME reduction.

For nondeterministic sources, guess one branch bit per source step and simulate it. In unary time this is an NP verifier; in binary time it is nondeterministic exponential time. Reducing a fixed source nondeterministic machine uses the same wrapper, and transforms accepting histories in both directions. No universal halting claim is used. Invalid descriptions are rejected. The time and length arguments above also prove all reduction bounds in \ipra.
\end{proof}

\begin{theorem}[C06: unary and binary bounded-space acceptance]\label{thm:C06}
Deterministic acceptance with a supplied space budget $s$, where escape or nonhalting within the configuration bound counts as rejection, is PSPACE-complete for unary $s$ and EXPSPACE-complete for binary $s$, under polynomial-time many-one reductions.
\end{theorem}
\begin{proof}
Use a fixed work alphabet and a fixed number of source tapes, or normalize the machine. Its configuration counter has $O(s+|e|+\log(n+2))$ bits; the simulator stores one configuration and that counter. It is polynomial space in the unary input length and exponential space in the binary input length. A fixed source polynomial-space machine is reduced by outputting its code, the input, and its polynomial budget in unary. An exponential-space source uses its exponential budget in binary, a polynomial-length word. Halting source deciders accept before the configuration bound exactly when their original computation accepts. The budgeted semantics also defines an answer on nonhalting source programs. Soundness, completeness, and the storage bound are the same finite configuration argument used in the space section.
\end{proof}

As a consequence of the explicit separations in Section~\ref{sec:hierarchy}, binary bounded-time acceptance is outside P, and binary bounded-space acceptance is outside PSPACE. The change from unary to binary encoding is substantial: these are not polynomially interconvertible representations of the budget.

\section{Reductions, complete problems, and constructive search}\label{sec:np}

The historical completeness constructions originate with Cook and Karp \cite{Cook1971,Karp1972}. Our 2SAT treatment uses the implication-graph criterion associated with Aspvall, Plass, and Tarjan \cite{AspvallPlassTarjan1979}, while proving only the explicit polynomial bound established below.

The results in this section sharpen the verifier and circuit framework into a
collection of explicit reductions and algorithms. Cook's original paper and
Karp's reduction paper are the historical sources for the completeness
framework; the proofs below give the finite constructions and witness maps
needed here. The implication-graph criterion for 2SAT is the one treated by
Aspvall, Plass, and Tarjan. We use a deliberately simpler polynomial algorithm,
not their sharper linear-time implementation. Bibliographic pointers are
listed in the source ledger accompanying this report.

\subsection{Conventions that affect correctness}

A Boolean instance has an explicitly specified finite variable universe
$0,\ldots,v-1$. Its representation lists those variables, or gives $v$ in
unary, so $v$ is at most the input length. An assignment is a bit vector of
length exactly $v$, including values of unused variables. This detail is
essential for counting. Ordinary syntactic formulas can enforce the same
convention by appending $x_i\vee\neg x_i$ for every designated unused variable.
A literal is a pair $(i,\epsilon)$, where $i<v$ and $\epsilon$ indicates whether
the bit is negated. A clause is a finite list of literals; an empty clause is
false and an empty conjunction is true. In $k$CNF every clause has width at
most $k$. Repeating literals pads nonempty short clauses to width exactly
three without changing assignments. Graphs below are finite, simple,
undirected graphs with numbered vertices and an explicit adjacency table.
Numbers in arithmetic decision problems are binary, and lists retain their
positions even if two entries have equal values.

A verifier presentation consists of a fixed deterministic machine $V$,
nonnegative-coefficient polynomial bounds $q,T$, and the decision bit
\[
 A_V(x)=\bigvee_{y<2^{q(n)}} V_T(x,y),\qquad n=|x|.
\]
Here $y$ is decoded as exactly $q(n)$ bits and $V_T$ is the acceptance bit after
$T(n)$ steps, with absorbing halting states. We can enlarge $T$ so that
$T(n)\ge n+q(n)+1$. There is no unbounded evaluation in this definition.
All instances are validated before processing. A reduction on valid instances
is extended to malformed input by a fixed answer convention and fixed valid
yes/no instances. In particular, the complement of a totalized language also
complements that convention; ``invalid'' is not silently identified with both
answers.

\begin{theorem}[N01: constructive algebra of reductions]\label{thm:N01}
Let $A,B,C$ be represented decision bits. If polynomial-time transducers $f,g$
satisfy $A(x)=B(f(x))$ and $B(z)=C(g(z))$, their composed transducer is a
polynomial-time reduction $A\leq_m^p C$. Identity and composition are explicit,
and polynomial-time decidability is preserved downwards. Logarithmic-space
many-one reductions have the same composition property.
\end{theorem}
\begin{proof}
Suppose $f$ uses at most $p(n)$ steps and outputs at most $p(n)$ bits, while
$g$ uses at most $r(m)$ steps on $m$ bits. A machine runs $f$, copies its output
into a new input area, and runs $g$. A multitape implementation costs at most
$c(p(n)+r(p(n))+n+1)^2$ for a fixed coding constant $c$; this deliberately
loose bound covers the copies and elementary simulation. Polynomial
composition and the displayed pointwise equalities prove the claim. To
combine a reduction with a decider, replace $g$ by that decider. No decision
about whether a machine is polynomial time is used: its bound is part of the
supplied presentation.

For logspace, storing $f(x)$ would be wrong. Instead, simulate $g$ with a
virtual input tape. Its input-head position uses $O(\log n)$ bits, because
$f$ has polynomial output length. Whenever $g$ asks for position $j$, restart
$f$ on $x$, count the emitted symbols to $j$, and provide that symbol. Save
$g$'s work tape while doing this. The sum of the two logarithmic work bounds
is logarithmic. A halting logspace transducer has polynomial time and output
length: ignoring its write-only output tape, it has only polynomially many
configurations, and a repeated configuration would repeat forever. One may
instead supply the resulting polynomial clock explicitly, which is the
preferred internal representation. Thus the recomputation also has a
polynomial bound. The invariant is equality of the simulated control,
work-tape, and virtual-input contents with the ordinary computation of $g$ on
$f(x)$, for all positions below the explicit clock. It is a bounded predicate.
\end{proof}

\begin{theorem}[N02: verifier and branching presentations of NP]\label{thm:N02}
Polynomial-time binary nondeterministic machines and polynomial-length,
polynomial-time verifiers define the same decision bits. The translations
send accepting witnesses to accepting witnesses by primitive-recursive maps.
\end{theorem}
\begin{proof}
For a binary branching machine with time bound $T(n)$, use a string of exactly
$T(n)$ bits to select the branch at each step. An absent transition rejects;
a deterministic transition ignores its bit; a halted computation ignores all
remaining bits. The deterministic verifier simulates these choices for the
supplied clock. An accepting path of length at most $T(n)$ supplies its choice
bits, padded by zeros. Conversely, the simulation supplies the selected
accepting path. The path code and the bit string have elementary bounds in
$T$, and equality of the two current configurations is the induction
invariant. A finite branching degree greater than two is compiled into a
fixed binary choice tree, rejecting unused leaves, at constant-factor cost.

For the other direction, a nondeterministic machine first guesses exactly
$q(n)$ witness bits and then runs $V$ for its supplied time. It accepts exactly
when some witness does. These are existential witness translations; they do
not assert a bijection between all padded paths and witnesses. Such a
bijection is needed only in the counting theorem below, where we fix the
verifier convention from the outset.
\end{proof}

\subsection{Circuit compilation and Cook--Levin}

A circuit is a list of input nodes followed by constant, NOT, binary AND, and
binary OR gates. Every non-input gate refers only to earlier nodes. A selected
node is its output. Evaluate in list order, recording one bit for each node.
After $i$ gates, the recorded prefix is the unique prefix satisfying those
$i$ gate equations for the fixed input bits: the next gate has already
computed arguments and a uniquely determined output. This bounded prefix
invariant proves both existence and uniqueness without an infinitary
valuation object.

\begin{theorem}[N03: uniform finite computation circuits]\label{thm:N03}
For a fixed verifier $V$ and a supplied polynomial clock $T(n)\ge n+q(n)+1$,
there is a polynomial-time constructor $x\mapsto C_x$ with $q(n)$ input bits
such that, for every $q(n)$-bit vector $y$,
\[
 C_x(y)=V_T(x,y).
\]
There is a constant $c_V$ with at most $c_V(T(n)+1)^2$ gates and a polynomial
bit-length bound for $C_x$. The construction is logspace for fixed $V,q,T$.
\end{theorem}
\begin{proof}
We describe the compiler sufficiently explicitly to remove an appeal to a
universal circuit theorem. Suppose $V$ has $k$ tapes, a fixed finite alphabet,
and a fixed finite control. During $T$ steps, every head stays within an
interval containing the initially occupied cells and $T$ additional cells
on each side. Give every tape an interval of length
$L=n+q(n)+2T+3$, with its origin shifted by $T+1$.

At time $t$ use one bit $S_{h,j,a}$ for each tape $h$, cell $j<L$, and alphabet
symbol $a$, one head-position bit $H_{h,j}$, and one bit $Q_r$ for each control
state. Legal encodings are one-hot in the symbol, head, and state positions.
The initial layer consists of constants and the input witness wires, or their
negations, so its one-hot encoding is explicit. For example a witness cell
has symbol-1 bit $y_i$, symbol-0 bit $\neg y_i$, and zero bits for the other
symbols. The fixed input $x$ is hardwired.

For each tape and symbol form the scanned-symbol bit
\[
 U_{h,a}=\bigvee_{j<L}(H_{h,j}\wedge S_{h,j,a}).
\]
Implement this OR by a binary chain. For each transition rule $\rho$, form a
selector $D_\rho$ by ANDing its state bit with its scanned-symbol bits. There
are only constantly many selectors, since $V$ is fixed. The next-state bit is
the OR of selectors entering that state. For a cell-symbol bit use
\[
 S'_{h,j,a}=(\neg H_{h,j}\wedge S_{h,j,a})
 \vee\left(H_{h,j}\wedge
       \bigvee_{\rho\text{ writes }a\text{ on }h}D_\rho\right).
\]
If rule $\rho$ moves tape $h$ by $d_{\rho,h}\in\{-1,0,1\}$, define
\[
 H'_{h,j}=\bigvee_\rho(D_\rho\wedge H_{h,j-d_{\rho,h}}),
\]
with out-of-interval head bits set to zero. The reserved boundary cells ensure
that an actually reached head is never lost during the $T$ layers. Absorbing
halting states are implemented by rules that retain the tape and head and
repeat the state.

Every displayed conjunction or disjunction has constant length except the
$L$-term scanned-symbol chains. Consequently one layer uses at most $c_V L$
gates. Repeat the layer $T$ times and output the accepting-state bit. The
invariant says that every state, head, and cell-symbol bit in layer $t$ is the
corresponding bit of the $t$-step machine configuration, and that all one-hot
conditions hold. The initial layer proves the base case. A one-hot encoding
makes exactly one rule selector true; the three update formulas then give
precisely the tape, head, and state update. Bounded induction on $t\le T$
proves the claim for every witness, the latter quantified below $2^{q(n)}$.

Since $L\le 3T+3$, the gate bound is quadratic. Each node address has
$O(\log(T+2))$ bits, so the encoded circuit length is
$O_V((T+1)^2\log(T+2))$. A simple implementation that repeatedly scans its
partial output still has, for example, a bound $c'_V(T+2)^6$; no unit-cost
large-integer assumption is needed. More efficiently, preallocate a fixed
number of wire slots for each time, tape, cell, and local template. Unused
slots are constant gates. Loop counters and address arithmetic use
$O(\log(T+2))=O(\log(n+2))$ bits, and the bits of $x$ can be reread by an
input scan. This emits the circuit in logspace. The circuit's large values
are gate identifiers of logarithmic bit length, not evaluations of arbitrary
primitive-recursive syntax.
\end{proof}

\begin{theorem}[N04: Tseitin constraints have unique extensions]\label{thm:N04}
There is a linear-size conversion of an acyclic circuit $C$ to a 3CNF formula
$F_C$, with the original inputs retained as designated variables, such that
restriction and gate evaluation give inverse maps between satisfying
assignments of $F_C$ and inputs on which $C$ outputs one.
\end{theorem}
\begin{proof}
Use one variable for every circuit node. For a gate with output $z$, insert
the following clauses:
\[
\begin{array}{c|l}
 z=a\wedge b & (\neg z\vee a),\ (\neg z\vee b),\ (z\vee\neg a\vee\neg b)\\
 z=a\vee b & (z\vee\neg a),\ (z\vee\neg b),\ (\neg z\vee a\vee b)\\
 z=\neg a & (z\vee a),\ (\neg z\vee\neg a)\\
 z=1 & (z)\\
 z=0 & (\neg z).
\end{array}
\]
Append a unit clause asserting the output variable. For the AND line, if
$z=1$, the first two clauses force $a=b=1$; if $z=0$, the last clause says
that at least one input is zero. Conversely, those two cases verify all
three clauses whenever $z=a\wedge b$. The OR proof is dual; the NOT and
constant lines are two-case or one-case checks. These are finite bit case
splits, hence intuitionistically valid.

Gate evaluation extends every accepting input to an assignment satisfying
all clauses. Conversely, induction along the topological gate list shows
that a satisfying assignment has exactly the evaluation bit at every node.
Its output unit clause therefore forces acceptance. That induction also
proves uniqueness of the extension, including gates not connected to the
output. There are at most $3G+1$ clauses for $G$ gates. If the encoding derives
its variable universe solely from literal occurrence, add one tautology per
input; this changes the bound to $3G+q+1$ and prevents loss of unused inputs.
No new choices are introduced by repeating literals to reach width three.
All assignment maps are finite array operations, bounded by $2^{q+G}$.
\end{proof}

\begin{theorem}[N05: SAT, circuit SAT, and 3SAT are NP-complete]\label{thm:N05}
Every polynomial verifier has an explicit polynomial many-one reduction to
3SAT, with primitive-recursive forward and reverse witness maps. SAT and
circuit SAT are NP-complete as well.
\end{theorem}
\begin{proof}
Apply N03 to $V,x$ and N04 to $C_x$. An accepting witness $y$ gives the unique
satisfying assignment by gate evaluation. A satisfying assignment restricts
to $y$ and the simulation invariant gives $V_T(x,y)=1$. Thus
$A_V(x)=\operatorname{SAT}(F_{C_x})$ as equality of Boolean decisions.
The two compilers have polynomial composition. Membership in NP is witnessed
by a bit vector assigning every designated variable: evaluating a formula,
circuit, or CNF is polynomial in its explicit encoding. General SAT receives
the 3CNF formula unchanged, while the circuit target uses $C_x$ directly.
Every claim concerns a supplied verifier presentation; there is no search for
an unknown polynomial bound.
\end{proof}

\begin{theorem}[N06: circuit value is P-complete under logspace reductions]\label{thm:N06}
The problem of evaluating an explicitly supplied acyclic Boolean circuit with
specified input bits is in P and is hard for P under logarithmic-space
many-one reductions.
\end{theorem}
\begin{proof}
Evaluate the nodes in order, recording their bits. Array access by scanning
gives a polynomial tape-time bound, so no random-access assumption is needed.
For a fixed polynomial-time decider $M$ and its clock, apply N03 with no
witness inputs. The resulting constant-input circuit has value $M(x)$. The
last paragraph of N03 gives a logspace constructor. Its correctness is the
same layer invariant, so every presented P language reduces to circuit value.
\end{proof}

\subsection{Graph and arithmetic complete problems}

\begin{theorem}[N07: CLIQUE is NP-complete]\label{thm:N07}
Given a 3CNF $F$ with $m$ clauses, one can construct in polynomial time a graph
$G_F$ such that $F$ is satisfiable exactly when $G_F$ has an $m$-vertex clique.
\end{theorem}
\begin{proof}
Make a vertex for each literal occurrence, remembering both its clause and its
literal. Two vertices are adjacent exactly when they come from different
clauses and their literals are not complements. There are at most $3m$
vertices; a double scan writes the adjacency matrix in polynomial time.
From a satisfying assignment, select the first true literal occurrence in
each clause. Selected vertices are in different clauses and cannot carry
complementary literals, since both are true. They therefore form a clique.

Conversely, a clique contains at most one occurrence from each clause. A
clique of size $m$ therefore contains exactly one from each of the $m$
clauses. Assign a variable the value requested by its selected occurrence,
if there is one, and zero otherwise. Pairwise compatibility ensures that two
selected occurrences never request different values. Every clause then has
a true selected literal. Selection, compatibility checking, and assignment
formation are bounded scans. An empty clause makes an $m$-clique impossible;
if $m=0$, the empty assignment condition and empty clique both satisfy their
respective instances. Membership of CLIQUE in NP follows by checking a list
of $k$ distinct vertices and its at most $k^2$ pairs. Values $k$ larger than
the vertex count are rejected immediately.
\end{proof}

\begin{theorem}[N08: independent set and vertex cover are NP-complete]\label{thm:N08}
For a graph $G$ on $N$ vertices and $0\le k\le N$, the following equivalent
witness transformations are polynomial:
\[
\begin{split}
G\text{ has a }k\text{-clique}
 &\ \Longleftrightarrow\ \overline G\text{ has an independent set of size }k,\\
G\text{ has an independent set of size }k
 &\ \Longleftrightarrow\ G\text{ has a vertex cover of size }N-k.
\end{split}
\]
Consequently INDEPENDENT SET and VERTEX COVER are NP-complete.
\end{theorem}
\begin{proof}
For distinct vertices, complementing the adjacency bit changes ``all selected
pairs are edges'' into ``all selected pairs are nonedges.'' This leaves the
selected vertex list unchanged. For the second equivalence, let $S$ be the
selected independent set and output its vertex complement $C$. An edge
uncovered by $C$ would have both endpoints in $S$, contradicting independence.
Conversely, two adjacent vertices in the complement of a cover would give an
uncovered edge. Cardinalities satisfy $|C|=N-|S|$ by a bounded sum of membership
bits. For the common ``at least $k$'' and ``at most $N-k$'' conventions, trim
an oversized independent set by taking its first $k$ members; enlarge a
smaller cover when necessary. The input bounds make these operations finite.
Both target problems have polynomial verifiers checking their membership
bits and edges. Compose these transformations with N07.
\end{proof}

\begin{theorem}[N09: a parsimonious 3SAT-to-SUBSET-SUM reduction]\label{thm:N09}
For every 3CNF on $v$ designated variables with $m$ clauses, there is a
polynomial-time binary SUBSET-SUM instance with $2v+2m$ indexed positive
integers. Satisfying assignments and subsets summing to the target are in
explicit bijection. In particular, SUBSET SUM is NP-complete.
\end{theorem}
\begin{proof}
Use base eight, with $v$ variable columns and $m$ clause columns. For each
variable $x_i$, create two indexed numbers $a_i^0,a_i^1$. Both have digit one
in variable column $i$ and zero in every other variable column. In clause
column $j$, the digit of $a_i^b$ is the number of occurrences in clause $j$
that are made true by assigning $x_i=b$. This digit is at most three; repeated
literals are counted with their multiplicity. For each clause $j$, create
slack numbers $s_j^1,s_j^2$ with sole nonzero digit respectively one and two
in clause column $j$. The target has digit one in every variable column and
digit four in every clause column.

In any subset of the constructed list, a variable-column sum is at most two.
The total contribution from all variable numbers to a clause column is at
most three, because a literal occurrence contributes to exactly one of the
two numbers for its variable. The two slack numbers contribute at most three
more. Thus every column sum is at most six, strictly below eight. No carries
can occur in any subset sum. Equality with the target can consequently be
checked column by column, a statement proved by repeated quotient and
remainder by eight.

The variable columns force selection of exactly one of $a_i^0,a_i^1$ for
each $i$. These choices specify an assignment. If $t_j$ is the number of true
literal occurrences in clause $j$, then $0\le t_j\le3$, and its target column
requires a slack contribution $4-t_j$. The available slack sums are
$0,1,2,3$, each produced by a unique subset of the one- and two-slack numbers.
Hence the column can be completed exactly when $t_j\ge1$, and the completion
is unique. For $t_j=1,2,3$, select slack of total $3,2,1$, respectively.
This gives both directions and the bijection. If a clause is empty, its
column can never reach four, exactly as required.

Each integer has at most $3(v+m)$ binary bits because base eight is $2^3$.
The whole list has polynomial bit length and is written by nested bounded
loops over variables, clauses, and occurrences. An indexed subset is a bit
vector of length $2v+2m$; binary addition and comparison verify its sum in
polynomial time. Equal-valued entries remain distinct positions, so the
bijection also holds for unused variables and duplicate clauses. For the
zero-column case the list and target are empty and zero; there is one empty
subset and one zero-variable assignment. All nondegenerate generated numbers
are positive.
\end{proof}

\begin{theorem}[N10: PARTITION and 0--1 KNAPSACK are NP-complete]\label{thm:N10}
SUBSET SUM has explicit polynomial reductions to PARTITION and to the
0--1 knapsack decision problem with nonnegative integer weights and profits.
\end{theorem}
\begin{proof}
Let the subset-sum input be indexed numbers $a_1,\ldots,a_r$, target $T$, and
$S=\sum_i a_i$. If $T>S$, output a fixed no-instance. If $S=0$, decide the
instance directly and output a fixed yes- or no-instance. Otherwise
$S>0$ and $0\le T\le S$. Append the numbers $2S-T$ and $S+T$. The total is
$4S$, so equal partition asks for a sublist of sum $2S$. A subset of the
original list of sum $T$, together with $2S-T$, gives such a partition.
Conversely, a sublist summing to $2S$ contains exactly one appended number:
without them its sum is at most $S$, and with both its sum is at least $3S$.
If it contains $2S-T$, its original entries sum to $T$. If it contains $S+T$,
its original entries sum to $S-T$, so the complementary original entries sum
to $T$. The extra numbers have polynomial bit length.

For knapsack, give item $i$ weight and profit both $a_i$, capacity $T$, and
required profit $T$. The inequalities $\sum a_i\le T$ and $\sum a_i\ge T$
force equality. Conversely a subset sum of $T$ meets both constraints.
Checking subset sums, weights, and profits uses polynomial binary arithmetic,
so both target problems lie in NP.
\end{proof}

\begin{theorem}[N11: 0--1 integer feasibility is NP-complete]\label{thm:N11}
Feasibility of a finite system of integer linear inequalities with every
variable constrained to $\{0,1\}$ is NP-complete, even for the inequalities
arising from clauses of width at most three.
\end{theorem}
\begin{proof}
Replace each positive literal $x_i$ by the integer $x_i$ and each negative
literal by $1-x_i$. For each clause require the sum of its literal terms to
be at least one. On a bit assignment this sum counts exactly its true literal
occurrences. Therefore the same assignment satisfies the formula if and only
if it satisfies the system. The expansion has polynomial size, with bounded
coefficients even when a clause repeats a variable. An empty clause becomes
$0\ge1$. Membership in NP follows by guessing the bit vector and evaluating
all integer sums. If coefficients have $b$ bits and there are $r$ terms,
intermediate sums have at most $b+\lceil\log_2(r+1)\rceil+2$ bits, which
supplies the binary-arithmetic bound. This statement does not silently assert
a certificate bound for unrestricted integer variables.
\end{proof}

\begin{theorem}[N12: elementary NP presentations of neighboring problems]\label{thm:N12}
Graph isomorphism, Hamiltonian path, Hamiltonian cycle, and the bounded-factor
problem have explicit polynomial verifiers. The latter asks, on $N\ge2$,
whether $N$ has a divisor $d$ with $2\le d\le b$ and $d<N$.
\end{theorem}
\begin{proof}
For graph isomorphism on $r$ vertices, a witness lists $r$ vertex images.
Check that every image is below $r$, no two are equal, and all $r^2$
adjacency bits are preserved. The first two properties make the map a
permutation: a missing image would leave $r$ distinct images in $r-1$ slots,
contradicting the finite counting lemma. The witness uses
$O(r\log(r+2))$ bits and all checks are polynomial. Graphs of different sizes
are immediately rejected.

A Hamiltonian-path witness is a permutation of the vertices with an edge
between consecutive entries. For a cycle also check the last-to-first edge,
using the chosen convention for graphs of fewer than three vertices. These
verifiers directly express the defining objects. No NP-hardness of either
Hamiltonian problem is inferred merely from membership.

For bounded factor, guess the binary integer $d$, check the inequalities, and
compute $N\bmod d$ by binary long division. Its bit length is bounded by that
of $N$, and repeated shift, comparison, and subtraction give polynomial time.
The condition is equivalent to having a prime factor in the same range:
if a divisor $d\ge2$ is supplied, boundedly select its least divisor
$e\ge2$. A nontrivial factorization of $e$ would give a smaller divisor of
$d$, so $e$ is prime; also $e\le d$ and $e\mid N$. This equivalence does not
require computing $e$ in polynomial time, since the NP witness remains $d$.
It proves membership only, not a polynomial factoring algorithm or a
coNP presentation.
\end{proof}

\subsection{Tractable satisfiability fragments}

\begin{theorem}[N13: constructive 2SAT criterion and polynomial algorithm]\label{thm:N13}
A 2CNF formula without empty clauses is satisfiable if and only if no variable
and its complement are mutually reachable in its implication graph. There is
a polynomial algorithm returning either a satisfying assignment or a
variable together with the two contradictory implication paths.
\end{theorem}
\begin{proof}
Replace a unit clause $a$ by $a\vee a$. For each clause $a\vee b$, put edges
$\neg a\to b$ and $\neg b\to a$ on the $2v$ literal vertices. Every edge
$u\to w$ is paired with $\neg w\to\neg u$. An assignment satisfies the
formula exactly when it respects all implications. Such an assignment
respects every path by induction on its length. If $x_i$ and $\neg x_i$
are mutually reachable, either value of $x_i$ then forces its complement to
have that same truth value. Splitting on the bit of $x_i$ proves failure.

For the converse and an actual algorithm, compute reflexive reachability by
the recurrence
\[
 R_0(u,w)=[u=w]\vee E(u,w),\qquad
 R_{j+1}(u,w)=R_j(u,w)\vee(R_j(u,j)\wedge R_j(j,w)).
\]
The invariant says that $R_j$ describes paths whose intermediate vertices
lie below $j$. Path concatenation proves the forward update. For the reverse
update, erase repeated vertices from a path first; a simple path uses $j$
at most once as an intermediate vertex, and splitting there proves the
recurrence. Length is always bounded by $2v$. Recording a predecessor or
split index at the first successful update reconstructs paths by bounded
loops. Thus the table and any path certificates have explicit elementary
bounds; there is no unbounded reachability search.

Define $u\sim w$ when both $R(u,w)$ and $R(w,u)$ hold. Reflexivity,
symmetry, and transitivity follow from path operations. Represent each class
$C$ by its least literal number $r(C)$. Complementing every literal and
reversing every edge carries a path $u\leadsto w$ to
$\neg w\leadsto\neg u$; consequently complementation induces an involution
$C\mapsto\bar C$ on the classes. By hypothesis $C\ne\bar C$.

Here is an explicit ordering device that avoids an unproved topological-sort
subroutine. Let $d(C)$ count the classes $D$ with a path $D\leadsto C$,
including $C$ itself. Put
\[
 B=2v+1,\qquad K(C)=B\,d(C)+r(C).
\]
Different classes have different keys, since their representatives are below
$B$. If $C\leadsto D$ and $C\ne D$, every predecessor of $C$ is a predecessor
of $D$, while $D$ is not a predecessor of $C$ (otherwise the classes coincide).
Thus $d(D)\ge d(C)+1$, and $K(D)>K(C)$. This is a finite injection and
strict-cardinality argument over the list of representatives.

Assign a literal $u$ the value one exactly when
$K([u])>K([\neg u])$. The unequal keys make complementary literals receive
opposite bits. An edge $u\to w$ gives
$K([u])\le K([w])$ and, from its paired edge,
$K([\neg w])\le K([\neg u])$. If $u$ were true and $w$ false, these
inequalities would yield
\[
 K([\neg u])<K([u])\le K([w])
 <K([\neg w])\le K([\neg u]),
\]
an impossibility. Every implication is therefore respected, and the
assignment satisfies the formula. All quantities come from the finite
reachability table. Floyd's recurrence uses $O(v^3)$ Boolean operations;
representative scans and predecessor counts are polynomial too. Even a
sequential tape implementation with repeated table scans has a polynomial
bound. If there is an empty clause, return its direct inconsistency
certificate before constructing the graph. The case $v=0$ is handled by the
empty conjunction convention.
\end{proof}

\begin{theorem}[N14: Horn satisfiability has a least-model algorithm]\label{thm:N14}
Satisfiability of CNF formulas with at most one positive literal per clause is
polynomial-time decidable. On satisfiable input an explicit algorithm returns
the coordinatewise least satisfying assignment.
\end{theorem}
\begin{proof}
Write each clause as a rule $x_{i_1}\wedge\cdots\wedge x_{i_r}\to h$,
where $h$ is either a variable or false. Repeated antecedents do no harm.
Start with all variables false. In a fixed order scan all rules whose heads
are variables; whenever every antecedent is currently true, set the head
true. Repeat the complete scan $v$ times. A change adds a previously absent
true variable, and no variable is removed. If a full scan makes no change,
all later scans are identical. If the state after $v$ scans were not closed,
each of those scans would have made a change and a further change would be
possible, requiring more than $v$ distinct additions. Finite counting rules
this out. Thus the final state is closed under all positive-head rules.

At every addition, every satisfying assignment must make the newly added
variable true: it already makes all antecedents true by the induction
hypothesis and must respect this rule. Hence the computed true set is a
subset of every model. If any false-head rule now has all antecedents true,
every model would violate that rule, so return unsatisfiable. Otherwise all
rules are satisfied by the computed assignment; it is therefore a model,
and its containment in every model proves leastness. Facts are rules with
empty antecedent; an empty clause is the false-head rule with empty
antecedent and is rejected. At most $v$ scans of the input and polynomial
array accesses are needed. The invariant quantifies over assignments below
$2^v$ only; this makes its minimality proof an instance of bounded induction.
\end{proof}

\subsection{Search, conditional collapse, and counting}

\begin{theorem}[N15: constructive SAT search-to-decision]\label{thm:N15}
A supplied polynomial-time SAT decider yields a polynomial-time algorithm
which returns a satisfying assignment for any satisfiable formula and returns
``none'' for an unsatisfiable one. The same supplied decider yields such a
search algorithm for every polynomial verifier presentation.
\end{theorem}
\begin{proof}
Run the decider once on $F$. If it says no, return none. Otherwise maintain a
prefix assignment $a_0,\ldots,a_{i-1}$ for which the restricted formula is
satisfiable. Query whether this restricted formula with $x_i=0$ is
satisfiable. If yes, append zero. If no, append one. To justify the second
case constructively, take a witness for the currently true finite SAT bit;
its $i$th bit is either zero or one. Zero contradicts the no-answer, so one
supplies a witness for the selected restriction. The invariant is thus
preserved. After $v$ steps all variables are assigned and the restricted
formula's true value means exactly that this assignment satisfies $F$.

There are $v+1$ calls and each restricted encoding has length at most a fixed
polynomial in the original length. One may retain the original formula and
append unit clauses to specify the prefix. If the decider runs in time $p$,
a bound such as $(v+1)p(c(|F|+v+1)^2)+c'(|F|+v+1)^4$ suffices. The algorithm
never searches exponentially for the witness used in the correctness proof;
it uses the decider and finite bit case analysis.

For a verifier $V$, first construct $F_{C_x}$ by N03--N04, find its satisfying
assignment, and restrict to its designated witness-input variables. The
simulation theorem proves that this is an accepting witness. An initial
no-answer correctly certifies that none exists. The resulting algorithm is
a total search procedure with a distinguished failure result, not a claim
that every arbitrary multivalued selection becomes a prescribed function.
\end{proof}

\begin{theorem}[N16: the constructive completeness implication]\label{thm:N16}
A polynomial SAT decider implies a polynomial decider and a polynomial witness
finder for every presented NP problem. Conversely, if every NP presentation
has a polynomial decider, SAT does. Thus the standard statement
$\mathrm P=\mathrm{NP}$ is equivalent to $\mathrm{SAT}\in\mathrm P$ under the
numerical class-presentation convention of this report.
\end{theorem}
\begin{proof}
The forward decision transformation composes the supplied SAT machine with
the two compilers of N05. The witness-finding transformation is N15. These
are explicit functions of the SAT descriptor and the verifier descriptor;
their polynomial bounds are obtained by composition of the supplied
polynomials. The reverse implication specializes the class-equality
hypothesis to the concrete polynomial SAT verifier. The known inclusion
$\mathrm P\subseteq\mathrm{NP}$ supplies the other half of the equality.
These implications neither decide SAT in polynomial time nor settle the
class equality.
\end{proof}

\begin{theorem}[N17: an NP-complete coNP problem collapses NP and coNP]\label{thm:N17}
If a supplied NP-complete problem $K$ also has a coNP presentation, then
$\mathrm{NP}=\mathrm{coNP}$. In particular, $\mathrm P=\mathrm{NP}$ implies
$\mathrm{NP}=\mathrm{coNP}$.
\end{theorem}
\begin{proof}
For any presented NP language $A$, take its supplied completeness reduction
$f$ to $K$. Equality $A(x)=K(f(x))$ gives the equality of complement bits.
The coNP presentation of $K$ means that $\bar K$ has a polynomial verifier.
Compose it with $f$ to obtain a verifier for $\bar A$. Thus complementation
maps every NP presentation to another NP presentation. A coNP presentation
is, by definition, the complement of an NP presentation, so the same
transformation gives both inclusions. If $\mathrm P=\mathrm{NP}$, take
$K=\mathrm{SAT}$: a polynomial decider and its output-bit reversal give P,
hence NP, presentations for both $K$ and $\bar K$.
\end{proof}

\begin{theorem}[N18: parsimonious completeness for witness counting]\label{thm:N18}
For the counting convention
\[
 f_V(x)=\sum_{y<2^{q(|x|)}} [V_T(x,y)=1],
\]
counting satisfying assignments of circuits and of 3CNF formulas is
\#$\mathrm P$-complete under parsimonious polynomial many-one reductions.
Counting indexed subsets of a positive integer list with a specified sum is
also \#$\mathrm P$-complete under those reductions.
\end{theorem}
\begin{proof}
N03 preserves the accepting witness bit strings exactly. N04 gives each such
string one and only one satisfying gate extension, and retains all input
variables. Therefore $f_V(x)$ equals the number of satisfying assignments of
$F_{C_x}$, with no multiplicative factor. Circuit counting has the same
identity before Tseitin conversion. Both counts belong to \#$\mathrm P$
because their finite assignments are verified in polynomial time. Applying
the bijection of N09 gives equality with the number of indexed subsets
summing to its target. This target count is in \#$\mathrm P$ by its binary
subset verifier. Bounded induction on the finite witness lists proves that
inverse bijections preserve cardinalities. This does not infer that exact
counting is polynomial time, and it does not identify arbitrary padded
nondeterministic path counts with the fixed-witness convention.
\end{proof}

\begin{theorem}[N19: a factor-two vertex-cover approximation]\label{thm:N19}
There is a polynomial algorithm that returns a vertex cover $C$ satisfying
$|C|\le2|C^*|$ for every vertex cover $C^*$ of the same graph.
\end{theorem}
\begin{proof}
Scan the edges in a fixed order. If both endpoints are unused, add the edge
to a matching $M$ and mark its endpoints. Output all marked vertices. An
uncovered edge would have had two still-unmarked endpoints when scanned,
because marks only accumulate. It would therefore have been selected, a
contradiction; hence the output is a cover. The selected edges have disjoint
endpoints, so $|C|=2|M|$. Every cover $C^*$ contains at least one endpoint of
each selected edge. Choose the smaller numbered endpoint belonging to $C^*$
on each such edge. Disjointness makes this choice injective, proving
$|M|\le|C^*|$. The algorithm never needs to compute an optimum cover; the
inequality holds for each supplied cover witness and hence for the optimum
obtained, if desired, by finite minimization. A single edge scan with simple
mark updates is polynomial on the explicit graph encoding.
\end{proof}

\begin{theorem}[N20: vertex cover is fixed-parameter tractable]\label{thm:N20}
For a graph on $N$ vertices and parameter $k$, vertex cover of size at most
$k$ can be decided, and a cover produced, in time
$2^{k+1}\poly(N+|E|+1)$.
\end{theorem}
\begin{proof}
If there are no edges, return the empty cover. If an edge remains and $k=0$,
return failure. Otherwise choose the first edge $\{u,v\}$. Every cover
contains $u$ or $v$. Recurse on the graph with $u$ deleted and parameter
$k-1$, and on the graph with $v$ deleted and parameter $k-1$. A successful
branch returns its chosen endpoint together with the recursively produced
cover. If both branches fail, a proposed cover for the original graph would,
by its choice of endpoint, restrict to a cover in a failed branch, a
contradiction. This proves correctness by induction on the remaining depth,
with all covers bounded by $2^N$ and all graphs by their fixed matrix bound.

The recursion tree has at most $1+2+\cdots+2^k<2^{k+1}$ nodes. At each node
one performs a polynomial scan and graph update. It is a primitive-recursive
algorithm: enumerate the full binary tree through depth $k$ and compute
bottom-up, or use an explicit stack with the stated node clock. For $k\ge N$
return all vertices immediately. The parameter $k$ is thereby separated from
the input-size polynomial, which is the fixed-parameter conclusion.
\end{proof}

\begin{theorem}[N21: UNSAT and tautology are coNP-complete]\label{thm:N21}
UNSAT, propositional TAUT, and tautology for DNF formulas with terms of width
at most three are coNP-complete under polynomial many-one reductions.
\end{theorem}
\begin{proof}
The complement of UNSAT is SAT. The complement of TAUT has an NP witness: an
assignment on which the formula is false. These are bounded Boolean
complements, so finite assignment enumeration also identifies them exactly,
without using double-negation elimination for an arbitrary proposition.
For any coNP presentation $A$, its complement has an NP reduction $f$ to 3SAT.
Hence $A(x)=\operatorname{UNSAT}(f(x))$. Negate that 3CNF syntactically and
apply De Morgan laws to obtain a DNF whose terms have width at most three.
For every assignment the resulting DNF has the opposite truth bit, so
$f(x)$ is unsatisfiable exactly when this DNF is a tautology. Negation and
De Morgan conversion have linear output size for a CNF. The bounded universal
assignment test follows from the corresponding pointwise equality by a
finite conjunction fold. This proves the remaining hardness claims.
\end{proof}

\begin{theorem}[N22: Horn SAT is P-complete under logspace reductions]\label{thm:N22}
Horn satisfiability is P-complete, even when every clause has width at most
three. The hardness reduction uses only finite gate-local rules.
\end{theorem}
\begin{proof}
Take an acyclic circuit with its input values specified, as in N06. For each
node $i$ introduce two variables $t_i,f_i$, intended to assert that its value
is respectively true or false. For an input or constant node, add the unit
fact for its specified value. For a NOT gate $z=\neg a$, add the rules
$f_a\to t_z$ and $t_a\to f_z$. For an AND gate $z=a\wedge b$, add
\[
 t_a\wedge t_b\to t_z,\qquad f_a\to f_z,\qquad f_b\to f_z.
\]
For an OR gate $z=a\vee b$, add
\[
 t_a\to t_z,\qquad t_b\to t_z,\qquad f_a\wedge f_b\to f_z.
\]
These are Horn clauses of width at most three. Finally add the negative unit
clause $\neg f_o$, where $o$ is the output node.

Gate-order induction proves that the least closure of the positive-head
rules contains exactly one of $t_i,f_i$ at every node, namely the assertion
of its actual evaluated bit. The base case is the supplied input facts.
For each gate, the displayed rules force its correct assertion from its
already determined input assertions, and none of them triggers its incorrect
assertion. No other rule has that node as head. Every model must contain
this least closure, by the same addition invariant used in N14. If the
circuit outputs zero, $f_o$ is forced and the final negative unit is
impossible. If it outputs one, the least closure itself satisfies every
rule and $\neg f_o$, giving a model. Thus the Horn instance is satisfiable
exactly when the circuit value is one.

A constructor scans the gate list and emits at most three clauses per gate,
plus the input facts and final unit. Variable indices are the two integers
$2i$ and $2i+1$, so its counters and address arithmetic need only logarithmic
work space. Compose with N06 to obtain hardness for P. N14 gives polynomial
membership, and the bounded gate-prefix invariant gives the object-arithmetic
correctness proof. The construction makes no unsupported claim that all
Horn models have exactly one assertion per node; only their least closure
has that property, and that is sufficient for the reduction.
\end{proof}

\subsection{The object-arithmetic obligations in this section}

Each constructor above is a fixed algorithm on finite codes, and each proof
uses finite lists with displayed bounds. To lower the proofs to the repository's
object theory, instantiate the Boolean bounded-formula reflection and bounded
induction interfaces; do not reinterpret a semantic Lean function parameter as
an object-language variable. The especially useful invariants are the
following.

\begin{itemize}
\item For circuit evaluation and Tseitin: every gate below the current prefix
has its computed value; all assignment codes are below $2^{q+G}$.
\item For machine unrolling: all time indices are at most $T$, all tape
positions below $L$, and all witnesses below $2^q$; the complete configuration
array agrees bit by bit with the bounded machine run.
\item For graph reductions: every selected occurrence lies in its recorded
clause; selected literals have consistent requested bits; vertex and subset
codes are bounded by their explicit vertex counts.
\item For subset sum: all digit sums are below eight, exactly one variable
number is selected, and the slack digit equals four minus the true-occurrence
count. Quotient/remainder induction gives the no-carry equality.
\item For 2SAT: the reachability table at stage $j$ matches paths with allowed
intermediate vertices below $j$; representatives, predecessor counts, and keys
are all bounded finite scans.
\item For Horn SAT and SAT search: the true-variable set is contained in each
model below $2^v$, or the current partial assignment has a satisfying extension
below $2^v$, respectively.
\item For parameterized vertex cover: every branch state has a graph below the
original adjacency bound and a chosen-vertex list of length at most $k$; all
branches lie in the full tree with fewer than $2^{k+1}$ nodes.
\end{itemize}

The accompanying \texttt{lean/ReductionsOutline.lean} supplies finite semantic
definitions and precise target statements. It is an uncompiled outline, not a
certificate for these theorems. The Python companion checks the new witness
maps and algorithms against exhaustive small reference computations, including
unused inputs, duplicated literals, empty clauses, and equal-valued indexed
subset entries. Such checks validate implementation cases; the general proofs
are the arguments above.

\section{Space, reachability, complementation, and finite games}
\label{sec:space}

This section gives explicit finite forms of the space-complexity theorems.
The underlying classical results are due to Savitch and, for
complementation, independently to Immerman and Szelepcs\'enyi
\cite{Savitch1970,Immerman1988,Szelepcsenyi1988}.
Their interpretation here is more specific than an appeal to those results:
we give bounded algorithms, witness transformations, and the Boolean
invariants that permit quantifier-free induction.  The presentation of
space-bounded computation and quantified Boolean formulae is also consistent
with the author-hosted text of Arora and Barak \cite{AroraBarak}.

\subsection{Finite graph conventions and path elimination}

A graph has vertices $0,\ldots,N-1$, where $N\geq1$, and an adjacency bit
$E(u,v)$.  An explicit graph stores these $N^2$ bits; an implicit graph
provides a fixed bounded algorithm for the bit.  Throughout, $a,b<N$.
A walk of length $j$ is a list $v_0,\ldots,v_j<N$ such that
$E(v_i,v_{i+1})=1$ for every $i<j$.  Empty walks are allowed.  A list of
$j+1$ vertices can be encoded by $j+1$ base-$\max(2,N)$ digits, with a
separate length field.  Consequently every path-existence assertion with
$j\leq B$ is bounded by an explicitly primitive-recursive number, for
example $(B+2)\max(2,N)^{B+1}$ after a routine tuple encoding.

For a fixed source $a$, define reachability bits by
\[
 r_0(v)=[v=a],\qquad
 r_{i+1}(v)=r_i(v)\vee
       \bigvee_{u<N}\bigl(r_i(u)\wedge E(u,v)\bigr).
 \tag{S.1}
\]
The entire row $(r_i(0),\ldots,r_i(N-1))$ is an $N$-bit integer.
Equation (S.1), evaluated by nested bounded folds, is a primitive recursion
on $i$.  This definition gives a total numerical function even when a
later space-efficient implementation computes its values by a different
algorithm.

\begin{theorem}[S01: bounded reachability and constructive path shortening]
\label{thm:S01}
The bit $r_i(v)$ is one exactly when a walk from $a$ to $v$ of length at
most $i$ exists.  A witnessing walk can be obtained by primitive recursion
and bounded search.  Every walk in an $N$-vertex graph has a walk with the
same endpoints and length at most $N-1$.  Thus $r_{N-1}=r_N$, and their
common value decides unrestricted finite-walk reachability.
\end{theorem}
\begin{proof}
At $i=0$ the sole possible walk is the empty walk at $a$.  If (S.1) is
one at stage $i+1$, either the old bit is one, in which case retain its
walk, or select the least $u<N$ with $r_i(u)E(u,v)=1$ and append $v$ to
the already constructed walk to $u$.  Conversely, a nonempty walk of
length at most $i+1$ has a last predecessor reached in at most $i$
steps; a walk of length zero is retained by the first disjunct.  This
proves both directions with actual witnesses.

For shortening, scan a supplied walk from left to right while maintaining
a list of distinct vertices with the same current endpoint.  When the
next vertex is new, append it.  When it already occurs at position $j$,
delete the suffix after that occurrence.  The surviving prefix is still a
walk: in the repeated-vertex case its last vertex is precisely the vertex
just reached, so the next input edge, if any, remains usable.  A distinct
list of vertices below $N$ has length at most $N$.  One finite proof is to
form its $N$ incidence bits: when a new distinct vertex is inserted, their
sum increases by one, and each incidence bit is at most one.  Therefore
the final distinct walk has at most $N-1$ edges.  Every operation is a
bounded list scan or update.  Shortening a walk of length at most $N$
proves $r_N\leq r_{N-1}$; the reverse inclusion follows directly from
(S.1).  Shortening arbitrary supplied walks proves the final assertion.

For the arithmetic proof, fix an outer bound $B$ on the supplied walk
length.  The invariant says, for every prefix length at most $B$, that
the current list is distinct, consists of vertices below $N$, has the
same endpoint as the input prefix, and has valid consecutive edges.  All
its quantifiers are bounded by $B,N$ and the displayed list-code bound.
Its Boolean compilation is a legitimate quantifier-free induction
formula.  The path-construction invariant is handled in the same way.
\end{proof}

\subsection{Configuration graphs and logarithmic space}

Fix a machine with $q$ control states, $k$ one-sided work tapes, and
work alphabet of size $g$.  Its input tape is read-only.  A space budget
$s\geq1$ permits the first $s$ cells of each work tape; crossing the
budget is made rejecting.  A configuration records the control state,
the input-head position, the $k$ work-head positions and the work-tape
contents.  The number of possible ordinary configurations is at most
\[
 Q=q(n+2)(s+1)^k g^{ks}.                 \tag{S.2}
\]
The harmless $+1$ in each head factor permits a boundary sentinel.
Add a distinguished accepting sink and rejecting sink, and let $\ell$
be an explicitly computed bit length sufficient for every such record.
The vertex set can be padded to $N=2^\ell$: invalid records have no
outgoing edges.  A transition from any accepting configuration to the
accepting sink supplies a unique target.  For this fixed machine,
\[
 \ell\leq C\bigl(s+\lceil\log_2(n+2)\rceil+1\bigr)
 \tag{S.3}
\]
for a fixed natural constant $C$.  We use integer bit length; no real
logarithm is needed in the arithmetic definitions.

\begin{theorem}[S02: the finite configuration theorem]
\label{thm:S02}
Acceptance within the supplied space budget is equivalent to reachability
of the accepting sink in the above graph, and hence to existence of a
computation of fewer than $N$ transitions in that graph.  The
configuration graph and its edge test are primitive-recursive.  For
$s\geq\lceil\log_2(n+2)\rceil$, its vertices and edge test use $O(s)$
working bits.  A deterministic computation which eventually halts has
halted before repeating an ordinary configuration.
\end{theorem}
\begin{proof}
Write down the successive configurations of any accepting computation;
this is a graph walk, to which S01 applies.  Conversely, a graph walk
from the initial record to the accepting sink lists legal transitions,
so deleting the final sink edge gives an accepting computation.  Validity
of each record and each transition is decided by scanning the finite
transition table, checking head movements, and comparing unchanged tape
cells.  A scan needs only the two configuration records and bounded
indices.  Equation (S.3) therefore proves the claimed space bound.

If a deterministic computation repeats a configuration, its future is
repeated as well: induction on the number of further transitions proves
equality of the corresponding configurations.  A first halting time
cannot lie strictly after such a repetition.  This is a finite statement
about a supplied halting trace; it does not invoke a decision procedure
for unbounded halting.  Alternatively one can normalize any deterministic
space-bounded decider by rejecting at the explicit configuration clock.
The proof of equivalence uses its stated totality assumption.
\end{proof}

Whenever a complexity bound $s(n)$ is used operationally below, it is a
fixed primitive-recursive bound together with an $O(s(n))$-space procedure
for obtaining or marking the budget.  Equivalently, the theorem may be
read with an explicit budget parameter.  This qualification avoids silently
assuming that an arbitrary externally described numerical bound can be
computed in the advertised space.  The standard logarithmic, polynomial,
and exponential space bounds have the required constructors.

\begin{theorem}[S03: directed reachability is NL-complete]
\label{thm:S03}
Directed graph reachability is complete for $\mathrm{NL}$ under
logarithmic-space many-one reductions.  In particular,
\[
 \mathrm L\subseteq\mathrm{NL}\subseteq\mathrm P.
\]
\end{theorem}
\begin{proof}
For membership, nondeterministically choose a walk of at most $N-1$
edges, retaining only the current vertex, target, and step counter.
An explicit input graph can be rescanned to test an edge.  These
registers use $O(\log(N+2))$ bits.  S01 proves both acceptance directions.

For hardness, fix a logarithmic-space machine recognizing a language.
In (S.3), $\ell=O(\log(n+2))$, hence $N=2^\ell$ is polynomial in $n$.
Output its configuration graph in a fixed adjacency-matrix encoding.
The transducer runs through the pair of vertex codes $(u,v)$, computes
the corresponding edge bit, and writes it.  It stores only two
$O(\log n)$-bit records and bounded scan indices; it never stores the
whole graph.  The output length $N^2$ is polynomial.  S02 proves the
reduction equation.  Start and accepting vertices are included in the
output in the same space.

Deterministic logarithmic-space computation is a special case of
nondeterministic computation.  For the second containment, an explicit
graph can be solved deterministically by the $N$ iterations of (S.1),
using $O(N^3)$ elementary adjacency operations.  For the polynomial-size
configuration graph just constructed, the resulting time, including
edge generation or explicit graph construction, is polynomial in $n$.
The transducers and their simulation equations provide numerical
compiler witnesses for the two containments.
\end{proof}

\begin{corollary}[S04: even acyclic directed reachability is NL-complete]
\label{thm:S04}
The restriction of directed reachability to directed acyclic graphs is
NL-complete, by a logarithmic-space reduction whose correctness has
explicit path maps.
\end{corollary}
\begin{proof}
Given $G$ with $N$ vertices, form vertices $(v,i)$ for $v<N$ and
$0\leq i<N$.  Add $(u,i)\to(v,i+1)$ exactly when $i<N-1$ and either
$u=v$ or $E(u,v)=1$.  Every edge increases the second coordinate, so
no directed cycle is possible.  A path from $a$ to $b$ of length at
most $N-1$ can be padded by stays to give a path from $(a,0)$ to
$(b,N-1)$.  Conversely erase stay edges in a layered path.  For $N=1$
this is just the length-zero test.  The two bounded counters and a
rescan of $G$ give a logarithmic-space transducer.  Membership follows
from S03, because the general reachability algorithm applies.
\end{proof}

\subsection{Savitch's recursion and its actual workspace}

For $a,b<N$ define Boolean matrices $R_j$ by
\[
 \begin{split}
 R_0(a,b)&=[a=b]\vee E(a,b),\\
 R_{j+1}(a,b)&=\bigvee_{z<N}
                 \bigl(R_j(a,z)\wedge R_j(z,b)\bigr).
 \end{split}                                             \tag{S.4}
\]
As a numerical definition this is primitive recursion on the whole
$N^2$-bit matrix.  This observation establishes primitive-recursive
status; storing that matrix is \emph{not} the space-efficient algorithm.

\begin{theorem}[S05: Savitch's finite theorem]
\label{thm:S05}
The bit $R_j(a,b)$ is one exactly when there is a walk from $a$ to $b$
of length at most $2^j$.  There is a deterministic implementation of
$R_j(a,b)$ using
\[
 O\bigl((j+1)(\lceil\log_2(N+2)\rceil+
                    \lceil\log_2(j+2)\rceil)+w\bigr)
 \tag{S.5}
\]
working bits when the edge test uses $w$ working bits.  Its time is
bounded by an explicit primitive-recursive function; for a polynomial-time
edge test one may use
$C(2N+2)^{j+2}(\ell+j+2)^C$ with a suitable fixed $C$.
\end{theorem}
\begin{proof}
The assertion at $j=0$ is the definition of a walk of length at most
one.  Concatenate the two walks certified by a successful midpoint in
(S.4) to prove soundness at $j+1$.  Conversely split a walk after
$\min(2^j,\text{its length})$ edges.  Both pieces have length at most
$2^j$, so their common midpoint makes (S.4) true.  These operations
are bounded list transformations.  If a witness is wanted, search for
the least successful midpoint and recursively concatenate the witnesses.

For the space claim, implement the recurrence by depth-first evaluation.
A frame stores its endpoints, a midpoint counter, its level, and a
constant-size phase indicating whether it is trying a new midpoint or
returning from one of the two child calls.  If the first child is false,
advance the midpoint.  If it is true, evaluate the second child.  Return
true on the first pair of successful children, and false after the last
midpoint.  A parent need not store either child's recursion stack after
that child returns.  There are at most $j+1$ active frames.  Each has
$O(\log(N+2)+\log(j+2))$ bits; edge-test scratch space can be reused
at the leaves.  This proves (S.5).

Every frame has at most $2N$ child calls.  Bounding all control and
edge-test costs by the displayed polynomial gives the stated coarse
time bound by induction.  For arbitrary specified bounded edge code,
replace that polynomial by its supplied primitive-recursive clock.
The explicit time bound also allows the stack machine to be defined by
clocked primitive recursion, without using general well-founded recursion
as an unexamined arithmetic principle.

For the iPRA induction, fix an outer maximum $J$.  At levels $j\leq J$,
all walk witnesses are bounded by a base-$\max(2,N)$ code of length
$2^J+1$.  Thus the invariant universally quantifying over vertices,
levels and witnesses is bounded.  The stack-correctness invariant is
bounded by the displayed time and stack-code bounds.  Both compile to
Boolean equations for quantifier-free induction.
\end{proof}

\begin{corollary}[S06: the principal space equalities]
\label{thm:S06}
For the explicit budget conventions above and $s(n)\geq\log_2(n+2)$,
\[
 \mathrm{NSPACE}(s)\subseteq\mathrm{DSPACE}(s^2).
\]
Consequently,
\[
 \begin{gathered}
 \mathrm{NL}\subseteq\mathrm{DSPACE}(\log^2 n),\qquad
 \mathrm{NPSPACE}=\mathrm{PSPACE},\\
 \mathrm{NEXPSPACE}=\mathrm{EXPSPACE},\qquad
 \mathrm{PSPACE}\subseteq\mathrm{EXPTIME}.
 \end{gathered}
\]
\end{corollary}
\begin{proof}
Use the configuration graph with $N=2^\ell$ and evaluate
$R_\ell(a,b)$.  S01 says that this reaches every possible accepting
configuration, while (S.3) and (S.5) give $O(s^2)$ working bits.
The logarithmic corollary follows by substitution.  The square of a
polynomial is a polynomial.  The square of $2^{p(n)}$ is $2^{2p(n)}$,
which is still an exponential-space bound.  The reverse inclusions in
the two equalities hold because deterministic machines are special
nondeterministic machines.  Finally a deterministic polynomial-space
decider has only $2^{\poly(n)}$ configurations and cannot repeat one
before halting, by S02; simulation up to that clock is exponential time.
The same normalization also yields
$\mathrm{DSPACE}(s)\subseteq\mathrm{DTIME}(2^{O(s)})$ for the stated
fixed machines and budget constructors.
\end{proof}

\subsection{Inductive counting without negative guesses}

Write $A_i=\{v<N:r_i(v)=1\}$ and $c_i=\sum_{v<N}r_i(v)$.  These are
notation for a bit table and its bounded sum, not set-existence axioms.
A guessed path of length at most $i$ is checked online: retain its
current vertex and counter, guess the next vertex, check its edge, and
at a guessed stopping time require the desired endpoint.  Every branch
is stopped by the bound $i$.

\begin{lemma}[S07: forced enumeration]
\label{thm:S07}
Suppose the exact number $c_i$ is supplied.  There is a nondeterministic
$O(\log(N+2))$-space procedure which, on a requested $v$, certifies
$v\notin A_{i+1}$.  Every completed successful branch certifies that
fact, and if the fact holds a successful branch can be constructed by a
primitive-recursive procedure.  Here $0\leq i<N$.
\end{lemma}
\begin{proof}
Scan $u=0,\ldots,N-1$ in order, keeping a counter $d=0$.
At each $u$ guess whether to select it.  A selected vertex must be
accompanied by a checked walk from $a$ of length at most $i$; also check
$u\ne v$ and $E(u,v)=0$, then increase $d$.  An unselected vertex
requires no claim.  At the end require $d=c_i$.

Every selected vertex belongs to $A_i$, and no vertex is selected twice.
A subset of the finite $0/1$ table $r_i$ with the same sum $c_i$ equals
that table: if an omitted table entry were one, summing the componentwise
inequalities would give $d<c_i$.  Hence a successful branch has checked
$u\ne v$ and $E(u,v)=0$ for every $u\in A_i$.  Formula (S.1) now says
$v\notin A_{i+1}$.

Conversely, if $v\notin A_{i+1}$, select exactly the vertices in
$A_i$, and for each select the canonical path constructed in S01.
All checks pass, and exactly $c_i$ vertices are selected.  Computing this
successful transcript may use large temporary tables; the recognizing
machine only guesses it and checks paths one at a time.  This distinction
is essential.  The recognizing machine stores $i,v,u,d,c_i$ and a few
path registers, all logarithmic in $N$; it does not store $A_i$.
\end{proof}

\begin{theorem}[S08: constructive space complementation]
\label{thm:S08}
Nonreachability has a nondeterministic logarithmic-space recognizer whose
accepting transcripts are sound and whose successful transcript can be
constructed primitively recursively whenever nonreachability holds.
For the budgeted machine presentations above,
\[
 \mathrm{NSPACE}(s)=\mathrm{coNSPACE}(s)
 \quad(s\geq\log_2(n+2)),
 \qquad \mathrm{NL}=\mathrm{coNL}.
\]
These classes are closed under finite unions, intersections, differences,
and symmetric differences.
\end{theorem}
\begin{proof}
Start from $c_0=1$.  Given the correct $c_i$, compute a candidate
$c_{i+1}$ by scanning every $v<N$.  Guess either ``reachable'' or
``not reachable.''  In the first case verify a path of length at most
$i+1$ and increment the new counter.  In the second case run S07.
All unsuccessful guesses reject the branch.  Every branch completing
the scan has classified each vertex correctly, so its new count equals
$c_{i+1}$.  At least one branch completes: use the actual finite bits
$r_{i+1}(v)$, canonical positive paths, and the successful transcripts
from S07.  Induction on $i=0,\ldots,N-1$ yields the exact $c_N$ on
every branch which survives all scans.

For the final target $b$, enumerate vertices once more, selecting only
vertices accompanied by a path of length at most $N$, requiring every
selected vertex to differ from $b$, and requiring exactly $c_N$
selections.  The equal-cardinality argument forces enumeration of all
of $A_N$.  Therefore the final test can succeed exactly when
$r_N(b)=0$, which by S01 is nonreachability.  It never accepts a bare
unsupported assertion of nonreachability.

There are at most $N$ stages, $N$ queried vertices per stage, $N$
forced-enumeration vertices per negative query, and $N$ edges per checked
path.  Thus $O(N^4)$ edge tests suffice on every branch, apart from
polynomial control overhead.  Only a constant number of indices, counts,
and current vertices are live at once.  This proves logarithmic space
for explicit graphs.  For the configuration graph, every vertex has
$\ell=O(s)$ bits and its edge test has $O(s)$ workspace; replacing an
explicit adjacency lookup by that test proves the general statement.
The accepting transcript bound is exponential in a primitive-recursive
bound and remains primitive recursive.

The iPRA induction invariant is stronger than ``the count is correct''
on one favored computation.  It states that every transcript prefix
within the explicit branch-code bound which completes stage $i$ has
count $\sum_{v<N}r_i(v)$, and that the canonical transcript constructor
completes that stage.  All transcript, vertex, path and stage quantifiers
are bounded.  Its Boolean compilation justifies the induction and
produces positive witnesses for completeness.

For closure under union, first choose which machine to run.  For
intersection, run the two normalized recognizers sequentially and
accept only if both accept, reusing their workspace.  Given accepting
traces for both, concatenation gives an accepting combined trace;
projection gives the converse.  Complementation, intersection and union
then construct differences and symmetric differences.  Reapplying the
complement compiler gives the reverse class inclusion, so the displayed
equality is constructive at the level of presentation compilers.
\end{proof}

\begin{corollary}[S09: 2SAT is NL-complete]
\label{thm:S09}
Satisfiability of formulae with clauses of size at most two is
NL-complete under logarithmic-space reductions.  Its restriction to
clauses of the forms $\neg x_u\vee x_v$ and unit clauses is already
NL-hard.
\end{corollary}
\begin{proof}
The constructive criterion in Theorem~\ref{thm:N13} says that a 2CNF
formula is unsatisfiable precisely when some literal and its negation
are mutually reachable in its implication graph.  Guess the variable
and the two bounded paths.  This is an NL recognizer for 2CNF
unsatisfiability, since the implication graph's edge bits can be
computed by rescanning the formula.  S08 yields an NL recognizer for
satisfiability.

For hardness, transform a graph with source $a$ and target $b$ into
\[
 F_G=x_a\ \wedge\ \neg x_b\ \wedge\!
              \bigwedge_{E(u,v)=1}(\neg x_u\vee x_v).
 \tag{S.6}
\]
If $a$ reaches $b$, the implication clauses propagate truth from $x_a$
to $x_b$, contradicting the unit clause $\neg x_b$.  If $a$ does not
reach $b$, the assignment $x_v=r_N(v)$ satisfies every clause: a true
reachable tail of an edge has a reachable head, $x_a=1$, and $x_b=0$.
Thus nonreachability reduces to 2SAT with explicit witness extraction.
The formula can be streamed with two vertex counters.  S03 and S08
make nonreachability NL-hard: reduce the complement of an arbitrary NL
language to reachability and then complement the resulting equality.
All complements here are complements of specified Boolean predicates.
This proves the result without an unbounded excluded-middle step.
\end{proof}

\subsection{Quantified Boolean formulae with controlled syntax size}

A Boolean formula uses variables, negation, conjunction, disjunction,
and the finite quantifiers $\exists z\in\{0,1\}$ and
$\forall z\in\{0,1\}$.  Evaluation of these quantifiers means respectively
Boolean disjunction and conjunction of the two evaluations.  A closed
formula has a determinate evaluation bit by a bounded recursive traversal
of its finite syntax.  Malformed syntax is assigned value zero and
excluded by a validity bit.  The language TQBF consists of valid closed
formulae with value one.

\begin{theorem}[S10: bounded evaluation of QBF]
\label{thm:S10}
There is a primitive-recursive evaluator for closed Boolean formulae,
and a deterministic polynomial-space implementation.  Truth has an
explicit finite winning strategy for existential choices; falsity has
an explicit finite counterstrategy for universal choices.
\end{theorem}
\begin{proof}
Use depth-first syntax evaluation.  Store the current assignment,
a syntax-position stack, and a phase bit at each quantifier.  At an
existential quantifier evaluate both assignments in order until one
is true; at a universal quantifier require both to be true.  After a
child returns, restore the old value of the bound variable.  For syntax
length $m$, there are at most $m$ stack frames; positions and variable
indices use $O(\log(m+2))$ bits, and an assignment uses at most $m$
bits.  Even a deliberately wasteful implementation copying the assignment
at each frame uses $O(m^2+m\log(m+2))$ bits.  A coarse clock
$2^{Cm}(m+2)^C$ suffices: a branch has depth at most $m$, every node
has at most two children, and each local operation is polynomial.
Running this implementation to that clock is a fixed primitive recursion.

For a true existential node choose the least child with value one;
for a true universal node retain both true children.  At a false
universal node choose the least false child, and at a false existential
node retain both false children.  The recursive construction is bounded
by the evaluation tree.  Induction on its finite depth proves that the
resulting strategy wins against every opposing sequence of choices.
This supplies witnesses and counterwitnesses without invoking an
infinite game-determinacy principle.
\end{proof}

The next construction is needed because writing down both halves of
(S.4) recursively would produce an exponentially long formula.
Let $U,V,M,A,B$ denote vectors of $\ell$ Boolean variables, and let
$D_x(U,V)$ be a Boolean formula for equality or one configuration step
of the fixed machine on input $x$.  Such a formula has polynomial size
in $n+s$.  Explicitly, enumerate its finitely many transition rules
and all possible head positions, require the stated control-state and
scanned-symbol bits, prescribe the changed head and symbol fields,
and require equality of every untouched cell.  The number of work
heads and the transition table are fixed, so this enumeration is
polynomial.  Input symbols are the fixed bits of $x$.  Validity checks
exclude malformed records; sink transitions are added explicitly.

Define formulae with free endpoint vectors by
\[
 \begin{split}
 \Theta_0(U,V)&=D_x(U,V),\\
 \Theta_{j+1}(U,V)
 &=\exists M\ \forall z\ \exists A\ \exists B\,
       \bigl(H_z(U,V,M,A,B)\wedge\Theta_j(A,B)\bigr),
 \end{split}                                               \tag{S.7}
\]
where every displayed vector quantifier abbreviates $\ell$ Boolean
quantifiers, all newly bound variables are fresh, and
\[
 H_z=\bigl(\neg z\to(A=U\wedge B=M)\bigr)\wedge
        \bigl(z\to(A=M\wedge B=V)\bigr).
 \tag{S.8}
\]
Vector equality is a conjunction of bit equalities; implication in
(S.8) is the ordinary Boolean connective $\neg P\vee Q$.

\begin{theorem}[S11: polynomial-size quantified reachability]
\label{thm:S11}
For all endpoint assignments, $\Theta_j(U,V)$ has value $R_j(U,V)$.
Its syntax has one copy of $D_x$ and only $O(j\ell)$ extra connective
and variable occurrences.  With binary variable indices its encoded
length is bounded by
\[
 O\bigl((|D_x|+j\ell)\log((j+1)\ell+2)\bigr),
 \tag{S.9}
\]
where $|D_x|$ counts syntax nodes before variable-name encoding.  Construction,
renaming, and substitution of constant endpoints take polynomial time
in the output bound.
\end{theorem}
\begin{proof}
Fix $M$.  At $z=0$, the constraints in (S.8) force $A=U,B=M$;
at $z=1$ they force $A=M,B=V$.  These assignments exist explicitly
and are unique.  Thus the two universal branches are true exactly
when both $\Theta_j(U,M)$ and $\Theta_j(M,V)$ are true.
The outer existential quantifier then takes the disjunction over
$M<N$, exactly as in (S.4).  This proves the recurrence by Boolean
calculation, and induction proves the evaluation identity.  In
particular the midpoint must be chosen \emph{before} the universal
selector; choosing it separately after the selector would invalidate
the proof.

There is only one textual occurrence of $\Theta_j$ on the right of
(S.7).  Each level adds three vector blocks, one selector, and a
constant number of vector equalities.  These have $O(\ell)$ nodes.
The total number of fresh variables is $O((j+1)\ell)$, so variable names
cost $O(\log((j+1)\ell+2))$ bits.  This factor applies to the base
formula too, because its endpoint variables are renamed at the deepest
level.  This proves (S.9).
A bounded traversal renames free endpoint variables without expanding
or copying any recursive subformula.  All loops have polynomial bounds.
If a prenex convention is desired, pull the fresh inner quantifiers
outward across the conjunctions in (S.7); their variables are absent
from the corresponding guard.  Both directions of each transformation
are verified by the two Boolean quantifier cases.  This introduces no
duplication and no superpolynomial growth.
\end{proof}

\begin{theorem}[S12: TQBF is PSPACE-complete]
\label{thm:S12}
Every fixed polynomial-space machine presentation has a polynomial-time
many-one reduction to TQBF, with an iPRA proof of its numerical
correctness equation.  Together with S10 this proves PSPACE completeness.
\end{theorem}
\begin{proof}
For the polynomial-space machine and input $x$, compute its polynomial
space bound, configuration length $\ell$, initial record $a$, accepting
sink $b$, and the polynomial-size formula $D_x$.  Output the closed
formula $\Theta_\ell(a,b)$.  S11 says its value is $R_\ell(a,b)$;
S01 and S02 say that this is exactly the input's acceptance bit.
Here $\ell=O(\poly(n))$, so (S.9) is polynomial in $n$ and the
construction is polynomial time.  Every compiler loop, equality test,
and syntax substitution has an explicit polynomial bound.  Numerical
correctness is obtained by the bounded invariants already proved;
this is the stated non-formalized iPRA proof, not a claim that a Lean
object-derivation has already been generated.  S10 proves membership.
\end{proof}

\subsection{Succinct games with a polynomial horizon}

A useful immediate consequence is a precise game-complexity theorem.
An instance supplies a state width $r$, Boolean circuits
$\operatorname{owner}(q)$, $\operatorname{next}(q,b)$ and
$\operatorname{win}(q)$, an initial $r$-bit state, and a horizon $h$
in unary.  Players choose one bit per move.  The owner circuit says
whether the existential or universal player chooses, the next circuit
updates the state, and after exactly $h$ moves the existential player
wins iff $\operatorname{win}(q)=1$.  Earlier terminal positions can be
made absorbing.  Circuits are explicit finite acyclic circuits; their
outputs can be evaluated with polynomial workspace.

\begin{theorem}[S13: finite strategies and game completeness]
\label{thm:S13}
The winner of a supplied game is decided by a primitive-recursive
Boolean function.  One of the players has a canonical finite winning
strategy.  Deciding whether the existential player wins is PSPACE-complete
under polynomial-time reductions.
\end{theorem}
\begin{proof}
Write $W_0(q)=\operatorname{win}(q)$.  Given $W_i$, put
\[
 W_{i+1}(q)=
 \begin{cases}
 W_i(\operatorname{next}(q,0))\vee
 W_i(\operatorname{next}(q,1)),&\operatorname{owner}(q)=1,\\
 W_i(\operatorname{next}(q,0))\wedge
 W_i(\operatorname{next}(q,1)),&\operatorname{owner}(q)=0.
 \end{cases}
 \tag{S.10}
\]
As in (S.1), primitive recursion on the whole $2^r$-bit table gives
a total numerical definition.  Alternatively, a clocked depth-first
evaluation uses at most $h+1$ frames, each holding an $r$-bit state,
a move index and a return bit, plus the reusable circuit workspace.
Its space is polynomial in the input size because $h$ is unary.
Its time is at most $2^{h+1}$ times a polynomial local cost.

The winning-strategy construction is the one in S10: choose the least
winning child at one's own winning position and retain both children
at the opponent's position; reverse the roles at a losing position.
The entire tree has fewer than $2^{h+1}$ nodes and a primitive-recursive
code bound, for example $2^{(r+4)2^{h+1}}$ with a fixed tree encoding.
Backward induction proves that the chosen strategy wins against every
opposing play.  Bound all states, histories and strategy codes by
these displayed numbers, and the invariant is a bounded arithmetic
formula available to iPRA.

For hardness take a prenex QBF from S12.  A game state stores the
current variable index and the assignment prefix.  The owner is the
quantifier type at that index; the move sets that variable, and the
terminal circuit evaluates the quantifier-free matrix.  Its width,
transition circuits, and horizon are polynomial in the formula length.
Induction over the remaining quantifier prefix identifies game winning
with QBF truth.  This gives the reduction.  The succinct representation
is essential: the theorem concerns a polynomial-size description of
an exponentially large game tree, not an explicitly listed tree.
\end{proof}

\subsection{Formalization obligations particular to this section}

The companion \texttt{lean/SpaceOutline.lean} fixes numerical graph
and QBF data and the total recurrences.  It is an uncompiled outline,
not an asserted completion of these complexity proofs.  The shortest
route to a complete formalization has the following dependencies.

\begin{enumerate}
\item Establish bit-table access, bounded folds, finite path encoding,
      and list deletion at the numerical level.  Prove S01 using a
      bounded path-code invariant, then compile that invariant with
      ProveIt's bounded-formula reflection rather than using an
      induction principle for arbitrary formulas.
\item Define fixed-machine configuration packing, validity, the explicit
      accepting sink, and the local edge predicate.  Prove both trace
      maps in S02, including malformed-code behavior.  The machine code
      and its explicit clock are inputs to a bounded interpreter; an
      unclocked universal evaluator for primitive-recursive syntax is
      never needed.
\item Prove equality of the matrix recurrence (S.4) with the stack
      implementation.  A theorem about a large matrix alone does not
      prove Savitch's workspace bound.  Formalize the live-frame bound
      and the time clock separately.
\item Define the forced-enumeration transcript checker and the canonical
      successful-transcript constructor.  Prove the two implications in
      S07, and carry both universal soundness and constructed completeness
      through the stage induction of S08.  This is where an unjustified
      appeal to ``guessing nonreachability'' must be excluded.
\item Define fresh-variable generation, vector equality and the single-copy
      construction (S.7).  Prove its evaluation recurrence and its syntax
      count independently.  Prove the edge-formula compiler correct before
      combining these facts into S12.
\item Finally construct the actual object-language formulas and proof
      terms.  The target is ProveIt's intuitionistic derivability type
      \texttt{([] : List (Fm PTm 0))} followed by the
      \texttt{PRAAx} turnstile and the coded sentence.  Ambient Lean
      theorems about finite lists or functions are useful semantic
      checks, but do not by themselves inhabit that object-theory type.
\end{enumerate}

\section{Explicit separation witnesses and coarse hierarchies}
\label{sec:hierarchy}

The familiar hierarchy theorems originate in diagonal simulations \cite{HartmanisStearns1965,HennieStearns1966}. We prove concrete consequences with generous resource gaps. The advantage for \ipra\ is that every proposed smaller-resource machine receives an explicit mismatch input. We never infer $\exists x\,A(x)\ne B(x)$ merely from $\neg\forall x\,A(x)=B(x)$.

\subsection{Self-descriptions with numerical bounds}
For a program word $e$, a positive coefficient $c$, and an exponent $k$, form
\[
 W(e,c,k)=1^{|e|}\,0\,e\,1^c\,0\,1^k\,0.
\]
The initial unary field announces the exact program length, so zeros in the program do not confuse the next two unary fields. The last zero must be the last input symbol. The parser and constructor are PR, and direct concatenation equations prove
\[
 \operatorname{parse}(W(e,c,k))=(e,c,k),\qquad
 n=|W(e,c,k)|=2|e|+c+k+3.
\]
Thus $|e|,c,k\le n$. These unary descriptors are used only as deliberately padded diagonal inputs. They are not the efficient serializers used in reductions.

\begin{theorem}[H01: P is properly contained in EXPTIME]\label{thm:H01}
There is a fixed PR bit $D$ computed by a fixed machine in time $2^{K(n+1)^2}$, for a fixed constant $K$, such that \ipra\ proves, uniformly in $e,c\ge1,k$,
\[
 D(W(e,c,k))\ne
 \B\bigl(e,W(e,c,k),c(|W(e,c,k)|+1)^k\bigr).
\]
Consequently no polynomial-time presentation decides $D$, and $\mathsf P\subsetneq\mathsf{EXPTIME}$.
\end{theorem}
\begin{proof}
On an invalid self-description return zero. Otherwise recover $e,c,k$, set $n=|x|$, and compute
\[
 D(x)=1-\B(e,x,c(n+1)^k).
\]
Every operation is PR. At $x=W(e,c,k)$ the parser equation makes $D(x)$ exactly the complement of the candidate's bit, so the mismatch follows by the two Boolean cases. No halting promise for $e$ is needed.

It remains to give a time bound for the algorithm defining $D$, not merely to show it is PR. Write $T=c(n+1)^k$. For $n\ge1$, quantifier-free induction gives $n+1\le2^n$, and hence
\[
 T\le n(n+1)^n\le2^{n+n^2}\le2^{(n+1)^2}.
\]
Make a unary countdown by starting with $c$ marks and copying a block $n+1$ times in each of $k$ stages. All intermediate blocks have at most $T$ marks, and there are at most $n$ stages. Repeated full scans give the coarse bound $C_0(n+T+1)^4$ for this production. Theorem~\ref{thm:F06} gives the same form for the universal simulation. Thus the entire computation costs at most $C_1(n+T+1)^4$. With $q=(n+1)^2$, this is at most $16C_1 2^{4q}\le2^{Kq}$ for $K=4+16C_1$, enlarged for malformed and short inputs. These are elementary inequalities, with no eventual-growth premise. This clock lies in EXPTIME after increasing the fixed coefficient and degree.

The parsing, copying, and simulation invariants are those of the foundation section, bounded by the displayed terms. ProveIt already contains object-level arithmetic derivations of the two clock estimates, \texttt{DiagonalClock.clock\_exp} and \texttt{cost\_exp} \cite{ProveItDiagonal}. The numerical simulator and complete diagonal theorem still require the formalization work described later; the availability of those arithmetic lemmas is not a claim that the full theorem has been kernel checked.
\end{proof}

\begin{theorem}[H02: a coarse fixed-degree time hierarchy]\label{thm:H02}
For every fixed integer $k\ge1$, on the canonical model with read-only input and one work tape,
\[
 \mathsf{DTIME}_{\mathrm{can}}((n+1)^k)
 \subsetneq
 \mathsf{DTIME}_{\mathrm{can}}((n+1)^{4k+4}).
\]
There is a PR counterexample word for each proposed machine and constant in the lower class.
\end{theorem}
\begin{proof}
Use only descriptions $W(e,c,0)$, and define $D_k$ by complementing the bounded computation with clock $c(n+1)^k$. The fixed $k$ is part of the constructed machine, not read from its input. As $c\le n$, its simulated time is at most $(n+1)^{k+1}$. Unary countdown construction and the fourth-power interpreter bound cost $O((n+1)^{4k+4})$. On $W(e,c,0)$ the output differs from the candidate within that exact clock. The nonstrict inclusion follows simply from monotonicity of the clocks, with constants adjusted. This proves the displayed coarse hierarchy, not the optimized standard $t\log t$ time-hierarchy theorem. For P versus EXPTIME the difference between one and several tapes disappears by polynomial simulation.
\end{proof}

\subsection{Space-bounded diagonalization}
Normalize machines to one one-sided work tape over a fixed alphabet of size $g\ge2$, with a separate read-only input tape. A fixed machine with several tapes or another finite alphabet can be encoded in fixed-length blocks with head-marker bits. Folding a two-sided tape into two tracks makes it one-sided. The encoding uses at most a fixed constant times its old work space, plus constant bookkeeping; this is absorbed into the coefficient $c$. Sweeps may increase time, which is irrelevant here.

For a supplied budget $s\ge1$, the configurations of a normalized machine with $q$ nonhalting states are bounded in number by
\[
 C(e,n,s)=(q+2)(n+2)(s+1)g^s.
\]
The factors record state, input head, work head, and tape content. Reject attempts to use a cell outside the budget. Run for $C$ transitions, retaining absorbing halts, and denote the acceptance bit by $B_s(e,x)$. A binary counter to $C$ uses $O(\log q+\log(n+2)+s)$ bits. Thus the whole finite simulation uses $O(s+|e|+\log(n+2))$ space even if its time is enormous. Repeated configurations cannot precede the first accepting configuration in a deterministic trace; the finite pigeonhole argument therefore proves equivalence to eventual acceptance under the budget. This is the deterministic specialization of the space section's configuration lemma.

\begin{theorem}[H03: L is properly contained in PSPACE]\label{thm:H03}
There is a PR language $D_{\log}$ in PSPACE and a PR map $W$ such that every logarithmic-space descriptor $(e,c)$ disagrees with $D_{\log}$ on $W(e,c,0)$. Thus $\mathsf L\subsetneq\mathsf{PSPACE}$.
\end{theorem}
\begin{proof}
On a valid $W(e,c,0)$ of length $n$, put $s=c\ell(n)$, where $\ell(n)=\lceil\log_2(n+2)\rceil$, and output $1-B_s(e,x)$. On other words output zero. Because $c\le n$, we have $s\le n\ell(n)\le(n+2)^2$. Keep the program on the input and rescan it as necessary. Its source-state code uses at most $n$ bits. The work configuration and binary configuration-clock counter occupy $O(n+s)$ cells, so this is a polynomial-space decider.

For each supplied $(e,c)$ the parsing equation gives the exact same $s$ used by that descriptor on the displayed input. Complementation therefore gives an explicit disagreement. In particular it disagrees with every ordinary total logarithmic-space decider after the constant-factor normalization above. Inclusion follows by the inequality $\ell(n)\le n+1$. The proof uses bounded simulation and Boolean arithmetic, not any claim that an arbitrary logspace program eventually halts.
\end{proof}

\begin{theorem}[H04: PSPACE is properly contained in EXPSPACE]\label{thm:H04}
There is a PR language $D_{\mathrm{ps}}$ in EXPSPACE with a PR disagreement input for every polynomial-space descriptor. In particular $\mathsf{PSPACE}\subsetneq\mathsf{EXPSPACE}$.
\end{theorem}
\begin{proof}
Return zero on invalid self-descriptions. On $W(e,c,k)$, use the budget $s=c(n+1)^k$ and output $1-B_s(e,x)$. The H01 inequality gives $s\le2^{(n+1)^2}$. A configuration and its binary clock use $O(n+s)$ space, at most $2^{K(n+1)^2}$ for a fixed $K$. This belongs to EXPSPACE. Its time may be double exponential, but its space bound is the asserted resource.

For every $(e,c,k)$, the word $W(e,c,k)$ parses back to that descriptor, and the bit is its complement. Polynomial-space inclusion in exponential space follows from elementary clock domination, after enlarging fixed parameters. All budgets, counters, and trace bounds are PR functions even when double exponential.
\end{proof}

\begin{theorem}[H05: an exponential-time hierarchy step]\label{thm:H05}
Let $2\mathsf{EXPTIME}$ denote the union of deterministic time classes with clocks $2^{2^{p(n)}}$ for polynomials $p$. There is an explicit PR language in $2\mathsf{EXPTIME}$ outside EXPTIME, with PR counterexample inputs for all exponential-time descriptors. Hence $\mathsf{EXPTIME}\subsetneq2\mathsf{EXPTIME}$.
\end{theorem}
\begin{proof}
Return zero on invalid self-descriptions. On $W(e,c,k)$ complement $\B(e,x,c2^{(n+1)^k})$. Since $c,k\le n$,
\[
 \log_2(c2^{(n+1)^k})
 \le n+(n+1)^n\le n+2^{n^2}\le2^{(n+1)^2+1}.
\]
The inequalities can be expressed without logarithms as bounds on the clock itself. Fourth-power simulation and countdown production fit $2^{2^{K(n+1)^2}}$ after enlarging a fixed $K$. The parsing equation again gives a mismatch against every lower-class descriptor. The ordinary inclusion is the elementary inequality between the clock families. This proof generalizes, at each fixed tower height, by applying one more exponential to the same padded-descriptor bound; we count only the explicitly stated step here.
\end{proof}

\paragraph{What has not been inferred.}
These separations do not imply $\mathsf P\ne\mathsf{NP}$, $\mathsf P\ne\mathsf{PSPACE}$, or $\mathsf{PSPACE}\ne\mathsf{EXPTIME}$. They are witnesses against precisely the descriptor families stated. Nor does finite diagonalization prove that every negation of an unbounded agreement assertion has a PR counterexample function. The sharp deterministic and nondeterministic hierarchy theorems need their constructibility and simulation hypotheses; our coarser proofs avoid asserting those stronger results without developing the corresponding machinery.

\section{Alternation, the polynomial hierarchy, and parallel circuits}
\label{sec:alternation}

\subsection{Finite alternation}
An alternating machine has terminal accepting/rejecting configurations and nonterminal existential or universal configurations, each with two successors. Duplicate a successor to represent a unary transition. A finite time bound $T$ supplies the evaluation horizon. At an existential node take OR of the successor bits; at a universal node take AND. Terminal bits are retained. A nonterminal node at horizon zero has value zero. Thus every evaluation is a finite Boolean tree, including malformed or cyclic descriptions.

Although this definition mentions a recursively branching tree, its arithmetic value is PR. For a fixed horizon, bound every possible configuration code as in F05. Tabulate all its value bits at remaining time zero and update the entire table $T$ times. Alternatively traverse the binary tree of at most $2^{T+1}-1$ nodes with an explicitly bounded stack. Both implementations use only primitive recursion. The definition does not quantify over infinite strategies or arbitrary sets of configurations.

\begin{theorem}[A01: alternating polynomial time equals polynomial space]\label{thm:A01}
For clocked finite presentations, $\mathsf{AP}=\mathsf{PSPACE}$, with descriptor transformations and \ipra\ proofs of equality of acceptance bits.
\end{theorem}
\begin{proof}
For a time-$T(n)$ alternating machine, evaluate its tree by depth-first search. Each stack frame stores the current configuration, a phase bit, and the first child's value. There are at most $T(n)+1$ frames, each with $O(n+T(n)+1)$ cells for a fixed machine. Thus a polynomial $T$ gives polynomial space. The invariant says every completed frame equals the finite OR/AND recurrence; all frames, depths, configurations, and partial tree codes have displayed finite bounds. The traversal clock is a polynomial factor times $2^{T+1}$, which is PR. There is no demand that the evaluation take polynomial \emph{time}.

In the reverse direction, Theorem~\ref{thm:S12} transforms a fixed polynomial-space computation into a polynomial-length quantified Boolean formula. Its construction is polynomial time and uses the single-copy reachability recursion, so it does not unroll an exponential path. Evaluate each existential variable block by existential choices and each universal block by universal choices; save the assigned bits. At a leaf evaluate the quantifier-free matrix deterministically. Every branch uses a polynomial number of transitions in the input length. The QBF simulation theorem identifies its bit with the original space computation. Composing the two finite machine constructions gives the claimed transformations. This reconstructs the polynomial-time part of the classical alternation theorem \cite{ChandraKozenStockmeyer1981}.
\end{proof}

\subsection{Fixed levels of the polynomial hierarchy}
For a fixed $k\ge1$, a $\Sigma_k^{\mathsf P}$ presentation consists of a polynomial-time Boolean verifier $V$ and polynomial bounds $p_1,\ldots,p_k$ such that
\[
 A(x)=\bigvee_{u_1<2^{p_1(n)}}
       \bigwedge_{u_2<2^{p_2(n)}}\cdots
       Q_{u_k<2^{p_k(n)}} V(x,u_1,\ldots,u_k).
\]
The alternating folds begin with OR; $\Pi_k^{\mathsf P}$ begins with AND. These are first-order bounded formulas and PR bits. Set $\Sigma_0^{\mathsf P}=\Pi_0^{\mathsf P}=\mathsf P$. The unbounded union PH is terminology in the metatheory; every assertion below specifies what happens at each fixed finite level.

\begin{theorem}[A02: duality, monotonicity, and the space bound for PH]\label{thm:A02}
For every fixed $k$, complementation interchanges $\Sigma_k^{\mathsf P}$ and $\Pi_k^{\mathsf P}$. Each level is included in both appropriately prefixed next levels, and all finite levels are contained in PSPACE. In particular $\Sigma_1^{\mathsf P}=\mathsf{NP}$ and $\Pi_1^{\mathsf P}=\mathsf{coNP}$.
\end{theorem}
\begin{proof}
Complement the matrix bit and exchange each bounded OR with AND. Finite Boolean De Morgan identities prove the result by induction on the fixed number of folds; they do not use De Morgan laws for arbitrary unbounded predicates. A dummy one-bit block does not change a bit independent of that variable. Insert a dummy block at the beginning for the opposite leading quantifier, or at the end for same-leading monotonicity; merge adjacent blocks of the same kind by pairing their words. This gives the usual level inclusions with explicit polynomial witness lengths.

There is an alternating polynomial-time machine making the indicated block choices and then running $V$; Theorem~\ref{thm:A01} supplies a polynomial-space evaluator. Directly, a depth-first traversal needs only the $\sum_i p_i(n)$ current assignment bits, a bounded stack of counters, and polynomial verifier storage. The first-level equalities are the verifier definitions.
\end{proof}

\begin{theorem}[A03: bounded-alternation QBF completeness]\label{thm:A03}
For each fixed $k\ge1$, truth of polynomial-size quantified Boolean formulas with $k$ alternating blocks beginning existentially is $\Sigma_k^{\mathsf P}$-complete under polynomial many-one reductions. The universally beginning version is $\Pi_k^{\mathsf P}$-complete. The quantifier-free matrix may be an arbitrary Boolean formula; no restricted CNF claim is needed.
\end{theorem}
\begin{proof}
Membership follows by using the quantifier blocks as the alternating witness blocks and evaluating the matrix in polynomial time. For hardness, fix a source verifier $V(x,\vec u)$ and its clocks. The circuit construction in N03 yields a polynomial-size circuit with the bits $\vec u$ as inputs and $x$ fixed. Let $z$ contain the gate values, let $E(\vec u,z)$ be the conjunction of their local equations, and let $o(z)$ be the output bit. Unique extension of an acyclic circuit proves, for every assignment $\vec u$, both
\[
 V(x,\vec u)=1\ \longleftrightarrow\ \exists z\,(E(\vec u,z)\land o(z)),
\]
and
\[
 V(x,\vec u)=1\ \longleftrightarrow\ \forall z\,(E(\vec u,z)\to o(z)).
\]
The second equivalence uses the actual unique extension for the reverse implication; if the output is zero, that extension is a counterexample. Every quantifier is over a finite bit array.

If the final source block is existential, use the first matrix and put $z$ in that block. If it is universal, use the second and put $z$ in that block. Earlier blocks are unchanged. The conjunction of gate equations is a polynomial-length formula, so the construction is a polynomial-time printer with exactly $k$ blocks after harmless dummy variables. Bitwise equivalence under each fixed outer fold proves the reduction in \ipra. Complementing gives the dual statement, or repeat the same last-block construction.
\end{proof}

\begin{theorem}[A04: conditional collapses, stated constructively]\label{thm:A04}
For each fixed $k$, the numerical hypothesis $\mathsf P=\mathsf{NP}$ implies $\Sigma_k^{\mathsf P}=\Pi_k^{\mathsf P}=\mathsf P$. The hypothesis $\mathsf{NP}=\mathsf{coNP}$ implies $\Sigma_k^{\mathsf P}=\Pi_k^{\mathsf P}=\mathsf{NP}$ for $k\ge1$.
\end{theorem}
\begin{proof}
Under P=NP, an existentially bounded polynomial projection of a P predicate is an NP predicate, hence has a P presentation by the hypothesis. A universal projection is the complement of the existential projection of the complemented bit, and P is closed under complement. Eliminate the finitely many blocks from the inside out. Each use of the hypothesis supplies a numerical descriptor by existential elimination; the constructed final descriptor is the required witness. If a particular SAT decider is supplied, the Cook--Levin reduction provides an explicit compiler parameterized by that decider. A bare existential class-equality hypothesis is not silently replaced by a PR global choice function.

Under NP=coNP, complementing an NP predicate again has an NP presentation. Existentially projecting an NP predicate concatenates the projected word and its NP witness, so it remains NP. A universal projection is complement--existential--complement, hence remains NP. The same finite inside-out argument applies. In both proofs induction on the number of blocks is in the metatheory over the fixed formula construction; the internal inductions are only the bounded array and simulation inductions already justified. Neither hypothesis is asserted to hold.
\end{proof}

\subsection{Uniform Boolean circuits}
An acyclic Boolean circuit is a topologically ordered array of gates; its inputs precede its operations. A family is \emph{logarithmic-space uniform} if a fixed halting logspace transducer, given $1^n$, prints the circuit for inputs of length $n$. A fixed polynomial bounds the number of gates and the length of all gate records. $\mathsf{NC}^k$ uses fan-in at most two and depth $O(\ell(n)^k)$; $\mathsf{AC}^k$ allows polynomially bounded fan-in for AND and OR, with NOT of fan-in one. Constants, initial input gates, and a designated output are included. We use these uniform definitions throughout this subsection. Arbitrary nonuniform advice is discussed separately in the counting section.

\begin{theorem}[A05: balancing and sequential simulation]\label{thm:A05}
For every fixed $k\ge0$,
\[
 \mathsf{NC}^k\subseteq\mathsf{AC}^k\subseteq\mathsf{NC}^{k+1},
 \qquad \mathsf{NC}=\mathsf{AC}\subseteq\mathsf P,
\]
where NC and AC denote their finite-level unions with the uniformity convention above.
\end{theorem}
\begin{proof}
The first inclusion merely permits more inputs to a gate. Replace an AND or OR gate with $r\le n^a+a$ inputs by a balanced binary tree. Pad to a power of two with the identity bit (one for AND, zero for OR). Induction on tree height proves equal output. It has at most $2r$ gates and height $\lceil\log_2\max(2,r)\rceil=O(\ell(n))$. Replacing all gates therefore increases size by at most a polynomial and depth by at most this logarithmic factor. The printer enumerates each source gate and the binary-tree addresses; these counters need $O(\log(n+2))$ bits. To recover a source gate it reruns the source uniform printer with an output-position counter. Thus uniformity is preserved.

Finally a uniform polynomial-size circuit is evaluated sequentially in topological order, storing one bit per gate and looking up its predecessors by bounded scans. The uniform printer is polynomial time: its internal configuration count is polynomial and, being a halting deterministic transducer whose output is never read, it cannot revisit an internal configuration. Both printing and evaluation are polynomial time. A bounded prefix invariant says every stored gate value equals its Boolean expression. The union equalities follow from the two adjacent-level containments.
\end{proof}

\begin{theorem}[A06: shallow uniform circuits use logarithmic space]\label{thm:A06}
$\mathsf{NC}^1\subseteq\mathsf L$. More generally, a uniform bounded-fan-in circuit of depth $d(n)$ and polynomial size can be evaluated using $O(d(n)+\ell(n))$ work space, with a PR time bound when $d$ has one.
\end{theorem}
\begin{proof}
Unfold the output gate into a binary expression tree, duplicating shared subcircuits only conceptually. A depth-first traversal stores the left/right path from its root, the phase at each level, and a single saved first-child value for each active binary gate. These are $O(d)$ bits in total. Do \emph{not} store a logarithmic-size gate address in every frame. Instead recover the current gate by replaying the root-to-current path, keeping just one gate number. Given that number, rerun the uniform circuit printer to obtain its gate record. This uses $O(\ell(n))$ additional space; replay counters use $O(\log(d+2))\le O(d+1)$ bits.

At each return the saved bit equals the value of the already evaluated subtree. The path and phase invariant proves that the algorithm eventually returns the output value after traversing at most $2^{d+1}-1$ occurrences, with a polynomial overhead for gate recovery. For $d=O(\ell(n))$ the space is logarithmic. The general time is deliberately allowed to exceed polynomial when $d$ is larger. All traversal bounds are finite PR expressions, so the same correctness proof belongs to \ipra.
\end{proof}

\begin{theorem}[A07: NL has uniform quadratic-logarithmic-depth circuits]\label{thm:A07}
$\mathsf{NL}\subseteq\mathsf{NC}^2$. Consequently
\[
 \mathsf{NC}^1\subseteq\mathsf L\subseteq\mathsf{NL}
 \subseteq\mathsf{NC}^2\subseteq\mathsf P.
\]
\end{theorem}
\begin{proof}
First take an explicitly encoded graph with $N$ vertices. Let $R_0=I\lor E$ and define Boolean matrix squaring
\[
 R_{j+1}(u,v)=\bigvee_{z<N}(R_j(u,z)\land R_j(z,v)).
\]
Induction shows that $R_j(u,v)$ means existence of a walk of length at most $2^j$. After $\lceil\log_2\max(1,N)\rceil$ stages this is reachability, by simple-path shortening. Each stage uses $O(N^3)$ AND/OR gates and depth $O(\log(N+2))$ when the ORs are balanced. Total depth is $O(\log^2(N+2))$, size $O(N^3\log(N+2))$, and stage/vertex/tree-address counters produce it in logarithmic space.

For a fixed logspace machine, its configuration graph has $N=n^{O(1)}$ vertices. A configuration has $O(\ell(n))$ bits. Adjacency of two such configurations is tested by the fixed transition table and equality of all unmodified work-tape positions. Selecting the input symbol at a coded input-head position is a balanced multiplexer on the $n$ input bits, of depth $O(\ell(n))$ and polynomial size. All configuration fields and work symbols likewise use polynomially many equality and bit-selection gates with depth $O(\ell(n))$. Thus the initial graph matrix has a uniform polynomial-size circuit of logarithmic depth. Append the matrix-squaring circuit. If accepting configurations are multiple, OR over them. The total remains uniform NC$^2$ and agrees with the nondeterministic acceptance bit by the configuration graph theorem.
\end{proof}

The result is an upper-bound construction, not a proof that NC equals P or that its displayed containments are strict. The exact gate and table meanings used here are finite and numerical; no circuit family is quantified over as an arbitrary set inside \ipra.

\section{Finite randomness, counting, circuits, and lower bounds}
\label{sec:random-lower}

The central objects in this section are finite tables and finite computation
traces. A probability is a rational abbreviation for a ratio of two natural
numbers; every inequality below has an explicitly stated integer form.
The standard complexity results provide the targets
\cite{adleman1978,lautemann1983,shannon1949,katz2011}.
The proofs here reconstruct them with bounded searches, positive correctness
statements, and numerical induction invariants. In particular, the assertion
that a finite bad set is smaller than its ambient set produces its least
missing element. It is not an appeal to an unbounded choice principle.

\subsection{The counting primitives and the randomized presentations}

For a coded Boolean table or a fixed primitive-recursive Boolean predicate
$E(i,\vec a)$, put
\[
 C_E(N,\vec a)=\sum_{i<N}[E(i,\vec a)=1].
\]
The brackets have values zero and one. The defining recursion is
$C_E(0)=0$ and $C_E(N+1)=C_E(N)+[E(N)=1]$.
Sets in this section are abbreviations for such predicates on bounded
intervals; their sizes mean these sums. Tuples are encoded in a fixed radix,
so a tuple of $k$ elements below $Q$ is a number below $Q^k$.
The intended language is always specified by a Boolean function $L(x)$,
or by a fixed bounded definition already compiled to such a function.

\begin{lemma}[R01: finite counting calculus]\label{thm:R01}
The following statements have primitive-recursive witnesses and proofs by
quantifier-free induction in the bounded compilation described earlier.
\begin{enumerate}
\item $0\leq C_E(N)\leq N$, and $C_E(N)>0$ if and only if there is an
      $i<N$ with $E(i)=1$. If the count is positive, the least such $i$ is
      obtained by a bounded scan.
\item If $E(i)\leq F(i)$ for all $i<N$, then $C_E(N)\leq C_F(N)$.
      For a finite family $(E_j)_{j<k}$,
      $|\bigcup_{j<k}E_j|\leq\sum_{j<k}|E_j|$.
\item $|E\times F|=|E||F|$. More generally, a sum of products over a
      finite Cartesian product factors into the product of the sums.
\item If a permutation of $[0,N)$ carries $E$ onto $F$, their counts agree.
      If an explicitly described involution without fixed points exchanges
      $E$ and its complement, then $2|E|=N$.
\item If $|E|<N$, the least $i<N$ outside $E$ exists and is found by bounded
      search. The same assertion applies to a union whose displayed upper
      bound is strictly below $N$.
\end{enumerate}
\end{lemma}
\begin{proof}
For the first two assertions, add the next zero-or-one term to both sides
of the maintained inequality. For the union bound use the pointwise
inequality $[\bigvee_{j<k}E_j(i)]\leq\sum_{j<k}[E_j(i)]$, then sum in $i$.
For the product assertion, split the radix enumeration into consecutive
blocks of the second radix and apply distributivity; induction in the
number of blocks proves the equality. A permutation rearranges a finite
sum, as is proved by tracking the counts of its already visited indices.
An exchanging involution is a bijection between the two parts, so their
counts add to $N$ and are equal. Finally the scan for a first missing
element has invariant: either an already recorded index is missing, or
every earlier index belongs to $E$. If the scan ended with no missing
index, its count would be $N$, contradicting the strict bound. All these
invariants contain only bounded quantifiers over explicitly coded tables.
Their Boolean compilations, rather than induction on arbitrary formulas,
are the object-theory induction instances.
\end{proof}

A randomized polynomial-time presentation consists of a fixed deterministic
machine $M$, an explicit polynomial clock, and a polynomial $r(n)$ bounding
the number of independent random bits. Unused bits are ignored, so each
input of length $n$ has exactly
\[
 Q(n)=2^{r(n)}
\]
equally weighted seeds. Its clocked output is denoted $M(x,u)\in\{0,1\}$,
for $u<Q(n)$. Define
\[
 a_M(x)=\sum_{u<Q(n)}M(x,u),\qquad
 e_{M,L}(x)=\sum_{u<Q(n)}[M(x,u)\ne L(x)].
\]
All bounds are on binary input length; enumeration of the $Q(n)$ seeds is
used in the proof theory, not asserted to be polynomial time.

An RP presentation satisfies
\[
 L(x)=0\Longrightarrow a_M(x)=0,\qquad
 L(x)=1\Longrightarrow 2a_M(x)\geq Q(n).
\]
A coRP presentation has $a_M(x)=Q(n)$ on yes inputs and
$2a_M(x)\leq Q(n)$ on no inputs. A BPP presentation satisfies
$3e_{M,L}(x)\leq Q(n)$ for every $x$. A PP presentation satisfies the
positive equivalence
\[
 L(x)=1\quad\Longleftrightarrow\quad 2a_M(x)>Q(n).
\]
The left side is a decidable Boolean equality. These conventions prevent
an undecidable proposition from being inserted as an unspecified
``intended language'' behind two one-sided implications.

\begin{theorem}[R02: one-sided randomness and witnesses]\label{thm:R02}
$\mathrm{P}\subseteq\mathrm{RP}\subseteq\mathrm{NP}$ and
$\mathrm{P}\subseteq\mathrm{coRP}\subseteq\mathrm{coNP}$.
An RP computation repeated independently $k\geq1$ times, accepting if any
trial accepts, still has no false positives and has at most $Q^k/2^k$
false-negative seeds. The analogous AND construction amplifies coRP.
In particular RP and coRP are contained in BPP.
\end{theorem}
\begin{proof}
A deterministic decider may ignore its random seed. On an RP yes input,
$2a_M\geq Q>0$ gives $a_M>0$; Lemma~\ref{thm:R01} supplies an accepting
seed, of polynomial length, as an NP witness. Conversely an accepting
seed is impossible when $L(x)=0$. Complementing both the output and the
language gives the coRP assertion.

For amplification let $b=Q-a_M$ on a yes input, so $2b\leq Q$. Exactly
$b^k$ seed tuples make all $k$ trials reject, by the product identity.
Thus $2^k b^k\leq Q^k$. On a no input each trial rejects, so the OR still
rejects on every tuple. Two trials give error at most $1/4$, which is at
most $1/3$. AND amplification follows by exchanging zero and one.
The machine uses $kr(n)$ random bits and $k$ times the original clock,
plus polynomial overhead for counters and the accumulated Boolean result.
\end{proof}

\begin{theorem}[R03: integer BPP amplification]\label{thm:R03}
For every $d\geq1$, taking the majority of $12d$ independent BPP trials
(breaking ties in a fixed way) gives
\[
 2^d e_{\mathrm{amp},L}(x)<Q(n)^{12d}.
\]
The number of random bits is $12d\,r(n)$; if $d$ is polynomial in $n$,
the amplified machine is polynomial time.
\end{theorem}
\begin{proof}
Fix $x$, let $b=e_{M,L}(x)$, and put $k=2t$. For a seed tuple $\vec u$,
let $j(\vec u)$ be its number of erroneous trials. Factorization of finite
sums gives the exact identity
\[
 \sum_{\vec u\in[0,Q)^k}2^{j(\vec u)}=(Q+b)^k.
\]
Indeed each coordinate contributes weight two on its $b$ erroneous seeds
and weight one on its $Q-b$ correct seeds, whose sum is $Q+b$.
An erroneous majority has $j(\vec u)\geq t$, regardless of the tie rule.
If $E$ is its error count, then
\[
 2^tE\leq(Q+b)^{2t},\qquad
 3^{2t}2^tE\leq(4Q)^{2t},
\]
because $3b\leq Q$. Cancelling the positive factor $2^t$ yields
$9^tE\leq8^tQ^{2t}$. Now choose $t=6d$. The explicit integer check
$2\cdot8^6=524288<531441=9^6$, followed by induction on $d$, gives
$2^d8^{6d}<9^{6d}$. Combining the inequalities proves the result.
No logarithm, real exponential, Chernoff bound, or integration is used.
\end{proof}

\begin{theorem}[R04: elementary closure of randomized classes]\label{thm:R04}
BPP is closed under complement, union, intersection, and polynomial-time
many-one preimages. RP and coRP are each closed under union, intersection,
and polynomial-time many-one preimages.
\end{theorem}
\begin{proof}
For complement, replace $M$ by $1-M$: its error indicator relative to
$1-L$ is identical. For BPP union and intersection, first amplify both
machines with $d=3$. Each has error less than $1/8$. Run them with separate
seeds and take respectively OR or AND. A wrong combined answer requires
at least one wrong component answer, hence the union bound gives total
error less than $1/4<1/3$ after clearing the seed denominators.

For RP union, run both original machines and take OR. There are no false
positives, and on a yes input at least one component succeeds on at least
half its seeds. For RP intersection first amplify both to failure at most
$1/4$, and take AND. On a yes input, at most half the seed pairs fail by
the union bound; on a no input one component always rejects. Complement
these constructions to obtain the two coRP closures. Finally compute a
reduction $f(x)$ deterministically and run the chosen randomized machine
on $f(x)$. Polynomial output length composes with its clock and random-bit
bound. Its count inequality is exactly the old inequality for $f(x)$.
\end{proof}

\subsection{Counting functions, majority computation, and advice}

A $\#\mathrm P$ presentation is a polynomial-time Boolean verifier
$V(x,w)$ and a polynomial witness length $p(n)$, with
\[
 f(x)=\sum_{w<2^{p(n)}}V(x,w).
\]
This is an actual primitive-recursive natural-valued function. Witnesses
are bit strings of an exact declared length. A convention that merely
allows arbitrary trailing padding would change counts and is therefore
not silently interchangeable with this one.

\begin{theorem}[R05: algebra of witness counts]\label{thm:R05}
$\#\mathrm P$ is closed under addition, multiplication, and polynomially
bounded sums of uniformly polynomial-time verifiers. Its positive supports
are precisely NP. Every nonnegative polynomial-time computable function
of polynomial output length is a $\#\mathrm P$ function. A fixed
presentation $f(x)=\#\{w<2^{p(n)}:V(x,w)=1\}$ also gives a
$\#\mathrm P$ presentation of $2^{p(n)}-f(x)$.
\end{theorem}
\begin{proof}
For addition, pad both witness lengths to their maximum by requiring the
new padding bits to be zero. Prepend a tag bit selecting the first or
second verifier. There is exactly one padded witness for each original
witness, so the tag partitions the accepted witnesses into two disjoint
sets of sizes $f(x)$ and $g(x)$. For multiplication concatenate one witness
for each verifier and accept exactly when both verifiers accept. The
product identity in Lemma~\ref{thm:R01} gives count $f(x)g(x)$.

For a sum over $i<k(n)$, use an index of length
$\lceil\log_2(k(n)+1)\rceil$ and a common polynomial witness length;
reject invalid indices and require unused witness padding to be zero.
The accepted pairs partition by $i$. A positive witness count is
equivalent, with a bounded least-witness map, to existence of an accepted
witness, which is exactly the NP presentation. If $f(x)<2^{p(n)}$ is
computed in polynomial time, accept $w$ precisely when $w<f(x)$; its
count is $f(x)$. Finally complement the Boolean verifier to obtain the
count $2^{p(n)}-f(x)$. This last result is relative to a specified witness
cube, and does not assert closure under arbitrary subtraction.
\end{proof}

\begin{theorem}[R06: PP containments and complement]\label{thm:R06}
$\mathrm{NP}\cup\mathrm{coNP}\cup\mathrm{BPP}\subseteq\mathrm{PP}
\subseteq\mathrm{PSPACE}$, and PP is closed under complement. Every
$\#\mathrm P$ function is computable using polynomial work space and
exponential time.
\end{theorem}
\begin{proof}
For an NP verifier with $Q=2^{p(n)}$ witnesses and acceptance count $a$,
use a tag bit and $p(n)$ other bits. On tag zero accept; on tag one run
the verifier. Of the $2Q$ seeds, exactly $Q+a$ accept, which is more than
half precisely when $a>0$. Thus the constructed PP language is exactly
the NP support, with both implications witnessed by finite arithmetic.

For complement closure, first pad the random length of a PP machine so
it is at least one. Write $Q=2H$. Use $2Q$ seeds $(b,u)$ with $u<Q$.
For $b=0$ accept if the original machine rejects. For $b=1$ accept if
$u<H+1$. The new acceptance count is
\[
 Q-a+H+1=3H+1-a.
\]
It is strictly greater than half of $2Q$, namely $Q=2H$, precisely when
$a\leq H$. This is the complement of $a>H$. The construction fits the
seed space since $H\geq1$ and $H+1\leq Q$. It also gives
$\mathrm{coNP}\subseteq\mathrm{PP}$ by the preceding NP construction.
A BPP yes input has $3a\geq2Q$, whence $2a>Q$; a no input has
$3a\leq Q$, whence $2a\leq Q$. Its original machine therefore already
satisfies the PP specification.

To compute a count, enumerate all $p(n)$-bit witnesses, retaining the
current witness, a $p(n)+1$-bit accumulator, and the work tape of the
polynomial-time verifier. Reuse the work tape between witnesses. The
workspace is polynomial and the time is $2^{p(n)}$ times a polynomial.
Comparing the resulting count to the half-cube threshold proves the PP
containment. The invariant after witness $w$ is that the accumulator
equals the sum on earlier witnesses; its Boolean form is bounded.
\end{proof}

\begin{theorem}[R07: BPP has polynomial advice and circuits]\label{thm:R07}
From a BPP presentation one can construct, for every length $n$, a
polynomial-length advice string $a_n$ that makes a fixed polynomial-time
machine correct on all $n$-bit inputs. The function $n\mapsto a_n$ can be
chosen primitive recursively. Consequently
$\mathrm{BPP}\subseteq\mathrm P/\mathrm{poly}$ in the finite circuit
and advice readings.
\end{theorem}
\begin{proof}
Amplify with $d=n+2$. Let $m(n)=12(n+2)r(n)$ and let $Q'=2^{m(n)}$ be
the amplified seed count. For each $n$-bit input $x$, fewer than
$Q'/2^{n+2}$ seeds give a wrong answer. The union over the $2^n$ inputs
has fewer than $Q'/4$ members, after multiplying through by $2^{n+2}$.
It is therefore not the whole seed cube. Choose the least seed outside
this union and call it $a_n$. The search predicate is a conjunction over
$2^n$ inputs of finite machine computations and comparisons with $L(x)$.
Thus both the search bound $Q'$ and its Boolean predicate are primitive
recursive. No countable choice of unrelated advice strings is necessary.

The deterministic advice machine runs the amplified machine with seed
$a_n$. Its clock is polynomial and its output equals $L(x)$ for every
input of length $n$, by the definition of the selected seed. Unrolling
this fixed clocked computation into a Boolean circuit, with the advice
bits hardwired, produces a polynomial-size circuit. Conversely a circuit
description of polynomial length can be used as advice by a universal
circuit evaluator; the gates are evaluated in topological order in
polynomial time. This explains the equivalence of the two formulations
used here. It does not say that the advice can be found in polynomial
time. Nor does it say that arbitrary languages in $\mathrm P/\mathrm{poly}$
are recursive: the primitive-recursive advice selector is available here
because it is constructed from the given BPP presentation.
\end{proof}

\subsection{A fully finite proof that BPP lies at the second hierarchy level}

\begin{theorem}[R08: the translate-covering construction]\label{thm:R08}
Every BPP presentation has both a $\Sigma^\mathrm P_2$ presentation and a
$\Pi^\mathrm P_2$ presentation, with explicit polynomial bounds and
primitive-recursive witness searches. Hence
\[
 \mathrm{BPP}\subseteq\Sigma^\mathrm P_2\cap\Pi^\mathrm P_2.
\]
\end{theorem}
\begin{proof}
Let $r=r(n)$ be the original random-bit bound and set
\[
 d=16(r+1)^2,\qquad m=12dr,\qquad Q=2^m,\qquad k=m+1.
\]
Amplify as in Theorem~\ref{thm:R03}; call the resulting predicate $A(x,u)$.
On a yes input its rejecting set has size $b$ with $2^d b<Q$; on a no
input its accepting set has size $a$ with $2^d a<Q$.
We require $k<2^d$. To verify it without an asymptotic exception, put
$v=r+1\geq1$. Then $k\leq193v^3<2^{8+3v}\leq2^{16v^2}$, using
$v\leq2^v$ and $8+3v\leq16v^2$. All constants and inequalities are
integer ones.

Bitwise exclusive-or on $m$-bit words is denoted $\oplus$; translation
by a fixed word is a permutation. Consider
\[
 \exists s_0,\ldots,s_{k-1}<Q\quad
 \forall u<Q\quad \bigvee_{i<k} A(x,u\oplus s_i)=1.                 \tag{R8a}
\]
Its witness block has length $km$, its universal block has length $m$,
and its matrix is decidable by $k$ polynomial-time runs. It is therefore
a $\Sigma^\mathrm P_2$ formula.

On a yes input, fix $u$. Exactly $b^k$ of the $Q^k$ shift tuples have all
$u\oplus s_i$ in the rejecting set: each coordinate is a permutation of
the seed cube. At most $Qb^k$ shift tuples leave some $u$ uncovered, by
the union bound. Since $2b<Q$ and $k=m+1$,
\[
 Qb^k < Q\,Q^k/2^k = Q^k/2 < Q^k.
\]
Thus a least covering tuple exists. This proves (R8a) constructively.
On a no input, any $k$ translates of the accepting set have union size
at most $ka<Q$, since $k<2^d$ and $2^d a<Q$. The least uncovered $u$
refutes that tuple. This supplies the reverse implication with a bounded
counterexample function.

For completeness, the dual positive presentation is
\[
 \forall s_0,\ldots,s_{k-1}<Q\quad
 \exists u<Q\quad \bigwedge_{i<k} A(x,u\oplus s_i)=1.              \tag{R8b}
\]
On a yes input the $k$ translates of the small rejecting set cannot cover
the cube, so their least uncovered point supplies $u$. On a no input the
large rejecting set has a covering tuple by the already proved argument,
and that tuple refutes (R8b). This proves the $\Pi^\mathrm P_2$ reading
directly, without using unrestricted classical contraposition. The special
case $m=0$ is harmless: $Q=k=1$, and the strict error bounds force zero
errors. Every search in the argument is bounded by $Q$ or $Q^k$.
\end{proof}

\subsection{Zero-error computation through finite truncations}

A bounded zero-error trial returns one of $0,1,?$ in polynomial time.
It has \emph{trace soundness} if every definite answer equals $L(x)$,
and \emph{success density at least one half} if at least half its seeds
produce a definite answer. These are finite Boolean and counting tests.
For a possibly longer-running random machine $Z$, normalize it to inspect
at most one fresh random bit per step. For a prefix of $T$ bits, let
$c_T(x,u)$ be the smaller of its halting time and $T$, defined by a
$T$-step simulation. The finite expected-time certificate with bound $p$
is the arithmetical statement
\[
 \forall x\ \forall T\qquad
 \sum_{u<2^T}c_T(x,u)\leq p(|x|)2^T.                              \tag{R9}
\]
The certificate is accompanied by trace soundness at every truncation.
This is the precise zero-error expected-time reading used here.

\begin{theorem}[R09: zero-error equivalence and restart bounds]\label{thm:R09}
The following presentations are interconvertible by explicit machine
constructions: (i) a bounded zero-error trial of success density at least
one half; (ii) an RP presentation and a coRP presentation of the same
Boolean language; (iii) a trace-sound random machine with a polynomial
finite expected-time certificate (R9). Thus the finite arithmetical
reading of $\mathrm{ZPP}=\mathrm{RP}\cap\mathrm{coRP}$ is provable in
\ipra. Repeating (i) for at most $j$ rounds leaves failure density at most
$2^{-j}$, and its average number of rounds is at most two for every cap.
\end{theorem}
\begin{proof}
From (i), accept only on output one to get RP. Accept on output one or
question mark to get coRP. Trace soundness eliminates the forbidden
one-sided errors, and the definite-output density gives the required
one-half success inequalities.

Conversely let $R$ be an RP test for $L$ and let $S$ be an RP test for
$1-L$, obtained by complementing a coRP test. Run them using disjoint
seeds. Return one if $R$ accepts; otherwise return zero if $S$ accepts;
otherwise return question mark. On a yes input, $S$ never accepts and
$R$ accepts on at least half its seeds. On a no input the roles reverse.
Both a definite wrong answer and a conflict of two valid answers are
therefore excluded by trace soundness. This constructs (i).

Pad one trial to a fixed polynomial duration $p(n)\geq1$, so that
independent blocks of $p(n)$ fresh bits suffice. If $Q=2^{p(n)}$ and
$b\leq Q/2$ is the number of inconclusive blocks, exactly $b^j$ block
tuples fail all $j$ trials. Hence $2^j b^j\leq Q^j$. If $N_j$ is the
number of rounds used by the procedure capped at $j$, finite double
counting gives
\[
 \sum_{\vec u\in[0,Q)^j} N_j(\vec u)
 =\sum_{i=0}^{j-1} b^i Q^{j-i}
 \leq Q^j\sum_{i=0}^{j-1}2^{-i}\leq2Q^j.
\]
The displayed fractions abbreviate the equivalent inequality after
multiplication by $2^{j-1}$ when $j>0$; the case $j=0$ is zero.
The equality counts a round whenever all preceding rounds were
inconclusive. Running successive independent trials until one succeeds
therefore gives (R9) with bound $2p(n)$ at every finite step cap: charge
at most $p(n)$ steps for each entered round, and apply the same finite
tail sum, including the possibly incomplete last block. More explicitly, let
$j=\lceil T/p(n)\rceil$ and extend each $T$-bit prefix to $jp(n)$ bits.
For every extension, $c_T\leq p(n)N_j$. Summing over all extensions and
using their equal multiplicity $2^{jp(n)-T}$ gives
$2^{jp(n)-T}\sum_{u<2^T}c_T(x,u)\leq2p(n)2^{jp(n)}$; cancelling
the positive multiplicity gives (R9). The case $T=0$ is immediate. The trace remains
sound because any eventual answer came from a sound trial. This proves
(i)$\Rightarrow$(iii) without asserting termination on every infinite
random tape.

Finally assume (iii), increase $p$ to at least one if necessary, and cap
the machine at $T=2p(n)$ steps. If $B$ prefixes have not halted by that
cap, each contributes exactly $T$ to the sum in (R9). Thus
$TB\leq p(n)2^T$, and cancellation of $p(n)>0$ yields
$2B\leq2^T$. Return question mark on these prefixes and the observed
answer otherwise. This is (i). The argument is the finite counting form
of Markov's inequality; no measure or assertion of universal random-tape
termination enters it.
\end{proof}


\subsection{Circuit counting with an explicit hard family}

For this subsection, a circuit on $n$ input bits is a finite straight-line
program. It has input wires, two constant wires, and topologically ordered
gates. Each gate applies one of the sixteen binary Boolean functions to two
earlier wires. The output is a specified wire. This finite basis contains
AND, OR, NOT, and copying; using it instead of a smaller complete basis
changes the relevant size bounds only by fixed factors. A truth table is
an integer $h<2^{2^n}$, whose $x$th bit is the value on input $x<2^n$.

\begin{theorem}[R10: quantitative circuit counting]\label{thm:R10}
The number of $n$-input Boolean functions computable with at most $s$ gates
is at most
\[
 B(n,s)=(n+s+2)\bigl(16(n+s+2)^2\bigr)^s.
\]
If $B(n,s)<2^{2^n}$, a primitive-recursive bounded search produces a truth
table not computed by any such circuit.
\end{theorem}
\begin{proof}
Pad a shorter circuit to exactly $s$ gates with unused copy gates, retaining
its output wire. For an upper bound, allow every gate to choose its two
sources from all $n+s+2$ wire positions, even the positions that would be
invalid because they are not earlier. There are at most
$16(n+s+2)^2$ descriptions for each gate and at most $n+s+2$ choices of
output. Every valid padded circuit occurs in this bounded enumeration,
and different descriptions may compute the same function. Consequently
the number of computed truth tables is bounded by $B(n,s)$.

Enumerate a circuit's truth table by evaluating its gates in order on
each of the $2^n$ inputs. Invalid descriptions are marked invalid; they
are not used as extra circuits. Scan the $2^{2^n}$ truth tables and choose
the first one distinct from the tables of all valid descriptions. If no
such table existed, marking the computed tables in an initially zero
bit vector would mark all $2^{2^n}$ positions with at most $B(n,s)$ marks.
The invariant that the number of marked positions is at most the number
of processed descriptions contradicts the strict inequality. The scan,
evaluation, marks, and equality comparisons are all bounded and primitive
recursive. They provide the witness, not merely a double-negated existence.
\end{proof}

\begin{theorem}[R11: a concrete Shannon bound]\label{thm:R11}
For $n\geq4$ let
\[
 s(n)=\left\lfloor\frac{2^n}{16(n+1)}\right\rfloor.
\]
At most $2^{2^{n-1}}$ truth tables on $n$ bits have circuits of size at most
$s(n)$. In particular a least hard table $h_n$ exists, and at least
$2^{2^n}-2^{2^{n-1}}$ tables require more gates.
\end{theorem}
\begin{proof}
Put $s=s(n)$. Since $s\leq2^{n-4}$ and
$n+2+2^{n-4}\leq2^n$ for $n\geq4$, we have $n+s+2\leq2^n$.
The latter elementary inequality starts at $n=4$ and is preserved when
$n$ increases, as is verified after doubling the exponential term.
The description bound is therefore
\[
 B(n,s)\leq2^{n+(4+2n)s}.
\]
By definition $16(n+1)s\leq2^n$, and $4+2n\leq3(n+1)$ for $n\geq1$;
hence $(4+2n)s\leq3\cdot2^{n-4}$. Also
$n\leq2^{n-2}=4\cdot2^{n-4}$ for $n\geq4$, by a one-step induction.
Thus
\[
 n+(4+2n)s\leq7\cdot2^{n-4}\leq2^{n-1}.
\]
The previous theorem proves the stated count and constructs $h_n$.
The phrase ``almost all'' can consequently be given the exact finite
content that the fraction of tables with these small circuits is at most
$2^{-2^{n-1}}$. Only the cross-multiplied integer inequality is used in
\ipra.
\end{proof}

\begin{corollary}[R12: a primitive-recursive language outside polynomial circuits]
\label{thm:R12}
There is a primitive-recursive Boolean language $H$, decidable in
space $2^{O(n)}$, with no polynomial-size circuit family. For every
$c\geq1$ and $k\geq0$, the explicitly chosen length
\[
 n=2^{4(k+c+4)}
\]
has the property that no circuit with at most $c(n+1)^k$ gates computes
$H$ correctly on all inputs of that length. For each such circuit a least
mismatch input is primitive recursively obtainable.
\end{corollary}
\begin{proof}
Let $H(x)$ be the $x$th bit of $h_n$ on length $n\geq4$, with arbitrary
fixed values at the finitely many smaller lengths. Here $x$ is its
fixed-length binary value; the ambient word coding also records $n$.
The bounded construction of $h_n$ proves primitive recursiveness.

To verify the advertised length, set $u=k+c+4$ and $t=4u$. For $t\geq4$,
$2^t\geq t^2$, as follows by induction from $t=4$ and
$(t+1)^2\leq2t^2$. Moreover
$t^2>4+c+(k+1)(t+1)$, since the right side is at most
$4u^2+2u$ and $t^2=16u^2$. As $n=2^t$ and $n+1\leq2^{t+1}$,
\[
 16c(n+1)^{k+1}
 \leq2^{4+c+(k+1)(t+1)}<2^{2^t}=2^n.
\]
Therefore $c(n+1)^k\leq s(n)$. The defining property of $h_n$ excludes
all circuits up to that size. Scan the $2^n$ inputs to find a mismatch
with any specified candidate; the absence of a mismatch would contradict
the definition of the hard table.

For the space bound, store one candidate table of $2^n$ bits, one circuit
description of $O(s(n)\log(n+s(n)+2))=O(2^n)$ bits, the current circuit's
gate values, and binary loop counters. A brute-force comparison may
recompute every candidate circuit's value on each input, so no collection
of all truth tables is stored. The counters for the truth-table and
description enumerations themselves have $O(2^n)$ bits. The work space
is therefore $2^{O(n)}$; the time can be doubly exponential. This proves
$\mathrm{EXPSPACE}\not\subseteq\mathrm P/\mathrm{poly}$ by an explicit
recursive example. It asserts no reverse containment: unrestricted
nonuniform advice can encode nonrecursive languages.
\end{proof}

\begin{theorem}[R13: a constructive universal circuit upper bound]
\label{thm:R13}
Every truth table on $n\geq1$ inputs has a circuit with at most
$n2^n+n$ gates, constructed primitive recursively from the table.
For $n\geq64$ there is a compiler producing a circuit $C$ with
\[
 n\,\operatorname{size}(C)\leq18\cdot2^n.
\]
Together with Theorem~\ref{thm:R11}, this gives the order
$\Theta(2^n/n)$ for the maximum, over truth tables of length $2^n$, of
minimum circuit size. No real asymptotic expression is needed for its
integer formulation.
\end{theorem}
\begin{proof}
First precompute the $n$ negated inputs. For every input assignment with
table bit one, form the conjunction of its $n$ matching literals, then
OR these terms. At most $2^n$ terms are present, and the zero-term case
uses the constant zero wire. The bound $n2^n+n$ follows by charging
$n-1$ gates to each conjunction and at most one OR gate per term.
The correctness invariant says that each partial OR equals the union
of the already processed accepted assignments.

For the sharper bound put
$m=\lfloor\log_2 n\rfloor-1$ and $q=n-m$. Then
$n/4<2^m\leq n/2$. Split an input into a $q$-bit prefix $a$ and an
$m$-bit suffix $b$. Precompute, as separate circuits, all the
$2^{2^m}$ Boolean functions of $b$ by the preceding construction. Each
requires at most $m2^m+m\leq n^2$ gates, so this library costs at most
$n^2 2^{\lfloor n/2\rfloor}$ gates. For every prefix $a$, the input
truth table explicitly specifies one residual function $b\mapsto H(a,b)$,
whose output wire is now available.

Build all prefix indicators using a shared binary prefix tree. First
precompute the negations of the prefix input bits, at cost at most $n$.
From the constant-one empty indicator, form the indicators for both
extensions of every prefix by AND with the next literal. This costs
less than $2^{q+1}$ gates. For each full prefix, AND its indicator with
its residual-function wire, and OR the resulting $2^q$ terms. These last
operations cost at most $2^{q+1}$ gates. Exactly one indicator is true
on every input, so the output equals the specified residual function.
The total gate count is at most
\[
 n^2 2^{\lfloor n/2\rfloor}+n+4\cdot2^q.
\]
For $n\geq64$, $n^3\leq2^{\lfloor n/2\rfloor}$: check $64$ and $65$,
then use $(n+2)^3\leq2n^3$ for $n\geq8$ to advance by two.
Consequently the first term, multiplied by $n$, is at most $2^n$.
Also $n^2\leq2^n$. Finally $2^m>n/4$ gives
$n\cdot4\cdot2^q<16\cdot2^n$. Adding the three integer bounds gives
the claimed constant 18. All residual-table lookups and circuit wiring
are bounded operations, so this is a uniform primitive-recursive
compiler from a truth table to its circuit, although its input and output
can themselves be exponentially large in $n$.
\end{proof}

\subsection{Decision-tree and communication lower bounds}

A query tree is a coded finite rooted binary tree. Its internal vertices
specify which input bit to inspect, and its leaves carry output bits.
A comparison tree instead specifies comparisons between two input keys,
and its leaves give permutations of the keys. Repeated queries are
allowed; they can only increase the measured number of queries.

\begin{lemma}[R14: finite tree capacity]\label{thm:R14}
A binary tree of depth at most $h$ has at most $2^h$ leaves. An arbitrary
finite binary tree has at most $2^h-1$ leaves at depth strictly below $h$.
\end{lemma}
\begin{proof}
The first assertion follows by induction on $h$: a leaf contributes one,
and otherwise the two subtrees each contribute at most $2^{h-1}$.
Equivalently, map each leaf path to its extensions to length $h$; the
extension sets are nonempty and disjoint. For the second assertion,
encode a path of length $j<h$ by the integer $2^j$ plus its binary value.
There are only $1+2+\cdots+2^{h-1}=2^h-1$ such codes. Distinct leaves
give distinct path codes. These arguments can be implemented as bounded
counts over the tree code, so no induction on arbitrary set-valued trees
is required in the object language.
\end{proof}

\begin{theorem}[R15: comparison sorting requires many comparisons]
\label{thm:R15}
Any comparison tree that correctly sorts $n$ distinct keys has worst-case
depth $h$ satisfying
\[
 2^h\geq n!.
\]
If $n\geq4$, $a=\lfloor n/2\rfloor$, and
$\ell=\lfloor\log_2 a\rfloor$, then $h\geq a\ell$.
Under the uniform distribution on the $n!$ rank orders, the average
comparison count is also $\Omega(n\log n)$ in an explicit integer form:
for every $j$ with $2^j\leq n!/2$,
\[
 2\sum_{\pi}\operatorname{depth}(\pi)\geq j\,n!.
\]
\end{theorem}
\begin{proof}
Represent the $n!$ orders by permutations of ranks $0,\ldots,n-1$.
Each rank order has a unique permutation of positions that lists the
keys increasingly. Two distinct rank orders cannot reach the same leaf
of a correct tree, because that leaf specifies only one permutation.
The map from rank orders to leaves is therefore injective. Enumeration
of permutations, for example by factorial-radix codes, and a finite
occupancy count prove that there are at least $n!$ leaves. Apply
Lemma~\ref{thm:R14}.

The largest $a$ factors in $n!$ are each at least $a$, so
$n!\geq a^a\geq2^{a\ell}$. Strict monotonicity of powers of two gives
$h\geq a\ell$. This avoids Stirling's formula and real logarithms in
the proof. For the average bound, fewer than $2^j$ leaves can have depth
below $j$. Hence at least $n!/2$ of the $n!$ rank orders have depth at
least $j$, and summing their depths gives the displayed inequality.
Choosing $j=\lfloor\log_2(n!)\rfloor-1$ when this is nonnegative gives
the stated order bound. These conclusions are specifically about the
comparison model; integer-key algorithms using other operations are not
covered by the lower bound \cite{mit2010sorting}.
\end{proof}

\begin{theorem}[R16: exact deterministic query bounds]\label{thm:R16}
The deterministic query complexity of $n$-bit OR, AND, and parity is $n$.
For parity, every correctly answered input path must query all $n$ bits.
For OR the all-zero path must query all bits, and for AND the all-one
path must do so.
\end{theorem}
\begin{proof}
Suppose the all-zero OR path omits position $i$. Flip only that position.
Induction through the path shows every queried answer remains the same,
so the same leaf is reached, although the correct OR value changes.
The omitted position can be chosen as the least position absent from
the recorded query list. The AND argument starts from the all-one input.
For parity, flip any omitted position on any input; the transcript again
remains unchanged and parity changes. These are explicit error witnesses.
Reading all positions and folding OR, AND, or XOR gives matching upper
bounds. The empty-input case has zero queries with its stipulated
constant output.
\end{proof}

\begin{theorem}[R17: finite randomized query lower bounds]\label{thm:R17}
Suppose a randomized query algorithm uses a fixed finite uniform seed
space and at most $q$ queries on every input and seed, and errs on at
most one third of seeds for every input. For $n$-bit OR, $3q\geq n$.
For $n$-bit parity, $q\geq n$. The same parity conclusion applies to any
error bound strictly below one half.
\end{theorem}
\begin{proof}
Let the seed count be $R>0$. For OR, let $S_i$ be the set of seeds for
which the run on the all-zero input queries position $i$. Couple that
run to the run with only bit $i$ set, using the same seed. If the zero
run does not query $i$, induction on the transcript shows the two runs
are identical. They cannot both be correct because their OR values
differ. The two error sets each have size at most $R/3$, so
$R-|S_i|\leq2R/3$, or $3|S_i|\geq R$. On the other hand, each seed
occurs in at most $q$ of the sets $S_i$, whence
$\sum_i|S_i|\leq qR$. Summing the first inequalities gives
$nR\leq3qR$; cancel $R>0$.

For parity assume $q<n$. Fix a seed, obtaining a deterministic tree.
For each reachable leaf, choose the least input position not queried on
its path. Flipping that bit preserves the leaf and exchanges the two
parity values. It is an involution without fixed points on the inputs
reaching that leaf. Thus exactly half of those inputs are answered
correctly by the leaf's fixed output. Summing over leaves, each seed is
correct on exactly half of the $2^n$ inputs. Summing over the $R$ seeds
therefore gives exactly $R2^n/2$ correct input--seed pairs. But correctness
on at least two thirds of the seeds at every input gives at least
$2R2^n/3$ such pairs, a contradiction after multiplying by six. More
generally any strict majority guarantee contradicts the same exact
half-count. The argument uses finite double counting and an explicitly
computable pairing, not a minimax theorem.
\end{proof}

A deterministic two-party communication protocol is a finite binary tree.
Each internal vertex specifies the speaker and a Boolean function of
that speaker's input and the preceding transcript. A leaf gives the
output to both participants. Communication cost is the maximum transcript
length. Fixing this shared leaf-output convention matters for exact
additive constants.

\begin{lemma}[R18: rectangles and fooling sets]\label{thm:R18}
For each protocol leaf the input pairs reaching it form a Cartesian
product $X_0\times Y_0$. If a Boolean communication problem has a finite
set $F$ of one-inputs such that two distinct $(x,y),(x',y')\in F$ have
at least one zero cross-input $(x,y')$ or $(x',y)$, then every correct
protocol has at least $|F|$ one-leaves, and $2^h\geq|F|$ for its cost $h$.
\end{lemma}
\begin{proof}
At the root every input pair is possible. A message from Alice restricts
only the first coordinate of the current Cartesian product, and a
message from Bob restricts only the second. Induction on the transcript
proves the rectangle assertion. If two distinct elements of $F$ reached
one leaf, both cross-inputs would reach it as well. At least one has
value zero, contradicting the one label required by the two elements of
$F$. Thus $F$ injects into the one-leaves. Use the finite tree bound.
All sets can be represented by bounded Boolean matrices and row/column
bit vectors \cite{blais2025communication}.
\end{proof}

\begin{theorem}[R19: deterministic equality needs $n+1$ bits]
\label{thm:R19}
For $n\geq1$, equality of two $n$-bit strings has deterministic
communication complexity exactly $n+1$ under the shared leaf-output
convention.
\end{theorem}
\begin{proof}
The $2^n$ diagonal pairs $(x,x)$ form a fooling set: for $x\ne x'$, the
cross-pairs have unequal coordinates. Therefore at least $2^n$ leaves
must be labelled one. Since $n\geq1$, a zero-input exists and must reach
some zero-leaf. The tree has at least $2^n+1$ leaves, so $2^h>2^n$ and
$h\geq n+1$. For the upper bound Alice sends her $n$ input bits; Bob
compares with his input and sends the resulting answer bit. If only Bob
is required to know the final answer, this last bit can be omitted and
the corresponding exact bound becomes $n$. The $2^n$ diagonal pairs
still require distinct transcripts: a shared transcript for $(x,x)$ and
$(y,y)$ with $x\ne y$ would give Bob the same view on $(x,y)$ as on
$(y,y)$ and force a wrong answer. Thus the lower bound $n$ persists.
\end{proof}

\begin{theorem}[R20: a constant-communication public-randomness equality test]
\label{thm:R20}
With $kn$ shared random bits, equality has a protocol communicating
$k+1$ bits, having no false negatives and false-positive density exactly
$2^{-k}$ on unequal inputs. In particular three communicated bits suffice
for error at most one third, independently of $n$.
\end{theorem}
\begin{proof}
View the shared bits as $k$ independent masks $u_j\in\{0,1\}^n$.
Alice sends the $k$ parities $\langle x,u_j\rangle$ over the two-element
field. Bob compares each with $\langle y,u_j\rangle$ and sends one if
all match, zero otherwise. Equal strings always pass. For unequal strings
let $i$ be their least differing coordinate. Flipping bit $i$ of a mask
exchanges agreement and disagreement of the two inner products. Exactly
half the masks collide, by the pairing lemma. Independence, expressed
by the product count, gives exactly $2^{kn-k}$ colliding mask tuples.
Taking $k=2$ gives error $1/4$ and cost three. The bound uses public
randomness explicitly; it is not a claim about private-coin protocols.
\end{proof}

\subsection{Finite costs and amortized analysis}

\begin{theorem}[R21: best, weighted average, and worst cost]
\label{thm:R21}
Let $c_0,\ldots,c_{N-1}$ be natural costs with $N>0$, and let
$w_0,\ldots,w_{N-1}$ be natural weights of total $W>0$. If
$b=\min_i c_i$ and $z=\max_i c_i$, then
\[
 bW\leq\sum_{i<N}w_ic_i\leq zW.
\]
Both extrema and the weighted sum are primitive recursive on the coded
finite lists. These inequalities give best $\leq$ weighted average
$\leq$ worst after rational interpretation.
\end{theorem}
\begin{proof}
The bounded scan computing the minimum and maximum retains the least
index achieving the running extremum. For every $i$, multiplication of
$b\leq c_i\leq z$ by the nonnegative integer $w_i$ preserves order.
Add these inequalities and use $\sum_i w_i=W$. The proof is induction
on the list length with a numerical partial-sum invariant.
\end{proof}

There is no automatic ordering of an amortized bound between an
input-distribution average and a worst-case single-operation cost.
Amortized analysis controls totals along valid sequences of state changes;
it is not a third statistic of the same finite input distribution. The
following statement is the exact useful principle.

\begin{theorem}[R22: the potential method over natural numbers]
\label{thm:R22}
For a finite state sequence $s_0,\ldots,s_m$, let $c_i$ be the cost of
the transition $s_i\to s_{i+1}$, and let $\Phi(s)$ be a nonnegative
integer potential. If
\[
 c_i+\Phi(s_{i+1})\leq a_i+\Phi(s_i)\qquad(i<m),
\]
then
\[
 \sum_{i<m}c_i+\Phi(s_m)
 \leq\sum_{i<m}a_i+\Phi(s_0).
\]
In particular $\sum_{i<m}c_i\leq\sum_{i<m}a_i+\Phi(s_0)$.
\end{theorem}
\begin{proof}
For $m=0$ the inequality is equality. Add the next transition inequality
to the induction hypothesis and cancel the common intermediate potential
$\Phi(s_m)$. This gives the inequality for $m+1$. The final consequence
uses $\Phi(s_m)\geq0$. No subtraction in the natural numbers, and no
assumption that each amortized charge is a realized single-operation
cost, is needed.
\end{proof}

\begin{corollary}[R23: binary counters have constant amortized bit-flip cost]
\label{thm:R23}
Starting at zero, $m$ increments of an unbounded binary counter use at
most $2m$ bit flips. More precisely, if $\nu(m)$ is the number of one bits
in the final counter, the total number of flips is $2m-\nu(m)$.
\end{corollary}
\begin{proof}
If the current value has $t$ trailing ones, the next increment changes
those $t$ ones to zero and the following zero to one. Its cost is $t+1$,
and the potential $\Phi=\nu$ changes from $\nu$ to $\nu-t+1$.
Thus $c_i+\Phi(s_{i+1})=2+\Phi(s_i)$ exactly. Sum this identity and use
$\Phi(0)=0$ to obtain total cost plus $\nu(m)$ equal to $2m$.
The trailing-one count, population count, and binary increment are all
bounded digit scans. This is a bit-flip statement: a particular tape
implementation may require additional head-motion cost, which must be
accounted for separately.
\end{proof}


\subsection{Lean targets and the remaining proof-lowering work}

The file \texttt{lean/CountingOutline.lean} defines the semantic interfaces
and target propositions for this section. It is an \emph{uncompiled outline},
with no axioms or admitted theorem bodies, rather than a completed proof of
these results. Its two elementary recursion equations are the only theorem
bodies. The substantive targets are definitions of propositions whose proofs
remain to be implemented.

The first layer defines \texttt{sumBelow}, \texttt{countBelow}, least-witness
specifications, product and union count targets, and finite random-machine
families. The proof order is to establish the two recursion equations,
induct on loop bounds for addition and monotonicity, prove radix product
factorization, and then implement first-witness search. The weighted identity
in R03 uses the same product lemma with values one and two in place of bits.
R07 and R08 then require only the finite union bound, XOR permutations,
and the machine-to-circuit compiler already developed in the NP section.
The zero-error interface explicitly records truncated cost and answer;
its operational instantiation must prove consistency of prefixes,
absorbing halting states, and the fact that an unfinished $T$-step trace
has cost exactly $T$. These conditions cannot be replaced by unrelated
semantic cost and answer functions.

The second layer defines gates, topologically valid circuits, a circuit
interpreter, truth tables, the description bound, and hard-table and
universal-upper-bound targets. The implementations follow the proof:
bounded description enumeration, a marking-count invariant, and least
missing-table search; then literal conjunctions, prefix indicators, and
residual-function selection. The final layer supplies finite query trees,
comparison trees, and communication protocols, together with their
interpreters and depth functions. The principal lemmas are preservation of
transcripts under an unqueried bit flip and preservation of Cartesian
products along protocol paths. Counts of leaves, permutations, and mask
pairs finish their lower bounds.

All these Lean function and inductive-type parameters are a semantic
presentation of the mathematics. For an actual \ipra{} derivation, instantiate
them with the fixed primitive-recursive programs, encode every finite object
as a natural number, and replace every bounded property by its Boolean
compiler. Prove the loop invariants with the repository's quantifier-free
induction rule and check the resulting object-language derivation. A kernel
proof of an ambient Lean target alone would not establish that last fact.
The thirteen independent checks in \texttt{code/random\_lower\_checks.py}
exercise finite instances and exact identities; they supply implementation
evidence and do not replace universal proofs.

\section{A finite arithmetic treatment of a quantum gate model}
\label{sec:quantum}

The main article directly links quantum machines, BQP, and QMA. A full treatment of arbitrary real amplitudes would require an additional representation and approximation theory. We can nevertheless prove substantial quantum-complexity statements by finite signed counts. This section makes the gate model explicit rather than hiding that issue.

Use $q$ qubits, initially a computational-basis word containing the classical input and zero ancillas. Allowed gates are NOT $X$, classical reversible Toffoli, and Hadamard $H$. Toffoli maps $(a,b,c)$ to $(a,b,c\mathbin{\mathrm{xor}}ab)$. NOT flips a bit. On a basis bit $a$, Hadamard creates the two possible outputs $b\in\bits$ with sign $(-1)^{ab}$ and factor $1/\sqrt2$. Distinct control and target wires are required. A constant-one ancilla, made with $X$, lets Toffoli implement controlled NOT. The circuit contains a polynomial number of qubits and gates, is printed by a polynomial-time uniform procedure, and one designated output qubit is measured at the end. This defines the explicitly named class $\mathsf{BQP}_{HT}$ when every input has acceptance probability at least $2/3$ or at most $1/3$ as specified by its language bit.

The subscript is deliberate. Hadamard plus classical reversible gates are an important universal model in the standard theory \cite{Shi2003}, but this report does not reprove the approximation/universality theorem identifying every other conventional presentation of BQP with this one. Our theorems concern exactly these finite gate descriptions and their positive gap promises. They do not permit a real amplitude to serve as an uncomputable advice parameter \cite{AdlemanDeMarraisHuang1997}.

\begin{theorem}[Q01: exact path-sum identity]\label{thm:Q01}
Let a circuit contain $h$ Hadamard gates. For each $u\in\bits^h$, simulate the classical reversible gates and use $u_i$ as the chosen output at the $i$th Hadamard. Let $y(u)$ be the final basis word and let $\sigma(u)\in\{1,-1\}$ be the product of its Hadamard signs. Define
\[
 s_y=\sum_{u<2^h}[y(u)=y]\,\sigma(u),\qquad
 A=\sum_{y:\,y_{\rm out}=1}s_y^2.
\]
Then the exact acceptance probability is $A/2^h$, and $\sum_y s_y^2=2^h$. These identities, including $0\le A\le2^h$, have finite \ipra\ proofs without introducing square roots as arithmetic objects.
\end{theorem}
\begin{proof}
The usual amplitude notation $s_y/(\sqrt2)^h$ is only explanatory; the arithmetic state is the array of signed integers $s_y$ together with the exponent $h$. Initially one array entry is one and the others are zero. A reversible classical gate permutes entries. At a Hadamard on one wire, the old entries indexed by the two values of that wire, say $a,b$, become $a+b,a-b$, while $h$ increases by one. The identity
\[
 (a+b)^2+(a-b)^2=2(a^2+b^2)
\]
shows that the sum of squares doubles. Thus it equals $2^h$ by induction on gate position. The signs $(-1)^{ab}$ give exactly the same two-entry update when the branch sums are expanded, proving the path formula by the same bounded induction.

There are $2^q$ entries, each with absolute value at most $2^h$, and at most $2^h$ paths. Signed-array and path-code bounds are explicit exponentials, so every invariant is bounded. Final acceptance sums precisely the squared entries whose output bit is one; the normalization identity gives the bounds. All calculations use signed integer addition and multiplication.
\end{proof}

\begin{theorem}[Q02: classical random computation embeds exactly]\label{thm:Q02}
$\mathsf{BPP}\subseteq\mathsf{BQP}_{HT}$, and in particular $\mathsf P\subseteq\mathsf{BQP}_{HT}$. The reduction preserves each input's exact rational acceptance probability.
\end{theorem}
\begin{proof}
Take the polynomial-size circuit obtained by unrolling a clocked randomized verifier, with its $r$ random bits exposed as inputs. Apply Hadamard once to each of $r$ initially zero ancillas, giving one branch of weight $2^{-r}$ for each seed. Evaluate the classical Boolean circuit reversibly by storing every intermediate gate output on a fresh zero ancilla. AND is a Toffoli into a zero target. Repeated inputs are simplified first: AND$(a,a)$ and OR$(a,a)$ are copies, implemented by a controlled NOT with a distinct reserved constant-one control. All Toffoli targets and controls are distinct wires. To compute NOT into a fresh target, set that target to one and XOR the source into it. OR follows from $a\lor b=\neg(\neg a\land\neg b)$ with fresh ancillary outputs. Thus each gate has a constant-size reversible implementation, and the whole transformation has polynomial size and a polynomial-time printer.

Keep the original seed register throughout; there is no need to erase the recorded history. Different seeds remain different orthogonal basis words. Hence no cross terms mix their output probabilities. Measuring the output bit accepts on exactly the seeds for which the classical verifier accepts, with probability equal to their count divided by $2^r$. The bit-level reversible identities and the injection from seeds into final basis words prove the arithmetic statement. A deterministic decider is the case of no random seeds.
\end{proof}

\begin{theorem}[Q03: a PP simulation by an explicit branch gap]\label{thm:Q03}
$\mathsf{BQP}_{HT}\subseteq\mathsf{PP}$, in the exact finite-majority definition of PP used in the counting section.
\end{theorem}
\begin{proof}
Append two Hadamard gates on a fresh unused zero wire. They compose to the identity: the two-entry update twice is $(a,b)\mapsto(2a,2b)$, and the amplitude denominator gains a factor two (the probability denominator gains a factor four). Probabilities do not change, and now $h\ge2$.

For two branch words $u,v\in\bits^h$, put
\[
 g(u,v)=[y(u)=y(v),\ y(u)_{\rm out}=1]\,\sigma(u)\sigma(v)
 \in\{-1,0,1\}.
\]
Expanding the squares in Q01 gives $\sum_{u,v}g(u,v)=A$. This contribution is computable in polynomial time from $u,v$ by two gate-list simulations; it needs no amplitude array.

Construct a fair-coin machine using a selector bit, the $2h$ bits of $(u,v)$, and one extra bit. In selector mode zero, accept both extra-bit branches if $g=1$, reject both if $g=-1$, and accept exactly one if $g=0$. Its number of accepting branches minus rejecting branches is $2A$.

In selector mode one there are $2^{2h+1}$ assignments to the other bits. Interpret them as an integer $j$ and accept exactly when
\[
 j<2^{2h}-2^{h-1}.
\]
Comparison of these $O(h)$-bit numbers takes polynomial time. This mode has acceptance-minus-rejection gap $-2^h$. The entire machine therefore has gap
\[
 2A-2^h.
\]
It accepts with strict majority exactly when $A/2^h>1/2$. The promised $2/3$ versus $1/3$ gap makes this precisely the supplied language bit. The number of random bits is $2h+2$, polynomial in the original input length. The proof is finite counting and signed arithmetic, hence available in \ipra. This gives the path-counting content of the standard containment while retaining our stated gate restriction.
\end{proof}

\begin{theorem}[Q04: exact simulation in polynomial space]\label{thm:Q04}
A polynomial-size $HT$ circuit's integer numerator $A$ and the comparison $2A>2^h$ can be computed in polynomial space. Consequently $\mathsf{BQP}_{HT}\subseteq\mathsf{PSPACE}$ independently of the PP containment.
\end{theorem}
\begin{proof}
Enumerate the $2^q$ possible final basis words $y$ with a $q$-bit counter. For each such word, enumerate all $2^h$ Hadamard choice words, simulate the $L$ gates using a $q$-bit classical configuration and a sign bit, and add its sign exactly when the final word equals $y$. The partial sum needs at most $h+2$ signed bits. Square the result and, for an accepting $y$, add it to an accumulator. Even before using normalization, the crude upper bound $2^q 2^{2h}$ needs only $q+2h+1$ bits. Store the polynomial-size gate list and the current input and counters. All storage is polynomial in $q+h+L$, hence in input length.

The nested counters give an explicit exponential time bound, for example $2^{q+h}$ times a fixed polynomial in $q+h+L$. Prefix-sum invariants prove that the final accumulator equals Q01's $A$. This algorithm's use of exponential time is compatible with polynomial space. As before, the proof never postulates an infinite random tape or a completed real Hilbert space.
\end{proof}

\section{Lean formalization: definitions, exact targets, and remaining work}
\label{sec:formalization}

\subsection{The endpoint is an object-language derivation}
The repository's inspected toolchain is \texttt{leanprover/lean4:v4.32.0}. Its \texttt{ConciseIPRA.PR} module defines the arity-indexed syntax
\begin{lstlisting}
inductive PR : Nat -> Type
  | zero : PR 0
  | succ : PR 1
  | proj {k} (i : Fin k) : PR k
  | comp {m k} (f : PR m) (gs : Fin m -> PR k) : PR k
  | prec {k} (f : PR k) (g : PR (k + 2)) : PR (k + 1)
\end{lstlisting}
It then defines arithmetic terms \texttt{PTm n}, formulas \texttt{Fm PTm n}, arithmetic axioms \texttt{PRAAx}, and intuitionistic derivations. For a closed numerical sentence $\varphi$, the final goal has type
\begin{lstlisting}
([] : List (Fm PTm 0)) |- [PRAAx] phi
\end{lstlisting}
where the ASCII display represents the repository's Unicode turnstile notation. The actual syntax is shown in \texttt{lean/integration/ObjectTargets.lean}. A host-level theorem about an arbitrary function \texttt{Nat -> Bool} does not have that type. Neither a semantic structure with a correctness field nor an unproved target proposition is an object-language proof.

The supplied files are deliberately separated into definitions, target statements, and completion notes. Their comments identify which predicates remain to be lowered into PR syntax and which resource-machine simulations are still to be proved. There is no claim that the files have been typechecked in this task. This meets the requested outline stage, and leaves a concrete rather than rhetorical completion path.

\subsection{A complete numerical convention for the interpreter}
The independent \texttt{CoreOutline.lean} and \texttt{machine\_checks.py} use this canonical source model. The input is read only. Its symbols are $0$ for the left endmarker, $1$ for data zero, $2$ for data one, and $3$ for the right endmarker. The one-sided work tape has base-three digits, with $0$ blank, $1$ data zero, and $2$ data one. States $0,1$ are absorbing rejection/acceptance; the initial state is $2$. Initially both heads are at position zero, so the input head is on the left endmarker, and the work tape is blank. Left moves clamp at zero; input moves also clamp at the right endmarker. The transition writes before moving.

A source table has $q\ge3$ states and $12q$ rules. A rule with next state $a<q$, write symbol $b<3$, input move $u<3$, and work move $v<3$ is encoded as
\[
 27a+9b+3u+v<27q,
\]
where move values $0,1,2$ mean left, stay, right. The program payload is $q-3$ unary ones followed by zero, then $12q$ fixed-width rule fields of width $w=\lceil\log_2(27q)\rceil$, with least-significant bit first within each field. Its exact length is $(q-2)+12qw$. The word marker retains zero suffixes. Invalid length, fields, headers, or input codes receive the rejecting interpretation. The rule index is $12a+3\,\text{inputSymbol}+\text{workSymbol}$. The table is explicit, so parsing and a rule lookup have the polynomial scan implementations used in F06.

For a space budget $s$, the initial configuration or any resulting work head outside $\{0,\ldots,s-1\}$ causes rejection. With $s\ge1$, the bound $(q+2)(n+2)(s+1)3^s$ suffices. All definitions are finite numerical arithmetic. The Lean record presentation is a convenient semantic view; for object-language work pack its fields into a mixed-radix record using the explicit horizon or budget bounds. The per-transition equation between that packed record and the semantic record is the first required representation theorem.

\subsection{Module responsibilities}
\begin{longtable}{p{0.25\textwidth}p{0.67\textwidth}}
\toprule
File or target & Definitions and work to finish \\
\midrule\endhead
\texttt{CoreOutline.lean} & Marker words, bounded folds, canonical rules/configurations, fixed-table decoding, time and budgeted iteration, polynomial clocks and self-descriptions. Prove codec equations, invalid defaults, trace uniqueness, and physical simulation bounds. \\
{\footnotesize\path{ReductionsOutline.lean}} & Explicit variable universes, literals, clauses, acyclic gates, Tseitin constraints, subset-sum digits, graph/SAT predicates and assignment constructors. Prove gate truth tables, prefix uniqueness, witness bijections, and the declared reduction/solver obligations. \\
\texttt{SpaceOutline.lean} & Numeric graph tables, bounded paths, reachable rows, midpoint matrices, QBF syntax/evaluation, the single-copy selector, and finite games. Prove path/counter invariants, freshness, resource-stack implementation, and descriptor compilers. \\
\texttt{CountingOutline.lean} & Finite counts, random-program presentations, amplification inequalities, majority predicates, advice/covering targets, circuit descriptions, costs and lower-bound obligations. Prove exact inequalities before translating the quantifiers. \\
\texttt{QuantumOutline.lean} & Signed Hadamard paths, numerator counts, and the PP branch construction. Prove normalization, reversible embedding, and polynomial branch length. \\
{\footnotesize\path{integration/}\newline\path{ObjectTargets.lean}} & Actual \texttt{PR}/\texttt{PTm}/\texttt{Fm}/\texttt{Prf} target definitions, plus direct reuse points in ProveIt. These are specified tasks, not invented completed derivations. \\
\bottomrule
\end{longtable}

The high-level statements for A01--A07 reduce to these modules' machine, circuit, and QBF primitives. Add explicit printers for quantified block expansion, balanced gate replacement, configuration adjacency, and matrix squaring; prove their output lengths and logspace counters separately. A semantic \texttt{List.map} is not a proof of a logarithmic-space printer.

\subsection{How to lower a proof, step by step}
\paragraph{1. Construct the function syntax.}
For each numerical operation, build a finite \texttt{PR k} term by zero, successor, projection, composition and recursion. Reuse the repository's existing arithmetic, division, sum, pairing, and bounded-search terms where their exact types match. Define machine execution by primitive recursion on its explicit step counter. A polynomial evaluator with a numerical exponent is PR; an arbitrary total evaluator of PR syntax is not introduced. Decode fixed finite circuit/QBF/table grammars by bounded loops on their numerical codes. Functions such as Savitch's recurrence can be defined by iterating the entire finite relation table, even though a different machine implementation is required for the small-space bound.

\paragraph{2. Fix formulas and reflect their finite tests.}
Write the actual \texttt{BFm} syntax of each bounded invariant. Examples are
\begin{align*}
 &\forall i\le j\ \operatorname{trace}(z,i)=\Run(e,x,i),\\
 &\forall i<g\ \operatorname{gateConstraint}(C,i,a)=1,\\
 &\forall u,v<N\ R_j(u,v)=[\text{a coded path of length at most }2^j],\\
 &\forall u<N\ [\text{successful counted enumeration marks }u\iff R_i(u)].
\end{align*}
The bracketed path assertion quantifies over codes below $\max(2,N)^{2^K+1}$ for a fixed external horizon parameter $K\ge j$; the transcript assertion similarly has an explicit loop-count/field bound. No quantifier is left tacitly unbounded inside the induction formula. The repository's \texttt{BoundedFormulas.bf\_reflect}, \texttt{bf\_choice}, and \texttt{bf\_least}, and the scoped \texttt{BoundedSyntaxInduction.induction}, provide exactly the required proof rule. Their input is fixed syntax, and their outputs are derivations, not truth assumptions \cite{ProveItBounded}.

\paragraph{3. Prove the local algebra and preservation.}
Gate constraints use at most eight Boolean cases. Base-eight subset-sum columns use integer digit bounds below eight and a prefix division invariant. Inductive counting uses the equality between the reported count and the sum of the reachable bitmap, not an unverified claim that a guessed absent vertex is absent. Quantum gates use the signed identity $(a+b)^2+(a-b)^2=2(a^2+b^2)$. Random amplification uses finite distributivity and explicit integer inequalities, not analytic probability. Each yields defining-equation proofs, finite case splits, or previously proved arithmetic lemmas in the object theory.

\paragraph{4. Apply only the permitted induction.}
Reflect the invariant to a Boolean term $b(j)=1$. Prove its base and successor cases with the local hypotheses, apply \texttt{PRA.qfInd}, and reflect back. A proof by structural recursion on a fixed piece of syntax is a metatheoretic constructor of finite derivations. It must not be confused with induction on an arbitrary unbounded formula inside iPRA. Likewise, statements for each fixed PH level are a family of finite proof constructions, not an unjustified internal induction on semantic class membership.

\paragraph{5. Verify a real machine for the resource bound.}
Construct an explicit transducer or stack machine implementing the semantic function. Record every tape and counter. Prove simulation of one instruction and a PR bound on its overhead, then sum these bounds. For Savitch, a whole-matrix evaluator proves definability but not its quadratic-space theorem; the separate depth-first stack bound is necessary. For nonreachability, a materialized Python bitmap is an independent checker, not the claimed O(log $N$)-workspace machine. For circuit counting, an exponential truth table is compatible with EXPSPACE but does not supply a polynomial-time family. These distinctions are proof obligations, not engineering details that can be omitted.

\paragraph{6. Assemble the closed endpoint and audit it.}
First prove the compiler theorem with numerical parameters free; instantiate clock and size bounds; close all variables by universal introduction; export the actual \texttt{Prf PRAAx} term. Run Lean with the pinned toolchain and inspect the generated theorem type and axiom dependencies. Check both that there are no admitted proof constants and that the conclusion is the intended object-language formula. A theorem in Lean's metatheory can legitimately use metatheoretic tools to construct an object proof; the decisive evidence is the resulting derivation type. None of these final kernel checks is reported as completed for the new main theorems in this archive.

\subsection{Representative exact obligations}
\paragraph{Diagonal mismatch.}
The object sentence universally quantifies numerical $e,c,k$, guards valid word code and $c>0$, and asserts
\[
 D(W(e,c,k))\ne B(e,W(e,c,k),c(|W(e,c,k)|+1)^k).
\]
It does not merely assert that no polynomial machine agrees everywhere. Prove parser inversion, substitute into the defining equation for $D$, and use Boolean complement inequivalence. Then separately attach the concrete exponential machine and cost derivation. The existing diagonal-clock inequalities can be imported without treating them as proofs of parser or simulator correctness.

\paragraph{Cook--Levin with counted witnesses.}
For the compiled circuit $C(v,x)$ and CNF $F(v,x)$, give PR assignment maps $\operatorname{extend}$ and $\operatorname{restrict}$. Prove, for all input witness words below their explicit bound, that extension satisfies $F$ exactly when the source verifier accepts; prove every satisfying target assignment is that unique extension of its restriction. This establishes decision equivalence and the stronger bijection needed by N18. The explicit input-variable universe prevents unused witness bits from being lost.

\paragraph{Nonreachability and QBF.}
Express an online nondeterministic run as a finite transcript, with each guessed path checked directly. Soundness is a bounded assertion about every transcript under its fuel bound. Completeness constructs a transcript by bounded least-path search; the constructor may be expensive, while the verifier's storage remains small. For QBF, prove every newly allocated variable is fresh and that the single recursive occurrence has the intended selector meaning. A mere semantic recurrence without a polynomial-size syntax printer would not prove PSPACE completeness.

\paragraph{Probability and advice.}
For each word $x$ and parameter $d$, the amplified-error statement is an integer inequality between exact counts. To obtain advice of length $a(n)$, prove that the number of bad lists is strictly below $2^{a(n)}$, use bounded least search, and prove the selected list succeeds on all $2^n$ inputs. This gives a PR advice selector and a nonuniform polynomial circuit, with no claim of polynomial-time advice construction. For a zero-error program, the expected-time assertion remains a family of finite truncation inequalities at every numerical cap.

\subsection{Recommended completion order}
Complete the common numerical codecs and trace theorem first. Next complete circuit evaluation/Tseitin and finite graph reachability, which have small local invariants and extensive independent checks. Then implement the physical simulation and printer interfaces. At that point Cook--Levin, graph/integer completeness, bounded-time completeness, and diagonal mismatch can be assembled. Finish the Savitch stack and inductive-counting verifier next, then QBF/alternation/parallel circuits. The exact-counting kernel supports PP, amplification, advice, lower bounds, and the quantum gap proof. This order minimizes duplicated encodings and avoids proving later semantic claims against incompatible representations.

The archive's source manifest contains the repository revision and inspected blob identities. The intended integration directory is a new packet under \path{Logic/PrimitiveRecursiveArithmetic/docs/proofs/foundations/}, with eventual checked modules in \path{Derivations/Lean/IPRADerivations/}. No repository write or commit was performed for this delivery.

\section{Coverage, source corrections, and boundaries}
\label{sec:coverage}

\subsection{Selected immediate-link coverage}
The accompanying \texttt{sources/coverage\_map.json} records 59 selected immediate outgoing topic links. Each was observed on the main article and opened through that actual parent link. It records the target URL, theorem identifiers, and whether the treatment is complete for the selected result, conditional, limited, or restricted to an explicit model. The following compact table summarizes the mathematical groups; the machine-readable ledger retains the full link provenance.

\begin{longtable}{p{0.31\textwidth}p{0.17\textwidth}p{0.44\textwidth}}
\toprule
Article topic & Results & What is established here \\
\midrule\endhead
Turing machines; determinism; multi-tape machines & F05--F06 & Total clocked simulation, trace uniqueness, and specified polynomial simulations. \\
Complexity classes; P; NP; coNP & C01--C04, N02 & First-order presentations, core containments, and closure constructions. \\
Time and space measures; input encoding & C05--C06 & Complete bounded-computation problems with unary versus binary budgets. \\
Polynomial and logspace reductions & N01, N03, N06 & Composition, resource bounds, and explicit printers. \\
SAT; NP-hardness and completeness & N03--N05, N07--N11 & Cook--Levin and graph/integer reductions with witness maps. \\
Graph isomorphism; Hamiltonian paths; optimization & N12 & Explicit NP certificates for the stated decision versions. \\
SAT fragments; structural complexity & N13--N14, N22, S09 & 2SAT and Horn algorithms; NL and P completeness. \\
Function problems and search & C04, N15--N17 & Correct graph-versus-search qualification and conditional polynomial search. \\
Counting problems; \#P & N18, R05--R06 & Parsimonious completeness, counting algebra, supports, majority classes. \\
Vertex cover; parameterized complexity & N19--N20 & A factor-two approximation and a bounded branching algorithm. \\
L; NL; nondeterministic space & S01--S09 & Graph completeness, path shortening, Savitch, and complement closure. \\
PSPACE; NPSPACE; exponential space & S06, S10--S13 & Space equalities, compact QBF, and finite-game completeness. \\
Time and space hierarchies & H01--H05 & Explicit counterexamples and the five stated coarse separations. \\
Alternating machines; complexity classes & A01--A04 & AP = PSPACE, fixed PH levels and QBF completeness, conditional collapses. \\
AC; NC; parallel computation & A05--A07 & Balancing, sequential/gate-replay evaluation, and NL in NC$^2$. \\
RP; ZPP; BPP; probabilistic machines & R01--R09 & Exact amplification, finite zero-error formulation, PP, advice and PH containment. \\
Boolean circuits; circuit complexity & R10--R13 & Counting, a PR hard family, and matching-order worst-case circuit size. \\
Decision-tree complexity; sorting & R14--R17 & Comparison and query lower bounds, including finite randomized bounds. \\
Communication complexity & R18--R20 & Rectangles, fooling sets, exact deterministic equality and public-coin fingerprints. \\
Best/average/worst; amortized analysis & R21--R23 & Finite cost bounds, the potential identity, and exact binary-counter cost. \\
Quantum machines; BQP & Q01--Q04 & Exact path sums and BPP/PP/PSPACE containments for the specified $HT$ model. \\
\bottomrule
\end{longtable}

\subsection{Statements in the article that require care}
\paragraph{The graph of a function is not automatically an efficient search oracle.}
Checking a completed proposed output and finding an output are different computational tasks. C04 states the precise bounded-search result, and N15 proves the stronger prefix-extension property needed for SAT self-reduction.

\paragraph{Amortized cost is not another point in a universal average-cost ordering.}
For a finite distribution, minimum cost is at most weighted average cost and that is at most maximum cost. Amortization bounds total cost along admissible operation sequences and has no unconditional position between the average and worst costs of unrelated individual inputs. R21--R23 give the correct separate assertions.

\paragraph{Encoding and machine hypotheses affect every resource claim.}
Unary and binary clock fields produce different complete problems. A unit-cost RAM with unrestricted enormous integers is not covered by a polynomial simulation of finite-tape bit operations. Church--Turing and practical-feasibility theses are interpretive theses, not arithmetic theorems proved by our finite machine simulations. Quantum programs in this report use a specified finite gate basis; unrestricted amplitudes are excluded.

\paragraph{Some source-page performance descriptions are stale.}
The retrieved article still discusses a subexponential graph-isomorphism bound as the best known bound. Babai's primary work gives the stronger quasipolynomial result \cite{Babai2016}. We do not reproduce the outdated claim, and we do not count the quasipolynomial algorithm as proved here. N12 gives only the elementary NP certificate that its statement explicitly promises.

\subsection{Advanced results deliberately not claimed}
\begin{longtable}{p{0.28\textwidth}p{0.64\textwidth}}
\toprule
Result or subject & What remains beyond this packet \\
\midrule\endhead
Ladner's theorem & A complete delayed-diagonalization construction and a careful intuitionistic analysis of the numerical separation hypothesis. One cannot silently extract an unbounded mismatch witness from a negated universal agreement statement. No nonderivability claim is made. \\
Sharp time hierarchy and nondeterministic time hierarchy & The optimized universal simulation, time constructibility, domination hypotheses, and the nondeterministic diagonal machinery. H01--H02 and H05 prove their displayed coarser consequences. \\
General sharp space hierarchy & The uniform space-constructibility and separation-modulus bookkeeping at arbitrary bounds. H03--H04 prove concrete large-gap separations; S08 proves complement closure under its supplied budget hypothesis. \\
AKS primality and advanced factoring & Their number-theoretic certificates, polynomial identities, and bit-complexity analysis. No unsupported polynomial primality or factoring theorem is imported. \\
IP = PSPACE and PCP & Full arithmetization, finite-field sum-check, low-degree testing, and quantitative soundness machinery. The finite counting foundation is useful but does not supply these proofs. \\
Full QMA and general quantum universality & Quantum-witness representation, gap-preserving approximations, and universal-gate compilation. Q01--Q04 prove the exact stated finite model. \\
Continuous and analog complexity & A representation of continuous inputs, moduli, approximation costs, and a precise input/output promise. A finite discretization alone is not the unrestricted continuous theorem. \\
Blum speedup and unrestricted recursive complexity measures & An explicit partial-recursion framework, the axioms for the measure, and the corresponding priority/diagonal argument. No total PR evaluator of arbitrary recursive programs is assumed. \\
Open class equalities & No assertion of P = NP, P $\ne$ NP, L = P, P = BPP, NP = coNP, or an unproved strictness between intermediate classes. Conditional compiler theorems are kept conditional. \\
\bottomrule
\end{longtable}
P versus NP remains listed as unsolved by the official Clay Mathematics Institute source \cite{ClayPNP}. Omitting an established theorem from this finite collection does not imply that it is false, nonconstructive, or unprovable in iPRA. It means its complete proof and audit have not been supplied here.

\subsection{Reproducible artifacts and validation}
The archive includes every LaTeX section, the compiled PDF, the uncompiled Lean outline modules, exact-arithmetic Python implementations and finite checks, the 86-entry claims ledger, the 59-topic coverage ledger, repository/source metadata, build instructions, and validation logs. \texttt{code/verify\_all.py} runs the five independent mathematical check programs and records their outcomes. Each check's numerical extent is reported in its own JSON or log; the aggregate does not turn experiments into universal proofs.

The reduction checks compare satisfying assignments with independently evaluated circuits, graph predicates, and subset selections. The space checks compare graph reachability with counted nonreachability and with actual quantified syntax, including syntax-growth bounds. The probability checks use exact integer or rational arithmetic. The quantum checks compare signed path sums with an independent $a+b\sqrt2$ evaluator and explicitly enumerate the PP branch construction. The machine checks compare packed base-three tapes with a separate sparse-tape interpreter. These tests target different representations of the same finite mathematics.

The PDF is built from the included source with pdfLaTeX and checked for unresolved references, clipping/overflow, and render quality. Visual inspection complements text extraction. The verification log records that Lean compilation and complete object-language proof export were not performed; neither successful Python tests nor a well-rendered PDF changes that status.

\begin{thebibliography}{99}
\addcontentsline{toc}{section}{References}
\bibitem{WikiRoot}
Wikipedia contributors. \emph{Computational complexity theory}, and the selected immediate outgoing subject links recorded in \texttt{sources/coverage\_map.json}. Retrieved 22 November 2026. Used to delimit the requested subject matter, not as the authority for the arithmetic proofs.
\url{https://en.wikipedia.org/wiki/Computational_complexity_theory}.

\bibitem{ProveItPR}
Vladimir Reshetnikov, \emph{ProveIt}, revision \texttt{3eb195cab4cc3d552a9eca78880c6d428f3e44f9}. \path{Logic/PrimitiveRecursiveArithmetic/ConciseIPRA/Lean/ConciseIPRA/PR.lean}. Inspected source for \texttt{PR}, \texttt{PTm}, \texttt{PRAAx}, and quantifier-free induction.
\url{https://github.com/VladimirReshetnikov/ProveIt/blob/3eb195cab4cc3d552a9eca78880c6d428f3e44f9/Logic/PrimitiveRecursiveArithmetic/ConciseIPRA/Lean/ConciseIPRA/PR.lean}.

\bibitem{ProveItBounded}
\emph{ProveIt}, same revision. \path{Derivations/Lean/IPRADerivations/Foundations/BoundedFormulas.lean} and \path{BoundedSyntaxInduction.lean}, under \path{Logic/PrimitiveRecursiveArithmetic}. Inspected bounded-formula and object-induction interfaces.
\url{https://github.com/VladimirReshetnikov/ProveIt/blob/3eb195cab4cc3d552a9eca78880c6d428f3e44f9/Logic/PrimitiveRecursiveArithmetic/Derivations/Lean/IPRADerivations/Foundations/BoundedSyntaxInduction.lean}.

\bibitem{ProveItDiagonal}
\emph{ProveIt}, same revision. \texttt{Foundations/DiagonalClock.lean} in the \texttt{IPRADerivations} library. Object derivations \texttt{clock\_exp} and \texttt{cost\_exp}.
\url{https://github.com/VladimirReshetnikov/ProveIt/blob/3eb195cab4cc3d552a9eca78880c6d428f3e44f9/Logic/PrimitiveRecursiveArithmetic/Derivations/Lean/IPRADerivations/Foundations/DiagonalClock.lean}.

\bibitem{ProveItBaseline}
\emph{ProveIt}, same revision. \emph{Computational complexity in intuitionistic primitive recursive arithmetic}, research packet dated 20 November 2026. README, claims ledger, foundational and diagonal sections, and object-iPRA roadmap inspected.
\url{https://github.com/VladimirReshetnikov/ProveIt/tree/3eb195cab4cc3d552a9eca78880c6d428f3e44f9/Logic/PrimitiveRecursiveArithmetic/docs/proofs/foundations/ipra_complexity_foundations_2026-11-20}.

\bibitem{Cook1971}
Stephen A. Cook. The complexity of theorem-proving procedures. \emph{Proceedings of the Third Annual ACM Symposium on Theory of Computing}, pp.~151--158, 1971. DOI: 10.1145/800157.805047. Retrieved transcription of the primary paper:
\url{https://www.cs.cmu.edu/~15455/resources/Cook1971-complx-thm-proof.pdf}.

\bibitem{Karp1972}
Richard M. Karp. Reducibility among combinatorial problems. In R. E. Miller and J. W. Thatcher, eds., \emph{Complexity of Computer Computations}, pp.~85--103, 1972. DOI: 10.1007/978-1-4684-2001-2\_9. Retrieved reprint:
\url{https://www.cs.umd.edu/~gasarch/BLOGPAPERS/Karp.pdf}.

\bibitem{AspvallPlassTarjan1979}
Bengt Aspvall, Michael F. Plass, and Robert Endre Tarjan. A linear-time algorithm for testing the truth of certain quantified Boolean formulas. \emph{Information Processing Letters} 8(3), pp.~121--123, 1979. DOI: 10.1016/0020-0190(79)90002-4.
\url{https://mathweb.ucsd.edu/~sbuss/CourseWeb/Math268_2007WS/2SAT.pdf}.
Our simpler bounded algorithm proves polynomial time, not the optimal linear bound.

\bibitem{Savitch1970}
Walter J. Savitch. Relationships between nondeterministic and deterministic tape complexities. \emph{Journal of Computer and System Sciences} 4(2), pp.~177--192, 1970. DOI: 10.1016/S0022-0000(70)80006-X.
\url{https://www.sciencedirect.com/science/article/pii/S002200007080006X}.

\bibitem{Immerman1988}
Neil Immerman. Nondeterministic space is closed under complementation. \emph{SIAM Journal on Computing} 17(5), pp.~935--938, 1988. DOI: 10.1137/0217058.
\url{https://epubs.siam.org/doi/abs/10.1137/0217058}.

\bibitem{Szelepcsenyi1988}
R\'obert Szelepcs\'enyi. The method of forced enumeration for nondeterministic automata. \emph{Acta Informatica} 26, pp.~279--284, 1988. DOI: 10.1007/BF00299636.
\url{https://doi.org/10.1007/BF00299636}.

\bibitem{HartmanisStearns1965}
Juris Hartmanis and Richard E. Stearns. On the computational complexity of algorithms. \emph{Transactions of the American Mathematical Society} 117, pp.~285--306, 1965. DOI: 10.1090/S0002-9947-1965-0170805-7.
\url{https://www3.risc.jku.at/people/schreine/courses/compcomp/HartmanisStearns1965.pdf}.

\bibitem{HennieStearns1966}
F. C. Hennie and R. E. Stearns. Two-tape simulation of multitape Turing machines. \emph{Journal of the ACM} 13(4), pp.~533--546, 1966. DOI: 10.1145/321356.321362.
\url{https://www.khoury.northeastern.edu/~viola/classes/papers/HennieStearns66.pdf}.

\bibitem{ChandraKozenStockmeyer1981}
Ashok K. Chandra, Dexter C. Kozen, and Larry J. Stockmeyer. Alternation. \emph{Journal of the ACM} 28(1), pp.~114--133, 1981. DOI: 10.1145/322234.322243. Author's publication record:
\url{https://www.cs.cornell.edu/kozen/Papers/papers_by_year.htm}.

\bibitem{adleman1978}
Leonard M. Adleman. Two theorems on random polynomial time. \emph{19th Annual Symposium on Foundations of Computer Science}, pp.~75--83, 1978. DOI: 10.1109/SFCS.1978.37. Bibliographic attribution; full-paper retrieval failed. The finite proof is developed in this report and compared with the official notes below.
\url{https://doi.org/10.1109/SFCS.1978.37}.

\bibitem{lautemann1983}
Clemens Lautemann. BPP and the polynomial hierarchy. \emph{Information Processing Letters} 17(4), pp.~215--217, 1983. DOI: 10.1016/0020-0190(83)90044-3. Publisher abstract retrieved.
\url{https://www.sciencedirect.com/science/article/pii/0020019083900443}.

\bibitem{shannon1949}
Claude E. Shannon. The synthesis of two-terminal switching circuits. \emph{Bell System Technical Journal} 28(1), pp.~59--98, 1949. DOI: 10.1002/j.1538-7305.1949.tb03624.x. Official publication record:
\url{https://www.nokia.com/bell-labs/publications-and-media/publications/the-synthesis-of-two-terminal-switching-circuits/}.
The report specifies its own finite circuit basis and explicit constants.

\bibitem{katz2011}
Jonathan Katz. \emph{Notes on Complexity Theory}. University of Maryland, 2011. Official instructor notes, especially the circuit, counting, and randomized-computation lectures.
\url{https://www.cs.umd.edu/~jkatz/complexity/f11/all.pdf}.

\bibitem{AroraBarak}
Sanjeev Arora and Boaz Barak. \emph{Computational Complexity: A Modern Approach}. Cambridge University Press, 2009. Author-hosted draft, used for orientation; access limitations are recorded in the source manifest.
\url{https://www.cs.princeton.edu/~arora/pubs/webimage.pdf}.

\bibitem{mit2010sorting}
Constantinos Daskalakis. \emph{6.006 Introduction to Algorithms, Lecture 10: Sorting}. Massachusetts Institute of Technology, Fall 2010. Official course slides.
\url{https://courses.csail.mit.edu/6.006/fall10/lectures/lec10.pdf}.

\bibitem{blais2025communication}
Eric Blais. \emph{Models of Computation, Chapter 21: Communication Complexity}. University of Waterloo, Winter 2025. Official instructor notes.
\url{https://cs.uwaterloo.ca/~eblais/cs365/w25/communication}.

\bibitem{Shi2003}
Yaoyun Shi. Both Toffoli and controlled-NOT need little help to do universal quantum computation. \emph{Quantum Information and Computation} 3(1), pp.~84--92, 2003. DOI: 10.26421/QIC3.1-7. Primary preprint:
\url{https://arxiv.org/abs/quant-ph/0205115}.
The universality theorem is cited as context, not assumed as an iPRA theorem in this packet.

\bibitem{AdlemanDeMarraisHuang1997}
Leonard M. Adleman, Jonathan DeMarrais, and Ming-Deh A. Huang. Quantum computability. \emph{SIAM Journal on Computing} 26(5), pp.~1524--1540, 1997. DOI: 10.1137/S0097539795293639. Publisher abstract retrieved.
\url{https://epubs.siam.org/doi/10.1137/S0097539795293639}.

\bibitem{Babai2016}
L\'aszl\'o Babai. \emph{Graph Isomorphism in Quasipolynomial Time}. Primary preprint, arXiv:1512.03547, with subsequent author revisions. Cited only to correct an outdated baseline in the source article; the algorithm is not proved here.
\url{https://arxiv.org/abs/1512.03547}.

\bibitem{ClayPNP}
Clay Mathematics Institute. \emph{P vs NP}. Official problem statement and status, retrieved 22 November 2026.
\url{https://www.claymath.org/millennium/p-vs-np/}.
\end{thebibliography}

\end{document}
